\section{Introduction}

Axiom~A flows are a class of dynamical systems introduced by Smale \cite{smale} to describe chaotic dynamical systems. It arises in numerous physical problems and it contains two very interesting examples: Morse gradient flows and Anosov flows. On the one hand, the first one is well-known for its link with topology: notably through the Morse inequalities \cite{morse}, stated by Morse in $1920$, which relate the Betti numbers of the manifold and the number of critical points of a Morse function. Given a Riemannian metric on the manifold, Smale \cite{smale-morse} later gave another proof using dynamical arguments and ideas going back to Thom \cite{thom}. On the other hand, Anosov flows were defined first by Anosov in \cite{anosov} to describe the properties satisfied by geodesic flows on negatively curved manifolds and their links with the topology of the manifold are more subtle.

From a purely dynamical point of view, Axiom~A flows are interesting because, once they satisfy a dynamical assumption called the \textit{strong transversality assumption}, they form an open subset of the set of vector fields and the flow induced by any small perturbation of the vector field $X$ is topologically conjugated to the flow induced by~$X$. Moreover, we~have for any vector field on a compact manifold without boundary the equivalence:
\[
\text{Axiom~A + strong transversality assumption} \iff \text{$\Ccal^1$-structurally stable},
\]
which motivated for a long time the study of Axiom~A flows. This equivalence is often referred to as the $\Ccal^1$-structural stability conjecture and it was solved by Robinson \cite{robinson} for the first implication and by Hu \cite{hu} (in dimension $3$) and Hayashi \cite{hayashi}, Wen \cite{wen} (in all dimensions) for the converse statement. The proof of the $\Ccal^1$-structural stability conjecture for diffeomorphisms was previously obtained by Robbin \cite{robbin} for the first implication and by Ma\~{n}e \cite{mane} for the converse statement. A good review on structural stability conjectures can be found in the book of Wen \cite{wen-book}. Up to now, a proof of the $\Ccal^r$-stability conjecture for $r > 1$ is out of reach due to regularity needed in the closing lemma and in the Franks lemma. There exist other transversality assumptions, such as the Kupka-Smale transversality, which was originally introduced (in the case of Axiom~A flows) by Smale under the terminology ``Axiom B' ''. The Kupka-Smale transversality is weaker than the strong transversality since its definition only deals with the periodic orbits of an Axiom~A flow. It also has the nice property to be generic, as it was proved by Palis-de Melo \cite{palis-demelo} on surfaces and by Peixoto \cite{peixoto} in the general case. Using the contributions of many authors for the proof of the stability conjectures, Aoki \cite{aoki} (for~Axiom~A diffeomorphisms) and Gan \cite{gan} (for Axiom~A flows) proved that the set of Axiom~A flows satisfying the strong transversality assumption is exactly the interior of the set of Kupka-Smale flows\footnote{These are the flows with a countable set of periodic orbits and with a finite set of fixed points which are all hyperbolic and which satisfy the Kupka-Smale transversality assumption.}. If we denote by $\KS(M)$ the set of Kupka-Smale flows on $M$, by $\ST(M)$ the set of Axiom~A flows satisfying the strong transversality assumption and by $\rSS(M)$ the set of $\Ccal^1$-structurally stable flows, then the previous discussion can be summarized by the equality
\[
\rSS(M) = \ST(M) = \Int \KS(M).
\]
Equivalently, an Axiom~A flow satisfying the Kupka-Smale transversality assumption is structurally stable if and only if it lies in the interior of $\KS(M)$.

In another direction, the concept of currents\footnote{In coordinates, differential forms with value in the set of distributions. We refer to the book of Schwartz \cite{schwartz} and the lecture notes of Laudenbach \cite{laudenbach-book} for a comprehensive introduction.} turns out to be very useful in the study of gradient flows. More precisely, Laudenbach \cite{laudenbach} and Harvey-Lawson \cite{harvey-lawson} gave a new interpretation of Morse homology in terms of currents by proving the following statement. Let us consider a smooth compact Riemannian manifold $(M,g)$ of dimension $n$ and a smooth Morse function $f$. If $x$ denotes a critical point of index $k\in\lcr1,n\rcr$, then the stable manifold $W^\rs (x)$ for the flow induced by $\nabla_g f$ is an embedded submanifold of dimension $k$ and we have (in the sense of currents)
\begin{equation}\label{eq: intro derivative of de Rham currents}
\partial [W^\rs (x)] = [\partial W^\rs (x)] = \sum_{\substack{\ind y = \ind x - 1 \\ y \text{ critical point}}} n(x,y)[W^\rs (y)]
\end{equation}
for some $n(x,y) \in \Z$ often called the \textit{instanton numbers}. For every $\ell\in\lcr0,n\rcr$, the space $\D'^{,n-\ell} (M)$ of currents of degree $n-\ell$ is defined as the topological dual of the space of differential $\ell$-forms $\Omega^{\ell} (M)$, \ie the space of smooth sections $\Gamma (M; \Lambda^{\ell} T^* M)$. An equivalent formulation for equation \eqref{eq: intro derivative of de Rham currents} is:
\[
\forall \omega \in \Omega^{k-1} (M), \quad \int_{W^\rs(x)} d\omega = \sum_{\substack{\ind y = \ind x - 1 \\ y \text{ critical point}}} n(x,y) \int_{W^\rs (y)} \omega.
\]
This relation is often presented in the following algebraic form. Consider the differential on the complex of critical points defined by
\[
\partial x = \sum_{\substack{\ind y = \ind x - 1 \\y \text{ critical point}}} n(x,y)\cdot y \quad \text{for every }x \in C^k (f) = \{ \text{critical points of index }k \},
\]
with the same numbers $n(x,y)$ as before. The cohomological complex $(C^*(f),\partial)$ is referred to as the Morse complex and is in fact quasi-isomorphic to the de Rham complex, in the sense that the cohomology groups are the same. A remarkable feature of \eqref{eq: intro derivative of de Rham currents} is that it gives a representation of this algebraic complex in terms of currents that are invariant by the gradient flow. Morse inequalities have been generalized to more general dynamical systems as one can witness in the book of Franks \cite{franks_homology}. Yet, to~the best of our knowledge, the previous algebraic procedure does not extend to Axiom~A flows and there is no analogue of its analytical version as constructed by Laudenbach. In~this direction, we~can mention the article of Ruelle and Sullivan~\cite{ruelle-sullivan} in which they constructed similar closed invariant currents for Axiom~A diffeomorphisms. Nevertheless, their construction was only local (near a basic set) and was not enough to recover the whole de~Rham cohomology. More recently, Dang and Rivière showed how to use the theory of Ruelle resonances \cite{DR0} to define a natural cohomological complex of currents associated with Morse-Smale and Anosov flows which are two particular examples of Axiom~A flows. Precisely, for a Morse-Smale gradient vector field $V = \nabla_g f$, they interpreted the Morse complex $(C^* (f),\partial)$ as the complex $(C^* (V),d)$ where
\begin{equation}\label{eq: intro def Morse complex}
C^k (V) = \Ran \biggl(\pi_0^{(k)}: \omega \!\in\! \Omega^{k} (M; \C) \mapsto \frac{1}{2 i \pi} \int_{\gamma_0}(\Lie_{V}^{(k)} + z)^{-1} (\omega) dz \in \D'^{,k} (M;\C) \biggr).
\end{equation}
Here, the right term is the set of generalized eigencurrents associated with the Ruelle resonance $0$ of the Lie derivative of $V = \nabla_g f$ which is usually defined by
\[
\Lie_V^{(k)}: u \in \Omega^k (M;\C) \mto \rest{\frac{d}{dt}}{t= 0} \left(\varphi^t \right)^* u \in \Omega^k (M;\C).
\]
Moreover, $\gamma_0$ denotes a sufficiently small closed curve surrounding $0$ and $d$ is the exterior derivatives acting on currents. In~order to make sense of this linear map $\pi_0^{(k)}$ for every $k\in\lcr0,n\rcr$, they proved in \cite[Prop.\,4.2]{DR0} the meromorphic continuation of the resolvent
\[
z \mto \int_M \left(\Lie_{V}^{(k)} + z \right)^{-1} (\phi) \wedge \psi = \int_M \int_0^{+\infty} e^{-tz} (\varphi^{-t})^*(\phi) dt \wedge \psi
\]
for all $(\phi, \psi) \in \Omega^k (M;\C) \times \Omega^{n-k} (M;\C)$ and for all $k\in\lcr0,n\rcr$. Moreover, they verified that their Morse complex coincides with the complex of currents defined by Laudenbach in \cite{laudenbach} by
\begin{equation}{\label{eq: intro Morse complex current Laudenbach}}
C^k (\nabla_g f) = \Vect_{\C} \bigl\{ [W^\ru (x)] \in \D'^{,k} (M), \: \nabla_g f(x) = 0,\: \ind (x) = k \bigr\}.
\end{equation}
This complex is in fact quasi-isomorphic to the de~Rham complex, \ie the cohomology induced by the Morse complex is isomorphic to the cohomology of the de~Rham complex. In~this article, we~associate a natural cohomological complex to every Axiom~A flow satisfying the strong transversality assumption using similar ideas. Namely, we~first prove:

\begin{thm}\label{thm 2 intro}
Let $V$ be an Axiom~A vector field which satisfies the strong transversality assumptions \eqref{Transversality assumption}. There exists $C_0 > 0$ such that, for every $k\in\lcr0,n\rcr$, the resolvent operator
\[
z \mto(\Lie_V^{(k)} + z)^{-1}: \Omega^k (M;\C) \to \D'^{, k}(M;\C)
\]
is holomorphic on $\Reel (z)\geq C_0$ and continues meromorphically to the whole complex plane $\C$.
\end{thm}

We call \textit{resonances} the poles of the resolvent operators $(\Lie_V^{(k)} + z)^{-1}$ viewed as a meromorphic function on $\C$. In~order to prove that $(\Lie_V^{(k)} + z)^{-1}$ admits a meromorphic extension, one needs to find good Hilbert spaces $\Hil_k$, called anisotropic Sobolev spaces, on which the Lie derivative operators $-\Lie_V^{(k)}$ have discrete spectrum of resonances. Moreover, due the anisotropic order of the Sobolev spaces, the generalized eigencurrents in such spaces are very regular along the unstable manifolds and are very irregular along the stable manifolds. Although the anisotropic Sobolev spaces strongly depend on the dynamics of the flow, the spectrum of $-\Lie_V^{(k)}$ is somehow intrinsic and does not depend on such spaces (as claimed in Theorem \ref{thm 2 intro}). This spectral approach has a long history that we shall briefly recall after the statement of the main theorems. Furthermore, using \textit{analytic Fredholm theory} \cite[App.\,C]{DZ-book}, we~will extend the resolvent operator on the half planes $\Reel (z)\geq -C_k$ (for arbitrarily large~$C_k$) as a meromorphic family of Fredholm operators. In~particular, we~will find that the residue of $(\Lie_V^{(k)} + z)^{-1}$ at $z = 0$ is a finite rank projector whose explicit expression is given~by
\[
\pi_0^{(k)}:= \Res_{z = 0} (- \Lie_V^{(k)}) = \frac{1}{2 i \pi} \int_{\gamma_0} (\Lie_V^{(k)} + z)^{-1} dz: \Omega^k (M;\C) \to \D'^{,k}(M;\C).
\]
Here, $\gamma_0$ denotes a positively oriented closed curve which surrounds the resonance $0$ and no other resonances. The quasi-isomorphism with the de~Rham complex is then given by the maps $\pi_0^{(k)}$. Precisely, we~prove:

\begin{thm}\label{thm intro}
Let $V$ be an Axiom~A vector field which satisfies the strong transversality assumption \eqref{Transversality assumption}. The complex $(C^* (V),d)$ defined in \eqref{eq: intro def Morse complex} is finite dimensional, is~quasi-isomorphic to the de~Rham complex $(\Omega^* (M;\C),d)$ and the quasi-isomorphism is given by the complex morphism $\pi_0^{(*)}$.

Moreover, for every $k\in\lcr0,n\rcr$, there exist a Hilbert space $\Omega^k (M;\C) \subset \Hil_k \subset \D^{\prime,k} (M;\C)$ (with continuous injection) and a positive integer $m_k (0)$ such that
\[
C^k (V) = \Ker \Bigl(\rest{(-\Lie_V^{(k)})^{m_{k} (0)}}{\Hil_{k}} \Bigr).
\]
\end{thm}

In addition, the complex $(C^* (V), d)$ of Theorem \ref{thm intro} generalizes the usual Morse complex, in the sense that \eqref{eq: intro Morse complex current Laudenbach} is satisfied, when $V$ is reduced to a Morse-Smale gradient flow $\nabla_g f$. Theorem \ref{thm intro} is a consequence of Theorem \ref{thm 2 intro} which relies on the construction of the Hilbert space $\Hil_k$ presented earlier and which is the main technical issue of this article. This last theorem shows the existence of a finite dimensional complex representing the de~Rham cohomology and generated by dynamical currents that are almost invariant by the Axiom~A flow, in the sense there may be Jordan blocks (the constant $m_k (0)$ may be greater than $1$).

The Hilbert spaces $\Hil_k$ of the statement are anisotropic Sobolev spaces adapted to the spectral analysis of transfer operators,
\[
L^t (u) = u \circ \varphi^{-t} \qquad \forall u \in \Ccal^{\infty} (M),
\]
as it was initiated by Ruelle \cite{ruelle, ruelle_book}. Such an operator extends to the space of differential forms by setting
\[
L_{(k)}^t (u):= (\varphi^{-t})^* (u), \qquad \forall u \in \Omega^k (M)
\]
and is related with the Lie derivative operator by the formula $\frac{d}{dt} L_{(k)}^t = \Lie_V^{(k)} L_{(k)}^t$ for every $t \in \R$.

\Subsection{Earlier results on meromorphic continuations for Axiom~A flows}

These functional constructions were originally made by Ruelle using Markov partitions in view of studying the mixing properties of these dynamical systems and the analytical properties of the zeta functions associated with their periodic orbits. Similar results have been obtained by Bowen, Fried, Rugh, Dolgopyat and others \cite{bowen2}, \cite{fried_meromorphic}, \cite{rugh}, \cite{dolgopyat} also relying on symbolic dynamics.

Building on some earlier work \cite{BKL} with Blank and Keller for Anosov diffeomorphisms, Liverani introduced in \cite{liverani} Banach spaces of distributions with anisotropic Hölder regularity adapted to contact Anosov flows. Among other things, these spaces enabled him to prove the meromorphic continuation of the resolvent in some half plane slightly beyond $\re(z)=0$. This approach was further developed in subsequent works with Butterley \cite{butterley-liverani} and Giuletti-Pollicott \cite{GLP} to show the meromorphic continuation of $(\Lie_V^{(k)}+z)^{-1}$ to the whole complex plane for every $k\in\lcr0,n\rcr$ and for any smooth Anosov flow. We also refer to \cite{gouezel-liverani_2006, gouezel-liverani} for earlier results with Gouëzel in that spirit for Axiom~A diffeomorphisms, where the resolvent acts on compactly supported functions near a compact locally maximal hyperbolic set. We can also mention \cite{baladi-tsujii_2007} for a different approach (still for diffeomorphisms and still locally near a hyperbolic set called basic) by Baladi and Tsujii using anisotropic Sobolev spaces and methods from Fourier/microlocal analysis.

From another perspective, the general theory of semiclassical resonances \cite{DZ-book}, \cite{helffer-sjostrand} was used to derive alternative approaches to construct Hilbert spaces adapted to the dynamics. First, for smooth Anosov diffeomorphisms, Faure, Roy and Sjöstrand recovered in \cite{FRS} the existence of a discrete spectrum for the transfer operator. Then for general Anosov flows, Faure and Sjöstrand constructed in \cite{FS} Hilbert spaces, referred to as \textit{anisotropic Sobolev spaces}, on which the Lie derivative $\Lie_V^{(0)}$ has discrete spectrum on a large part of the complex plane. Their analysis used the machinery of microlocal analysis as a toolbox and it reduced in some sense the problem to a dynamical question, \ie constructing an escape (or Lyapunov) function adapted to the dynamics on the cotangent space. In~\cite{FS}, the meromorphic extension of the resolvent was obtained for $k = 0$ and this result was extended to every $k\in\lcr0,n\rcr$ in~\cite{dyatlov-zworski} by Dyatlov and Zworski in view of applications to the Ruelle zeta function. More information on the spectrum (\eg band structure) has also been obtained by Faure and Tsujii \cite{tsujii_2010,tsujii_2012, faure-tsujii2013,faure-tsujii2017,faure-tsujii2019,faure-tsujii2021} in the context of Anosov flows and contact Anosov flows using these kinds of methods.

Subsequently, the meromorphic extension of the resolvent has been extended to Axiom~A flows by Dyatlov-Guillarmou \cite{dyatlov-guillarmou_2016, dyatlov-guillarmou_2018} under the assumption that the resolvent acts on differential forms supported near a fixed basic set. Although this analysis was enough to prove the meromorphic continuation of the Ruelle zeta function for Axiom~A flows, it does not seem to be sufficient to deduce Theorem \ref{thm 2 intro} (and thus Theorem \ref{thm intro}) which does not require any support restriction. In~that direction, Dang and Rivière \cite{DR0,DRI,DRII,DRIII} proved meromorphic continuation for globally supported test forms in the case of Morse-Smale gradient flows, a generic subset of Morse gradient flows which satisfy the strong transversality assumption. In~this series of articles, Dang and Rivière gave a complete description of Pollicott-Ruelle resonances, giving a band structure for the spectrum, computed the dimensions of eigenspaces by making explicit the eigenvectors in terms of de~Rham's currents and gave a new proof of Morse-Smale inequalities. In~particular, the link with the topology was made possible through their \textit{global} construction of Sobolev spaces adapted to the dynamics of Morse-Smale flows. In~this article, we~gather the approaches of Dyatlov-Guillarmou and Dang-Rivière so that we obtain the meromorphic continuation of the resolvent acting on globally supported forms for general Axiom~A flows. Finally, we~emphasize that, besides its applications to topology, Theorem \ref{thm 2 intro} also answers a question raised by Baladi in \cite[p.\,152]{baladi} for hyperbolic diffeomorphisms.

\subsection{Back to topology}

Recently, the developments of these analytical tools led to much progress on the link between the topology of the manifold and the spectrum of the Lie derivative, at least for the examples where the functional setup was globally defined, namely Morse-Smale and Anosov flows. We recall here some of these advances.
\begin{itemize}
\item \emph{Contact Anosov flows in dimension $3$.} In that geometric framework, Dyatlov and Zworski \cite{dyatlov-zworski_2017} computed the dimension of $\Ker (\Lie_V^{(k)})^{m_k(0)}$ for every $k\in\lcr0,3\rcr$ and expressed it in terms of the Betti numbers of the manifold. They used this to generalize earlier results of Fried \cite{fried_lefschetz} on the order of vanishing of the Ruelle zeta function. In~particular, their computation holds true for any geodesic flow acting on the unitary cotangent bundle $S^*\Sigma =: M$ of a compact negatively curved surface $\Sigma$. Borns-Weil and Shen \cite{borns-weil-shen} extended \cite{dyatlov-zworski_2017} to the non-orientable case and Hadfield \cite{hadfield} showed a similar result for compact negatively curved surfaces with boundaries.
\item \emph{Anosov flows in high dimension.} Küster-Weich \cite{kuster-weich} computed the dimension of $\Ker (\Lie_V^{(1)})^{m_1(0)}$ in terms of the first Betti number for hyperbolic manifolds of dimension $\neq 3$. Their result also holds for perturbations of hyperbolic metrics.
\item \emph{Perturbation of Anosov flows.} Ceki\'{c}-Paternain \cite{cekic-paternain} gave the first examples of Anosov flows in dimension $3$ which preserve a volume form where the vanishing order of the Ruelle zeta function jumps under perturbation of the flow. Again this was achieved by computing explicitly the dimension of the spaces appearing in the cohomological complex of Theorem \ref{thm intro}. In~dimension $5$, Ceki\'{c}-Dyatlov-Küster-Paternain~\cite{CDKP} found a similar result for geodesic flows on nearly hyperbolic $3$-manifold (the unitary cotangent bundle is $5$-dimensional).
\item \emph{Fried's conjecture \cite{fried_lefschetz,fried_meromorphic}.} Dang-Guillarmou-Rivière-Shen \cite{DGRS} established, in the case of Anosov flows, a criterion in terms of the spaces appearing in Theorem \ref{thm intro} to ensure that the value at $0$ of the twisted Ruelle zeta function is locally constant. It allowed them to prove Fried's conjecture on the Reidemeister torsion for nearly hyperbolic $3$-manifolds. This was further pursued by Chaubet and Dang \cite{chaubet-dang} who used the cohomological complex of Theorem \ref{thm intro} to define a dynamical torsion for contact Anosov flows in any dimension.
\item \emph{Morse-Smale flows.} Dang-Rivière \cite{DRIII} proved Theorem \ref{thm intro} in the case of Morse-Smale and Anosov flows. In~the specific case of Morse-Smale gradient flows \cite{DR_witten}, they also considered the Lie derivative operator as a limit of the Witten Laplacian and they obtained the Ruelle spectrum as a limit of the Witten spectrum. It enabled them to recover the Witten-Helffer-Sjöstrand instanton formula and to prove the Fukaya conjecture on Witten deformation of the wedge product.
\end{itemize}

\subsection{Outline of the proof} We use the microlocal approach to Pollicott-Ruelle resonances of the Lie derivative operator $\Lie_V$ as it was developed by Faure and Sjöstrand. Recall that the proof of Theorem \ref{thm 2 intro} relies on the construction of Hilbert spaces adapted to the dynamics. Following \cite{FS}, defining such spaces can be reduced through some microlocal procedure to the construction of an escape function. More precisely, one has to exhibit a family of functions that are decreasing along the Hamiltonian flow of $H(x,\xi)= \xi(V(x))$ on the cotangent bundle $T^*M$ of $M$. The existence of such decreasing functions, called energy or Lyapunov functions, is already known for the flow on $M$ as soon as $V$ is an Axiom~A flow satisfying the strong transversality assumption. We can cite for example the articles of Conley \cite{conley}, Wilson \cite{wilson} for flows and Pugh-Shub \cite{pugh-shub} for Axiom~A diffeomorphisms satisfying Smale's transversality assumptions. One of the main novelty of this article is to do the same for the induced Hamiltonian flow on $T^*M$. It was already done by Faure-Sjöstrand \cite{FS} for Anosov flows, by Dyatlov-Guillarmou \cite{dyatlov-guillarmou_2016} near a basic set and by Dang-Rivière in \cite{DR0,DRI} for Morse-Smale flows. To construct a decreasing function along this Hamiltonian flow, Dang and Rivière highlighted in the case of gradient flows \cite{DR0} that one needs to prove the compactness of the conormal distribution\enlargethispage{.5\baselineskip}
\[
\bigcup_{x \in M} \left\{ \xi \in S_x^* M \:: \: \xi (T_x W^\ru (x_{-})) = 0, \text{ for } x_- \text{ the critical point s.t } x \in W^\ru (x_{-}) \right\},
\]
where $W^\ru (x_-)$ denotes the unstable manifold of the critical point $x_-$. Nevertheless, to~do so, they made a restriction on the class of Morse-Smale flows, namely, the existence of $\Ccal^1$\textit{-linearization charts} near critical points. Such a restriction is not available for more general Axiom~A flows and we need to proceed differently. In~particular, we~note that our proof allows to remove this linearization assumption in the specific case of Morse-Smale flows. To prove a similar result for Axiom~A flows, we~proceed in three steps.
\begin{itemize}
\item We recall the definition of the strong transversality assumption for Axiom~A flows which generalizes the one used for Morse-Smale gradient flows.
\item Then, we~generalize the compactness result for conormal distributions without using $\Ccal^1$-linearizing charts. This step will require an analysis similar to the local analysis near basic sets performed by Dyatlov-Guillarmou \cite{dyatlov-guillarmou_2016}.
\item We deduce the existence of a \emph{global} escape functions for Axiom~A flows which satisfies the strong transversality assumption by adapting the construction of Faure-Sjöstrand \cite{FS}.
\end{itemize}

Concerning the proof of Theorem \ref{thm intro}, we~recall that that there is a strong analogy with the Hodge-de~Rham Laplace operator\footnote{The derivative $d^*$ denotes the formal adjoint of $d$ in $L^2 (M; \Lambda^k T^* M)$.} $\Delta = d \circ d^* + d^* \circ d = (d + d^*)^2$ acting on differential forms $\Omega^* (M)$ if we remark that
\[
\Lie_V = d\circ \iota_V + \iota_V \circ d = (d + \iota_V)^2.
\]
Note also that both operators $\Delta$ and $\Lie_V$ commute with the exterior derivative $d$. These analogies are at the heart of the proof of Theorem \ref{thm intro}.

\Subsection{Organization of the article}

\begin{itemize}
\item In Section \ref{section Dynamical preliminaries}, we~recall the definition of an Axiom~A flow and introduce the dynamical tools we will need. Furthermore, we~present in this part a few key notions for our analysis which turn out to be related: Smale's order relation on basic sets, strong transversality assumption, filtrations (with open sets) and unrevisited neighborhoods. We also explain how to bypass the $\Ccal^1$-linearizing charts used in Dang-Rivière's articles.
\item In Section \ref{section Escape function}, we~present a possible construction of an escape function and we state a generalization of the compactness result for conormal distributions which takes into account the neutral direction given by the flow direction. The results stated in this part were in fact the most challenging ones to prove.
\item In Section \ref{section Anisotropic Sobolev spaces}, we~define anisotropic Sobolev spaces, in which the Lie derivative operators $\Lie_V^{(k)}$ have discrete spectrum (see Theorem \ref{thm: discrete sprectrum} from which Theorem \ref{thm 2 intro} derives).
\item In Section \ref{section topology}, we~recall how the methods from \cite{DR0, DRIII} can be adapted to deduce Theorem \ref{thm intro} from Theorem \ref{thm 2 intro}.
\item In Section \ref{section: Dynamical proofs}, \ref{section: Compactness result and energy functions for the Hamiltonian flow} and \ref{section: Proof of escape function} we give the proof of the dynamical results such as the construction of energy functions for Axiom~A flows, the proof of the compactness of conormal distributions and the construction of the global escape functions.
\end{itemize}

\subsubsection*{Acknowledgements} The author would like to warmly thank Gabriel Rivière for many explanations about his work with Nguyen Viet Dang and for his careful reading and remarks which contributed a lot to improve this paper. We also thank the anonymous referee for the many suggestions and comments that also helped to improve the exposition of this article.

\section{Dynamical preliminaries.}\label{section Dynamical preliminaries}

Throughout this paper, we~denote by $(M^n,g)$ a smooth compact Riemannian manifold without boundary of dimension $n\geq 1$ together with some smooth Riemannian metric $g$. We also denote by $d_g$ the geodesic distance associated to the metric $g$ and by $|.|_g = \sqrt{g(. \:,.)}$ the norm induced on the fibers of the tangent bundle $TM$ or on the cotangent bundle $T^*M$. To a smooth vector field $V \in \Gamma (TM)$, we~can associate a flow $(\varphi^t)_{t \in \R}$ which solves the Cauchy problem:
\begin{equation}\label{eq: EDO flow}
\forall x \in M, \: \forall t \in \R, \quad
\left\{
\begin{array}{ll}
\dpl\dt \varphi^t (x) = V (\varphi^t (x)), \\[6pt]
\varphi^0 (x) = x.
\end{array}
\right.
\end{equation}
The system \eqref{eq: EDO flow} is highly non-linear in general, which makes difficult to predict the large-time behavior of trajectories, especially in the case of hyperbolic dynamics.

\begin{definition}[{\cite[p.\,796]{smale}}]
A point $x \in M$ is said to be \emph{non-wandering} if for every neighborhood $\Ucal$ of $x$ and every $T > 0$ there exists $t \in \R$ such that $|t| \geq T$ and $\varphi^t (\Ucal) \cap \Ucal \neq \emptyset$. The non-wandering points form a closed invariant subset of $M$, called the \emph{non-wandering set}, that we will denote by $\Omega:= \Omega (\varphi^t)$.
\end{definition}

We refer to Appendix \ref{Appendix: hyperbolic set} or to the books \cite{palis-demelo}, \cite{katok-hasselblatt}, \cite{FH} for a definition of hyperbolic set.

\begin{definition}[{Axiom~A flow, \cite[p.\,803]{smale}, \cite[Def.\,5.3.29]{FH}}]
A flow $\varphi^t: M \to M$ is said to be \emph{Axiom~A} if its non-wandering set $\Omega$ is hyperbolic and can be written as the union of the fixed points and of the closure of its periodic orbits with positive period, \ie\vspace*{-3pt}
\[
\Omega = \Fcal \sqcup \overline{\mathrm{Per} (\varphi^t)}.
\]
\end{definition}

From the definition, one can remark that an Axiom~A flow on a compact manifold must have a finite number of fixed points (since they all are hyperbolic) which are isolated in $\Omega$. It is known from the works of Smale and Bowen that an Axiom~A flow has a non-wandering set which splits into a finite number of hyperbolic invariant compact sets called basic sets:

\begin{prop}[{Spectral decomposition, \cite[\S II.5]{smale}, \cite{bowen2}, \cite[Th.\,5.3.37]{FH}}]
If $\varphi^t$ is an Axiom~A flow, then its non-wandering set $\Omega$ decomposes into a finite union of \emph{basic} sets $K_i$:\vspace*{-3pt}
\begin{equation}\label{eq: decomposition nonwandering set - spectral decomposition}
\Omega = K_1 \sqcup K_2 \sqcup \cdots \sqcup K_N,
\end{equation}
where $K$ basic means:
\begin{itemize}
\item $K$ is compact and hyperbolic;
\item $K$ is locally maximal: there exists some open set $O\subseteq M$ such that\vspace*{-3pt}
\[
K = \bigcap_{t \in \R} \varphi^t (O);
\]
\item $K$ is topologically transitive, \ie there exists a point $x \in K$ such that \\$\overline{(\varphi^t(x))}_{t \in \R_+} = K$.
\end{itemize}
\end{prop}

From now on, $\varphi^t$ will denote an Axiom~A flow on $(M,g)$. We call \textit{attractor} for $\varphi^t$ a basic set $K$ which satisfies\vspace*{-3pt}
\[
K = \bigcap_{t \in \R_+} \varphi^t (O),
\]
for some open set $O \supset K$. Similarly, we~call \textit{repeller} for $\varphi^t$ a basic set $K$ which satisfies\vspace*{-3pt}
\[
K = \bigcap_{t \in \R_-} \varphi^t (O),
\]
for some open set $O \supset K$.

\skpt
\begin{remark}\label{rk: fixed point isolated}
\begin{itemize}
\item The basic sets $K_i$ of the decomposition \eqref{eq: decomposition nonwandering set - spectral decomposition} are the maximal (for the inclusion) locally maximal, $\varphi^t$-invariant, compact, hyperbolic sets which are topologically transitive. A~basic set is equal to the closure of its periodic orbits which can possibly be reduced to a fixed point, \ie have period $0$.
\item If $N = 1$, then $M = \Omega (\varphi^t) = K_1$. This result is a consequence of the local maximality of basic sets. It can be proved using Proposition \ref{Appendix: prop intersection stable unstable sets of basic sets} (which uses the definition of a stable/unstable set of a basic set introduced in the next paragraph) together with the definition of the non-wandering set.
\end{itemize}
\end{remark}

\subsection{Stable and unstable manifolds} We begin by recalling some well-known facts concerning uniformly hyperbolic dynamics which can be found in \cite{katok-hasselblatt} or \cite{dyatlov}. Fix a basic set $K$. For all $\varepsilon > 0$ and all $z \in K$, the \textit{stable manifold}, \textit{weak stable manifold}, \textit{local stable manifold} and \textit{local weak stable manifold} at the point $z$ are defined by
\begin{align*}
W^\rs (z) &:= \{ x \in M: \: d_g(\varphi^t(x), \varphi^t(z)) \underset{t \to+\infty}{\to} 0 \}, \quad W^{\rso} (z):= \bigcup_{t \in \R} \varphi^t \left(W^\rs (z) \right), \\
W_{\varepsilon}^\rs (z) &:= \{ x \in W^\rs (z): \: d_g(\varphi^t(x), \varphi^t(z)) < \varepsilon, \: \: \forall t \in \R_+ \},\\
W_{\varepsilon}^{\rso} (z) &:= \{ x \in M: \: d_g(\varphi^t(x), \varphi^t(z)) < \varepsilon, \: \: \forall t \in \R_+ \}.
\end{align*}

By replacing $\rs$ by $\ru$ and $\varphi^t$ by $\varphi^{-t}$ in the previous equalities, we~could have defined similarly the $\emptyset$/weak/local/local weak unstable manifolds. From this remark, let us only deal with stable manifolds by keeping in mind that everything can be adapted for unstable manifolds. Thanks to the Hadamard-Perron theorem, also called stable manifold theorem, there exists $\varepsilon_0 \ll 1$ such that, for all $z \in K$, the sets $W_{\varepsilon_0}^{\rs/\rso} (z)$ are smooth submanifolds of $M$ of dimension $d_{\rs/\rso}$, which is constant on each basic set. Precise statements and proof of this result can be found in \cite[Th.\,6.4.9, p.\,267]{katok-hasselblatt} for the case of diffeomorphisms and in \cite[Th.\,5, p.\,34]{dyatlov} for the case of flows. In~general, stable manifolds are not embedded submanifolds but only immersed submanifolds, except in the case of Morse flows. Moreover, the stable manifold is related to the local stable manifold thanks to the following formula \cite[p.\,24]{dyatlov}:
\[
W^\rs (z) = \bigcup_{n \geq 0} \varphi^{-n} \left(W_{\varepsilon_0}^\rs (\varphi^{n} (z)) \right), \quad (n \in \N),
\]
which does not depend on $\varepsilon_0$ given by the stable manifold theorem. If $K$ denotes a basic set, then we define its \textit{stable set} by
\[
W^{\rs} (K):= \{ x \in M: \: d(\varphi^{ t} (x), K) \underset{t \to +\infty}{\to} 0 \}.
\]
Thanks to the shadowing lemma \cite[Th.\,18.1.6 p.\,569]{katok-hasselblatt} and to the local maximality of basic sets, this last set decomposes into the stable manifolds of elements of $K$, namely
\[
W^\rs (K) = \bigcup_{z \in K} W^\rs (z).
\]
A proof can be found in \cite[Th.\,5.3.25]{FH} for Axiom~A flows and in \cite[Prop.\,3.10, p.\,53]{bowen3}, \cite[Th.\,6.26, p.\,131]{wen} for Axiom~A diffeomorphisms. For every $\varepsilon >0$, we~can define the local stable set of $K$ by setting
\begin{equation}\label{eq: local decomposition of the stable manifold of a basic set}
W_{\varepsilon}^\rs (K) = \bigcup_{x \in K} W_{\varepsilon}^\rs (x).
\end{equation}

Now, let us present a lemma which was originally given by Smale.

\begin{lemma}[{Partition by stable manifolds, \cite[Cor.\,II.5.3]{smale} }]\label{lemma: spectral decomposition of the manifold} We have the following decomposition of $M$ in stable sets:
\begin{equation} \label{partition of the manifold in stable manifolds}
M = \bigsqcup_{i = 1}^N W^\rs (K_i).
\end{equation}
\end{lemma}

Lemma \ref{lemma: spectral decomposition of the manifold} dictates the behavior of the trajectories outside the non-wandering set. Precisely, if we take an element $x \in M$, then there exists a unique pair $(i,j) \in \lcr1,N\rcr^2$ such that $x \in W^\ru (K_i) \cap W^\rs (K_j)$ and the decomposition (\ref{partition of the manifold in stable manifolds}) provides elements $x_- \in K_i$ and $x_+ \in K_j$ such that $x \in W^\ru (x_-) \cap W^\rs (x_+)$. The point $x_+$ is unique modulo the equivalence relation on $K_j$ given by:
\[
z_1 \sim_j z_2 \iff z_1, z_2 \in K_j \text{ and } z_1 \in W^\rs (z_2).
\]
A similar remark holds for $x_-$.

Let us mention one more lemma concerning the closure of stable sets. The following lemma is illustrated in Figure \ref{fig: closure of local stable manifold} and emphasizes the fact that the closure of the local stable manifold of a basic set is still contained in the stable set of this basic set.

\begin{lemma}\label{lemma: closure of local stable manifold on a basic set}
For every $\varepsilon < \varepsilon_0 $, where $\varepsilon_0$ is the positive constant given by the stable manifold theorem for the basic set $K$, we~have $\overline{W_{\varepsilon}^\rs (K)} \subseteq W_{\varepsilon_0}^\rs (K)$.
\end{lemma}

\begin{proof}
Thanks to the stable manifold theorem, there exist $C, \lambda > 0$ such that for every $z \in K$ and every $\varepsilon < \varepsilon_0$,
\begin{equation}\label{eq: proof closure of local stable set, exp bound}
x \in W_{\varepsilon}^\rs (z) \iff \: \forall t \geq 0, \: d_g \left(\varphi^t (x),\varphi^t (z) \right) < \min \bigl(C e^{-\lambda t} d_g (x,z), \varepsilon \bigr).
\end{equation}
Now, if we fix $\varepsilon > 0$ and if we consider a sequence $((x_n,z_n))_n \in M \times K$ which converges to some limit $(x_{\infty},z_{\infty})$ such that $x_n \in W_{\varepsilon}^\rs (z_n)$ for all $n \in \N$, then the relation~(\ref{eq: proof closure of local stable set, exp bound}) applied to $(x_n, z_n)$ passes to the limit when $n \to +\infty$ and gives in particular $x_{\infty} \in\nobreak W_{\varepsilon_0}^\rs (z_{\infty})$.
\end{proof}

\begin{figure}[!t]
\centering
\def\svgscale{0.6}
\includesvg{closure_local_stable_man}
\caption{Illustration of Lemma \ref{lemma: closure of local stable manifold on a basic set}. Note that the closure of the global stable set of $K_j$ contains the global stable set of $K_i$.}
\label{fig: closure of local stable manifold}
\end{figure}

\subsection{Lifting the dynamics on the cotangent} Since we will use the analysis through escape functions developed in \cite[p.\,329]{FS}, we~start by introducing the lifted flow on the cotangent bundle $T^* M$
\begin{align*}
\Phi^t(x,\xi) = (\varphi^t(x), \left(D\varphi^{t} (x)^{-1} \right)^{\top}(\xi)), \quad \forall (x,\xi) \in T^* M
\end{align*}
and its restriction to the unitary cotangent bundle $S^* M = \{(x,\xi) \in T^* M, \: |\xi|_g = 1 \}$
\begin{equation}\label{eq: def kappa}
\widetilde{\Phi}^t = \kappa \circ \Phi^t \text{ on }S^* M,\text{ where } \kappa: T^* M \setminus 0_M \to S^* M, \: (x,\xi) \mto(x, \sfrac{\xi}{|\xi|_g}).
\end{equation}
Note that the flow $\Phi^t$ extends the flow $\varphi^t$ to the cotangent bundle in the sense that $\pi \circ \Phi^t = \varphi^t \circ \pi $, where $\pi$ denotes the projection from $T^*M$ to $M$ identified with the zero section, and that it is linear on each fiber. Moreover, the flow $\Phi^t$ sends $T^* M \setminus 0_M$ into $T^* M \setminus 0_M$ because the linear map $D\varphi^{t} (x)$ is invertible. Therefore, the flow $\widetilde{\Phi}^t$ is well-defined on $S^*M$ and is generated by a smooth vector field which will be denoted by $\widetilde{X}_H$. To summarize, we~have the following commutative diagrams:
%\begin{center}
%\begin{minipage}{0.45\textwidth}
\[
\begin{tikzpicture}
\matrix (m) [matrix of math nodes,row sep=3em,column sep=4em,minimum width=2em]
{
T^* M & T^* M \\
M & M \\};
\path[-stealth]
(m-1-1) edge node [left] {$ \pi $} (m-2-1)
(m-1-1) edge node [below] {$ \Phi^t $} (m-1-2)
(m-2-1) edge node [below] {$\varphi^t$} (m-2-2)
(m-1-2) edge node [right] {$ \pi $} (m-2-2);
\end{tikzpicture}
\hspace*{2cm}
%\end{minipage}
%\hspace{0.05\textwidth}
%\begin{minipage}{0.45\textwidth}
\begin{tikzpicture}
\matrix (m) [matrix of math nodes,row sep=3em,column sep=4em,minimum width=2em]
{
T^* M \setminus 0_M & T^* M \setminus 0_M \\
S^* M & S^* M \\};
\path[-stealth]
(m-1-1) edge node [left] {$ \kappa $} (m-2-1)
(m-1-1) edge node [below] {$ \Phi^t $} (m-1-2)
(m-2-1) edge node [below] {$\widetilde{\Phi}^t$} (m-2-2)
(m-1-2) edge node [right] {$ \kappa $} (m-2-2);
\end{tikzpicture}
\]
%\end{minipage}
%\end{center}

Since our analysis will take place in $T^* M$, we~also define the dual distributions associated with the neutral $E_\ro$, stable $E_\rs$ and unstable $E_\ru$ distributions\footnote{On a basic set $K$, we~also use the notations $E_{\rso}^*$ for $E_\rs^* + E_\ro^*$ and $E_{\ruo}^*$ for $E_\ru^* + E_\ro^*$.} which appear in the definition of hyperbolicity (see Appendix \ref{Appendix: hyperbolic set}) at any point $z \in \Omega$ by
\begin{align*}
(E_\ro^*(z))(E_\ru(z) &\oplus E_\rs(z)) = 0, \\
(E_\ru^*(z))(E_\ru(z) \oplus E_\ro(z)) = 0, \quad & \quad (E_\rs^*(z))(E_\rs(z) \oplus E_\ro(z)) = 0.
\end{align*}
This gives us a hyperbolic splitting of the cotangent bundle:
\begin{align*}
T_z^* M = E_\rs^* (z) \oplus E_\ru^* (z) \oplus E_\ro^*(z).
\end{align*}

\Subsection{Extension of the invariant distributions outside the non-wandering set}

Thanks to the partition's lemma \ref{partition of the manifold in stable manifolds} by stable manifolds, we~can extend the previous definitions outside the non-wandering set.

\begin{definition}\label{def distributions duales stables et instables sur Omega}
For every $x \in W^\rs (x_+) \cap W^\ru (x_{-})$ with $x_- \in K_i$ and $x_+ \in K_j$, we~define the spaces $E_\rs^*(x)$, $E_\rso^* (x)$, $E_\ru^*(x)$, $E_{\ruo}^* (x)$ to be the largest spaces satisfying:
\begin{align*}
E_\rs^*(x) (T_x W^{\rso} (x_+)) = 0, &\qquad E_\rso^* (x)(T_x W^{\rs} (x_+)) = 0,\\
E_\ru^*(x) (T_x W^{\ruo} (x_-)) = 0, &\qquad E_{\ruo}^* (x)(T_x W^{\ru} (x_-)) = 0 \:.
\end{align*}
Moreover, the definition does not depend on the choice of $x_+$ and $x_-$ in $W^\rs (x_+) \cap K_j$ and $W^\ru (x_-) \cap K_i$ respectively.
\end{definition}

\pagebreak[2]
\skpt
\begin{remark}
\begin{itemize}
\item If $x_+$ is a fixed point, then $W^{\rso} (x_+) = W^{\rs} (x_+)$ and consequently $E_\rs^*(x) = E_{\rso}^*(x) \subseteq \{\xi (V(x)) = 0 \}$.
\item The distributions $E_\rs^*, E_\ru^*, E_{\rso}^*$ and $E_{\ruo}^*$ are defined on the whole manifold and are $\Phi^t$-invariant.
\item Recall that fixed points are isolated (Remark \ref{rk: fixed point isolated}). So if $x_+$ is not a fixed point then we get that $\dim E_{\rso}(x_+) = \dim E_{\rs}(x_+) + 1$ and $\dim W^{\rso} (x_+) = \dim W^{\rs} (x_+) + 1$. Therefore, we~obtain $E_{\rs}^* (x) = E_{\rso}^* (x) \cap \{\xi (V(x)) = 0 \}$ and $\dim E_{\rso}^* (x) = \dim E_{\rs}^* (x) + 1$.
\item We can extend the \textit{neutral} distribution outside the non-wandering set by fixing, for every $x \in W^\ru (x_-) \cap W^\rs (x_+)$,
\[
E_\ro^* (x):= E_{\ruo}^* (x_-) \cap E_{\rso}^* (x_+) = \left\{ \xi \in T^* M, \: \xi \left(T_x W^{\ru} (x_-) + T_x W^{\rs} (x_+)\right) = 0 \right\}.
\]
\end{itemize}
\end{remark}

\subsection{Strong transversality assumption}

\begin{figure}[!t]
\begin{minipage}[t]{0.46\textwidth}
\definecolor{uququq}{rgb}{0.25,0.25,0.25}
\begin{tikzpicture}[line cap=round,line join=round,>=triangle 45,x=1.3cm,y=1.3cm, scale=0.85]
\clip(-2.5,-2) rectangle (2.6,1.9);
\draw (-2,0)-- (2,0);
\draw (0,1.5)-- (0,-1.5);
\draw [->] (0,1.5) -- (-1,1.5);
\draw [->] (0,1.5) -- (1,1.5);
\draw [->] (0,-1.5) -- (1,-1.5);
\draw [->] (0,0) -- (-1,0);
\draw [->] (0,0) -- (1,0);
\draw [->] (0,-1.5) -- (0,-0.75);
\draw [->] (0,1.5) -- (0,0.75);
\draw [->] (0,-1.5) -- (-1,-1.5);
\draw (-2,1.5)-- (2,1.5);
\draw (2,1.5)-- (2,-1.5);
\draw (2,-1.5)-- (-2,-1.5);
\draw (-2,-1.5)-- (-2,1.5);
\draw (-0.1,1.94) node[anchor=north west] {$K_1$};
\draw (-0.1,-1.49) node[anchor=north west] {$K_1$};
\draw (2.12,1.28) node[anchor=north west] {\large $\mathbb{T}^2$};
\draw (-2.2,1.94) node[anchor=north west] {$K_4$};
\draw (1.95,1.93) node[anchor=north west] {$K_4$};
\draw (1.93,-1.49) node[anchor=north west] {$K_4$};
\draw (-2.15,-1.49) node[anchor=north west] {$K_4$};
\draw (-0.07,0.43) node[anchor=north west] {$K_2$};
\draw (-2.6,0.24) node[anchor=north west] {$K_3$};
\draw (1.96,0.24) node[anchor=north west] {$K_3$};
\draw [->] (-2,0) -- (-2,0.79);
\draw [->] (-2,0) -- (-2,-0.82);
\draw [->] (2,0) -- (2,0.79);
\draw [->] (2,0) -- (2,-0.83);
\draw [shift={(-1,1.14)}] plot[domain=3.18:6.21,variable=\t]({1*0.92*cos(\t r)+0*0.92*sin(\t r)},{0*0.92*cos(\t r)+1*0.92*sin(\t r)});
\draw [shift={(-0.99,1.45)}] plot[domain=3.65:5.76,variable=\t]({1*0.81*cos(\t r)+0*0.81*sin(\t r)},{0*0.81*cos(\t r)+1*0.81*sin(\t r)});
\draw [shift={(-0.97,2.76)}] plot[domain=4.32:5.06,variable=\t]({1*1.67*cos(\t r)+0*1.67*sin(\t r)},{0*1.67*cos(\t r)+1*1.67*sin(\t r)});
\draw [->] (-0.95,1.09) -- (-1.06,1.09);
\draw [->] (-0.99,0.64) -- (-1.1,0.64);
\draw [->] (-0.97,0.22) -- (-1.08,0.23);
\begin{scriptsize}
\fill [color=uququq] (-2,1.5) circle (1.5pt);
\fill [color=uququq] (2,1.5) circle (1.5pt);
\fill [color=uququq] (2,-1.5) circle (1.5pt);
\fill [color=uququq] (-2,-1.5) circle (1.5pt);
\fill [color=uququq] (-2,0) circle (1.5pt);
\fill [color=uququq] (2,0) circle (1.5pt);
\fill [color=uququq] (0,1.5) circle (1.5pt);
\fill [color=uququq] (0,0) circle (1.5pt);
\fill [color=uququq] (0,-1.5) circle (1.5pt);
\end{scriptsize}
\end{tikzpicture}
\\
\refstepcounter{toto}\label{fig: Axiom A flow on the torus - without transversality assumption}\renewcommand{\baselinestretch}{.95}\normalfont
{\noindent\eqref{fig: Axiom A flow on the torus - without transversality assumption} \smaller Axiom~A flow which does not satisfy the (strong) transversality assumption: existence of a saddle-connection.}
\end{minipage}
\hfill
\begin{minipage}[t]{0.46\textwidth}
\definecolor{uququq}{rgb}{0.25,0.25,0.25}
\begin{tikzpicture}[line cap=round,line join=round,>=triangle 45,x=1.3cm,y=1.3cm, scale=0.85]
\clip(-2.4,-2) rectangle (2.7,1.8);
\draw (-2,0)-- (2,0);
\draw (0,1.5)-- (0,-1.5);
\draw [->] (0,1.5) -- (-1,1.5);
\draw [->] (0,1.5) -- (1,1.5);
\draw [->] (2,1.5) -- (2,0.75);
\draw [->] (2,-1.5) -- (2,-0.75);
\draw [->] (0,-1.5) -- (1,-1.5);
\draw [->] (-2,1.5) -- (-2,0.75);
\draw [->] (-2,-1.5) -- (-2,-0.75);
\draw [->] (0,0) -- (-1,0);
\draw [->] (0,0) -- (1,0);
\draw [->] (0,-1.5) -- (0,-0.75);
\draw [->] (0,1.5) -- (0,0.75);
\draw [->] (0,-1.5) -- (-1,-1.5);
\draw (-2,1.5)-- (2,1.5);
\draw (2,1.5)-- (2,-1.5);
\draw (2,-1.5)-- (-2,-1.5);
\draw (-2,-1.5)-- (-2,1.5);
\draw [shift={(-0.26,-0.32)}] plot[domain=1.66:2.79,variable=\t]({1*1.67*cos(\t r)+0*1.67*sin(\t r)},{0*1.67*cos(\t r)+1*1.67*sin(\t r)});
\draw [shift={(-1.87,2.03)}] plot[domain=4.83:5.8,variable=\t]({1*1.87*cos(\t r)+0*1.87*sin(\t r)},{0*1.87*cos(\t r)+1*1.87*sin(\t r)});
\draw [shift={(1.98,2.14)}] plot[domain=3.6:4.55,variable=\t]({1*2.03*cos(\t r)+0*2.03*sin(\t r)},{0*2.03*cos(\t r)+1*2.03*sin(\t r)});
\draw [shift={(-0.06,-0.76)}] plot[domain=0.53:1.37,variable=\t]({1*2.15*cos(\t r)+0*2.15*sin(\t r)},{0*2.15*cos(\t r)+1*2.15*sin(\t r)});
\draw [shift={(-1.8,-1.85)}] plot[domain=0.38:1.51,variable=\t]({1*1.72*cos(\t r)+0*1.72*sin(\t r)},{0*1.72*cos(\t r)+1*1.72*sin(\t r)});
\draw [shift={(-0.14,0.5)}] plot[domain=3.59:4.6,variable=\t]({1*1.91*cos(\t r)+0*1.91*sin(\t r)},{0*1.91*cos(\t r)+1*1.91*sin(\t r)});
\draw [shift={(1.78,-1.93)}] plot[domain=1.68:2.72,variable=\t]({1*1.78*cos(\t r)+0*1.78*sin(\t r)},{0*1.78*cos(\t r)+1*1.78*sin(\t r)});
\draw [shift={(0.28,0.24)}] plot[domain=4.77:5.87,variable=\t]({1*1.61*cos(\t r)+0*1.61*sin(\t r)},{0*1.61*cos(\t r)+1*1.61*sin(\t r)});
\draw [->] (-1.26,1.01) -- (-1.38,0.91);
\draw [->] (-0.76,0.53) -- (-0.83,0.48);
\draw [->] (0.66,0.59) -- (0.73,0.53);
\draw [->] (1.14,1.02) -- (1.24,0.95);
\draw [->] (0.69,-0.52) -- (0.8,-0.44);
\draw [->] (1.22,-1.06) -- (1.29,-1.01);
\draw [->] (-1.21,-1.09) -- (-1.29,-1.03);
\draw [->] (-0.72,-0.51) -- (-0.81,-0.44);
\draw (-0.1,1.94) node[anchor=north west] {$K_1$};
\draw (-0.1,-1.50) node[anchor=north west] {$K_1$};
\draw (2.12,1.28) node[anchor=north west] {\large $\mathbb{T}^2$};
\draw (-2.18,1.94) node[anchor=north west] {$K_2$};
\draw (1.95,1.93) node[anchor=north west] {$K_2$};
\draw (1.93,-1.49) node[anchor=north west] {$K_2$};
\draw (-2.15,-1.49) node[anchor=north west] {$K_2$};
\draw (-0.07,0.4) node[anchor=north west] {$K_3$};
\draw (-2.56,0.24) node[anchor=north west] {$K_4$};
\draw (1.96,0.24) node[anchor=north west] {$K_4$};
\begin{scriptsize}
\fill [color=uququq] (-2,1.5) circle (1.5pt);
\fill [color=uququq] (2,1.5) circle (1.5pt);
\fill [color=uququq] (2,-1.5) circle (1.5pt);
\fill [color=uququq] (-2,-1.5) circle (1.5pt);
\fill [color=uququq] (-2,0) circle (1.5pt);
\fill [color=uququq] (2,0) circle (1.5pt);
\fill [color=uququq] (0,1.5) circle (1.5pt);
\fill [color=uququq] (0,0) circle (1.5pt);
\fill [color=uququq] (0,-1.5) circle (1.5pt);
\end{scriptsize}
\end{tikzpicture}
\\
\refstepcounter{toto}\label{fig: Axiom A flow on the torus - with transversality assumption}\renewcommand{\baselinestretch}{.95}\normalfont
{\noindent\eqref{fig: Axiom A flow on the torus - with transversality assumption} \smaller Axiom~A flow which satisfies the\\ (strong) transversality assumption.}
\end{minipage}
\caption{Some Axiom~A flows on the 2-torus.}\label{fig:2}
\end{figure}

In this part, we~briefly recall the definition of the strong transversality assumption presented in the introduction, which turned out to be the good transversality assumption to prove the structural $\Ccal^1$\nobreakdash-sta\-bil\-i\-ty conjecture.

We say that $\varphi^t$ satisfies the strong transversality assumption if for every $(x_-,x_+) \in K_i \times K_j$ and for every $x \in W^\ru (x_-) \cap W^\rs (x_+)$ we have
\begin{equation}\label{Transversality assumption - first}
T_x W^{\ruo} (x_-) + T_x W^{\rso} (x_+) = T_x M.
\end{equation}
Since both spaces $T_x W^{\ruo} (x_-) $ and $T_x W^{\rso} (x_+)$ contain the flow direction $\R\cdot V(x)$, an equivalent formulation to (\ref{Transversality assumption - first}) is:
\begin{equation}\label{Transversality assumption}
T_x W^{\ru} (x_-) + T_x W^{\rso} (x_+) = T_x W^{\ruo} (x_-) + T_x W^{\rs} (x_+) = T_x M.
\end{equation}
Moreover, one can remark that the strong transversality assumption does not depend on the choice of $x_- \in K_i$ and $x_+ \in K_j$ such that $x \in W^\ru (x_-) \cap W^\rs (x_+)$ in (\ref{Transversality assumption - first}) or~\eqref{Transversality assumption}.
\begin{remark}If $x_-$ or $x_+$ is a fixed point then the strong transversality assumption reads
\[
T_x W^{\ru} (x_-) + T_x W^{\rs} (x_+) = T_x M.
\]
This is in particular the case for Morse-Smale gradient flows \cite{DR0} where both $x_-$ and~$x_+$ are fixed points.
\end{remark}

From \eqref{Transversality assumption} and directly from the definitions, we~can deduce the following \textit{disjointness} properties:
\begin{align*}
E_\rs^* \cap E_{\ruo}^* = 0_M \quad\text{and}\quad E_{\rso}^* \cap E_\ru^* = 0_M,
\end{align*}
where $0_M$ denotes the null section of $T^*M$.

\subsection{Order relation}

When Smale \cite{smale} defined Axiom~A flows, he exhibited a relation between basics sets of an Axiom~A flow. Precisely, for two basic sets $K_i$ and~$K_j$, he defined the relation $\leq$ by
\begin{equation}\label{order relation for basic sets}
\begin{split}
K_i \leq K_j &\iff W^\ru (K_i) \cap W^\rs (K_j) \neq \emptyset\\
&\iff \exists x \in M, \: \exists (x_-, x_+) \in K_i \times K_j, \: \: x \in W^\ru (x_-) \cap W^\rs (x_+).
\end{split}
\end{equation}
This relation is illustrated in Figures \ref{fig: order relation between two fixed points} and \ref{fig: order relation between two basic sets} and gives a graph structure on the family of basic sets, where the edges of the graph are the basic sets and an arrow between two basic sets (from $K_i$ to $K_j$) is given by the relation $K_i \leq K_j$. When $\varphi^t$ satisfies the strong transversality assumption, then the next theorem (originally due to Smale) states that the relation $\leq$ is a partial order relation. Since we were not able to find out an explicit proof in the literature and for the sake of completeness, we~present its proof in Appendix \ref{Appendix: proof of Smale ordering}.

\begin{figure}
\begin{minipage}[t]{.49\linewidth}
\def\svgscale{0.3}
\includesvg{deux_ens_basiques}
\caption{Order relation between two fixed points.}
\label{fig: order relation between two fixed points}
\end{minipage}
\begin{minipage}[t]{.49\linewidth}
\def\svgscale{0.4}
\includesvg{two_basic_sets}
\caption{Order relation between two basic sets.}
\label{fig: order relation between two basic sets}
\end{minipage}
\end{figure}

\begin{thm}[{Smale, \cite[Prop.\,8.5, p.\,784]{smale}}]\label{thm Smale}
If $\varphi^t$ is an Axiom~A flow, then we have for every basic set $K$
\[
W^\ru (K) \cap W^\rs (K) = K.
\]
Moreover, if $\varphi^t$ satisfies the strong transversality assumption \eqref{Transversality assumption}, then the relation $\leq$ defines a partial order relation. An equivalent definition of the relation $\leq$ is then given by
\begin{equation}\label{eq: thm Smale - equivalence order relation}
K_i \leq K_j \iff W^\rs (K_i) \subseteq \overline{W^\rs (K_j) } \iff W^\ru (K_j) \subseteq \overline{W^\ru (K_i) },
\end{equation}
and we have
\begin{equation}\label{eq: thm Smale - closure of stable and unstable sets}
\overline{W^\ru (K_i) } = \bigcup_{j, K_j \geq K_i} W^\ru (K_j) \: \text{ and } \: \overline{W^\rs (K_j) } = \bigcup_{i, K_i \leq K_i} W^\rs (K_i).
\end{equation}
\end{thm}

As we can see in Figure \ref{fig:2}\eqref{fig: Axiom A flow on the torus - without transversality assumption}, relation (\ref{order relation for basic sets}) may not be an order relation when the strong transversality assumption does not hold.
Another consequence of this theorem is that strong transversality implies the \textit{no-cycle property}: the induced graph has no cycle. This fact is necessary to ensure the existence of non-constant Lyapunov functions for the flow $\varphi^t$.

\subsubsection{Total order relation}\label{subsubsection total order relation}

In order to use mathematical induction, we~need to consider a total order on the family of basic sets. It can be achieved as soon as the Axiom~A flow $\varphi^t$ satisfies the no-cycle property. Therefore, if we assume the no-cycle property, then we define a total order relation from Smale's order relation $\leq$ on the basic sets as an order relation on $\lcr1,N\rcr$ compatible with the partial order relation~$\leq$ in the sense that
\begin{equation}\label{eq: total order relation, indices compatible}
K_i \leq K_j \implies i \leq j.
\end{equation}
\emph{From now on}, we~assume the no-cycle property and we fix a total order relation. In~Appendix \ref{Appendix: proof of Smale ordering}, we~prove that the set $\bigcup_{j\geq i }W^\ru (K_j)$ is compact for every $i\in\lcr1,N\rcr$. Be aware that $\overline{W^\ru (K_i)}$ and $\bigcup_{j\geq i }W^\ru (K_j)$ are not equal in general, see for example Figure \ref{fig: graph structures on the 2 torus}.

\begin{figure}[t]
\begin{minipage}[c]{.46\linewidth}
\centering
\def\svgscale{0.6}
\includesvg{graphe_structure_tore}
\end{minipage}
\begin{minipage}[c]{.46\linewidth}
\centering
\def\svgscale{0.6}
\includesvg{graphe_structure_tore2}
\end{minipage}
\caption{Both graph correspond to the Axiom~A flow on the \hbox{2-torus} of Figure \ref{fig:2}\eqref{fig: Axiom A flow on the torus - with transversality assumption}. A path in the oriented graph corresponds to a connecting orbit between the basic sets. The left one has its indices compatible with the graph structure contrary to the right one.}
\label{fig: graph structures on the 2 torus}
\end{figure}

\Subsection{Filtrations and unrevisited neighborhoods}

An important concept throughout our analysis is the concept of filtration. Even though the term of filtration usually refers to an increasing sequence of subcomplexes of a simplicial complex, we~give here an open version deeply related to Morse homology where the subcomplexes are open sets, \ie submanifolds of dimension $n = \dim M$ without boundary. To see the analogy, let us consider a Morse function $f: M \to \R$ which has~$N$ critical points $x_1,\dots, x_N$ that all satisfy $f(x_i) = i$ to simplify. Then, a~filtration is given by the family of open sets $(f^{-1} (]-\infty,i +\sfrac{1}{2}[))_{0 \leq i \leq N}$. For a general Axiom~A flow, we~give a definition of a filtration which appeared in the works of Smale \cite{smale}, Robbin \cite[Lem.\,7.9, p.\,471]{robbin}.

\begin{definition}[Filtration]\label{Def: filtration} Let $\varphi^t$ be an Axiom~A flow which satisfies the no-cycle property. Let us consider a total order relation on the basic sets in the sense of \eqref{eq: total order relation, indices compatible}. A sequence of open sets $(\Ocal_i^-)_{0 \leq i \leq N}$ is said to be a \emph{filtration} for $\varphi^{-1}$ if the following conditions hold:
\begin{enumeratei}
\item The sequence is increasing:
\[
\emptyset = \Ocal_0^- \subseteq \Ocal_{1}^- \subseteq \cdots \subseteq \Ocal_N^- = M.
\]
\item For every $i\in\lcr1,N\rcr$, the open sets and their closures are $\varphi^{-1}$-stable: $\varphi^{-1} \left(\Ocal_{i}^- \right) \subseteq \Ocal_{i}^- $, $\varphi^{-1} \left(\overline{\Ocal_{i}^-} \right) \subseteq \overline{\Ocal_{i}^-}$.
\item For every $i\in\lcr1,N\rcr$, we~have $K_i \subseteq \Ocal_i^- \setminus \overline{\Ocal_{i-1}^-}$.
\end{enumeratei}
\end{definition}

\begin{figure}[t]
\centering
\def\svgscale{0.6}
\includesvg{filtration}
\caption{Example of a filtration on the sphere $\Sn^2$.}
\label{fig: filtration on the sphere}
\end{figure}

\skpt
\begin{remark}\label{rk: filtration}
\begin{itemize}
\item Any filtration for $\varphi^{-1}$ induces a filtration for $\varphi^{1}$ by taking the interior of the complementary of each open set, \ie by setting $\Ocal_{i}^{+}:= \Int \left(\left(\Ocal_{N-i}^{-} \right)^c \right)$ for every $i$. Indeed, from $\varphi^{-1} \left(\Ocal_{N-i}^- \right) \subseteq \Ocal_{N-i}^-$, we~deduce
\[
\varphi^1 (\Ocal_{i}^+) = \varphi^1 \left(\Int \left(\left(\Ocal_{N-i}^- \right)^c \right) \right) \subseteq \varphi^1 \left(\left(\Ocal_{N-i}^- \right)^c \right) \subseteq \left(\Ocal_{N-i}^- \right)^c.
\]
So, $\varphi^1 (\Ocal_{i}^+)$ is an open subset of $\left(\Ocal_{N-i}^- \right)^c$, and by definition of the interior we get $\varphi^1 (\Ocal_{i}^+) \subseteq \Int \left(\left(\Ocal_{N-i}^- \right)^c \right) = \Ocal_{i}^+$.
Moreover, we~still have $\emptyset = \Ocal_0^+ \subseteq \Ocal_{1}^+ \subseteq \cdots \subseteq \Ocal_N^+ =\nobreak M$ and $K_{i} \subseteq \Ocal_{N-i+1}^+ \setminus \overline{\Ocal_{N-i}^+}$.
\item $\Ocal_i^- $ is a neighborhood of $\bigsqcup_{j\leq i} W^\rs (K_i)$. Indeed, $\Ocal_i^- $ contains every basic set $K_j$ for $j \leq i$ and if
$x \in W^\rs (K_j)$ then there exists $k \in \N$ such that $\varphi^k (x) \in \Ocal_i^-$. Thus, we~deduce $x \in \varphi^{-k} (\Ocal_i^-) \subset \Ocal_i^-$.
\item For every $i\in\lcr1,N\rcr$, let us define the set
\[
\Vcal_i:= \Ocal_i^- \cap \Ocal_{N-i+1}^+.
\]
One can check that $\Vcal_i$ is a neighborhood of $K_i$ which satisfies $\overline{\Vcal_i}\cap \Omega = K_i$ and the following property: for all $m \in \N$ and for all $x \in \Vcal_i$, if we have $x \in \Vcal_i$ and $\varphi^m (x) \in \Vcal_i$ then we must have $\varphi^k (x) \in \Vcal_i$ for all $k\in\lcr0,m\rcr$. In~the example presented in Figure~\ref{fig: filtration on the sphere}, the basic set $K_3$ belongs to $\Ocal_3^- \cap \Ocal_2^+$.
\end{itemize}
\end{remark}

This last remark brings us to the next definition.

\begin{definition}[{Unrevisited set, \cite[p.\,463]{robbin} and \cite{shub}}]Let $X$ be a smooth manifold. A set $W \subseteq X$ is called \emph{unrevisited} for a diffeomorphism $f: X \to X$ if for any integer $m \in \N$,
\[
x, f^m (x) \in W \implies \forall k \in \{0,\dots,m \},\; f^k (x) \in W.
\]
\end{definition}

We say that a set is unrevisited for the flow $\varphi^t$ if it is unrevisited for the time\nobreakdash-$1$ map~$\varphi^1$. According to the last point of Remark \ref{rk: filtration}, the existence of unrevisited neighborhoods is a consequence of the existence of a filtration for $\varphi^t$ and of a filtration for $\varphi^{-t}$. Moreover, the existence of a filtration can be deduced from the existence of a continuous Lyapunov function for the flow which goes back to the work of Conley \cite[p.\,20]{conley} and which can be found in \cite[Th.\,1.5.44 \& 5.3.47]{FH}:

\begin{thm}[Conley's fundamental theorem of dynamical systems]\label{prop: existence of continuous Lyapunov functions}
Let $\varphi^t$ be an Axiom~A flow which satisfies the no-cycle property.\footnote{We use here the fact that the chain recurrent set is equal to the non-wandering set when the Axiom~A flow satisfies the no-cycle property, which can be found in \cite[Th.\,5.3.47]{FH}. We refer to \cite[Def.\,5.3.42]{FH} for a definition of the no-cycle property.} There exists a continuous function $E: M \to \R$ which is constant on each basic set and take distinct values such that,
\[
E \left(\varphi^t (x) \right) < E (x), \qquad \forall t > 0, \quad \forall x \in M \setminus \Omega,
\]
and $E \circ \varphi^t \leq E$ everywhere for all $t \geq 0$.
\end{thm}

Considering the sublevel sets associated with such a function $E$, it is not hard to deduce

\begin{corollary}[Existence of a filtration]\label{cor: existence of filtrations} Under the same hypothesis, there exists a filtration for the flow $\varphi^t$ and there exists a filtration for the flow $\varphi^{-t}$ in the sense of Definition \ref{Def: filtration}.
\end{corollary}

Furthermore, intersecting two filtrations (one for $\varphi^t$ and one for $\varphi^{-t}$), we~get:

\begin{corollary}[Definition and existence of unrevisited neighborhoods]\label{prop: Robbin existence of unrevisited neighborhoods}Let $K$ be a basic set. We say that an open set $\Vcal$ of $K$ is an \emph{unrevisited neighborhood} of $K$ if $K \subset \Vcal$, $\Vcal$ is an unrevisited set, $\overline{\Vcal}$ is also an unrevisited set and $\Omega \cap \overline{\Vcal} = K$.
According to Theorem \ref{prop: existence of continuous Lyapunov functions} and Corollary \ref{cor: existence of filtrations}, there exist arbitrarily small unrevisited neighborhoods $\Vcal$ of $K$.
\end{corollary}

Filtrations and unrevisited neighborhoods play an important role in the dynamical proof of this article. The existence of Lyapunov functions, of filtrations and of unrevisited neighborhood are related and almost equivalent. To make it easier to read, we~sketch the nonlinear dependencies between these three notions:
\begin{enumerate}
\item We recall that if an Axiom~A flow satisfies the strong transversality assumption then it satisfies the no-cycle property.
\item According to Conley, for an Axiom~A flow satisfying the no-cycle property, there exists a continuous Lyapunov function, and this implies the existence of a filtration.
\item From this filtration, we~can obtain arbitrarily small unrevisited neighborhoods.
\item Once we have one filtration, we~can construct filtrations for any total order relation---see Lemma \ref{lemma: equivalence filtration and unrevisited}. This lemma mainly\footnote{In particular, its proof does not use the compactness of the sets $ \bigcup_{i \leq j} W^\ru (K_j)$ and $ \bigcup_{i \leq j} W^\rs (K_i)$ which is claimed in Theorem \ref{thm Smale}.} relies on Proposition \ref{Appendix: prop intersection stable unstable sets of basic sets} of Appendix~\ref{Appendix: proof of Smale ordering}.
\item Using filtrations for every total order relation constructed in Lemma \ref{lemma: equivalence filtration and unrevisited}, we~can deduce Corollary \ref{cor: corollary technical lemma - closure of unstable sets}. Namely, we~get the compactness of all the sets $ \bigcup_{i \leq j} W^\ru (K_j)$ and $ \bigcup_{i \leq j} W^\rs (K_i)$, for every total order relation defined in Section \ref{subsubsection total order relation}.
\item Once we have the compactness of these sets, we~can construct smooth Lyapunov functions adapted to the total order relation. Precisely, we~can construct Lyapunov functions with arbitrary values on the basic sets accordingly to the total order relation---see Proposition \ref{prop energy function for Axiom A flows}.
\end{enumerate}

\begin{remark}Theorem \ref{prop: existence of continuous Lyapunov functions} implies actually something stronger. Indeed, the filtration given by the sublevel sets of the continuous Lyapunov function is stable by $\varphi^{-t}$ for every $t > 0$ and, for every basic set $K$, there exists an open neighborhood $\Vcal \supset K$ such that: for all $T \geq 0$,
\[
x, \varphi^T (x) \in \Vcal \implies \forall t \in [0,T], \: \: \varphi^t (x) \in \Vcal.
\]
Even if this last definition seems to be more natural for flows, we~chose the diffeo-like definition which will be more convenient for the analysis of the lifted Hamiltonian dynamics on the phase space as we will see later on. Unrevisited neighborhoods will be a very important tool of our analysis. They will be a purely dynamical alternative to the $\Ccal^1$ linearizing charts near critical points which were used in \cite{DR0} for Morse-Smale gradient flows, as we can witness in Figure \ref{fig: unrevisited neighborhoods near a basic set}.
\end{remark}

Let us mention some properties satisfied by the unrevisited neighborhoods. All the results mentioned below about unrevisited neighborhoods were not precisely stated in the literature. They should be attributed to Conley, Robbin, Hirsch, Palis, Pugh, Shub to the best of our knowledge. Precisely,

\begin{figure}[!t]
\centering
\def\svgscale{0.7}
\includesvg{unrevisited_neighborhoods}
\caption{Illustration of some unrevisited neighborhoods near a basic set.}
\label{fig: unrevisited neighborhoods near a basic set}
\end{figure}

\begin{enumerate}\renewcommand{\theenumi}{{P\arabic{enumi}}}
\item\label{P1}
The intersection of two unrevisited sets is also an unrevisited.
\item\label{P2}
We have a uniform approximation of the stable and unstable manifolds by unrevisited neighborhoods in the sense of the following lemma.
\end{enumerate}

\begin{lemma}[Uniform convergence of unrevisited neighborhoods]\label{lemme: cv nonrevisite vers variete stable et instable}
Let $\Vcal$ be an unrevisited neighborhood of a basic set $K$ which satisfies $\overline{\Vcal} \cap \Omega = K$. The following equalities are satisfied:
\begin{equation}\label{eq: lemma unif convergence - closure of unrevisited}
\bigcap_{m \in \N} \overline{\Vcal} \cap \varphi^m (\overline{\Vcal}) = W^\ru (K) \cap \overline{\Vcal}, \qquad \bigcap_{m \in \N} \Vcal \cap \varphi^m (\Vcal) = W^\ru (K) \cap \Vcal.
\end{equation}
Therefore, the sequence of unrevisited neighborhoods $\Vcal \cap \varphi^m (\Vcal)$ is decreasing with respect to $m \in \N$ and tends to $W^\ru (K) \cap \overline{\Vcal}$ as $m$ tends to $+ \infty$ for the geodesic distance to a compact set, in the sense that
\begin{equation}\label{eq: lemma unif convergence - limit}
\sup_{y \in \overline{\Vcal} \cap \varphi^{m} (\overline{\Vcal})} d_g \left(y, W^{\ru} (K) \cap \overline{\Vcal} \right) \underset{m \to + \infty}{\to} 0.
\end{equation}
A similar statement holds for the stable manifolds if we replace $\varphi^{m}$ by $\varphi^{-m}$.
\end{lemma}

\begin{remark}\label{lemma: clos. of stable man on unrevisited neigh.}
For $K$ and $\Vcal$ as in the previous lemma, we~deduce from the first equation in (\ref{eq: lemma unif convergence - closure of unrevisited}) that the sets $ W^\rs (K) \cap \overline{\Vcal} $ and $W^\ru (K) \cap \overline{\Vcal} $ are compact sets (as claimed implicitly in (\ref{eq: lemma unif convergence - limit})). This yields the inclusions
\[
\overline{W^\rs (K) \cap \Vcal} \subseteq W^\rs (K) \cap \overline{\Vcal} \quad \text{and} \quad \overline{W^\ru (K) \cap \Vcal} \subseteq W^\ru (K) \cap \overline{\Vcal}.
\]
\end{remark}

\begin{proof}[Proof of Lemma \ref{lemme: cv nonrevisite vers variete stable et instable}]
Fix a basic set $K$ and an unrevisited neighborhood $\Vcal$ of $K$ such that $\overline{\Vcal} \cap \Omega = K$. Note first that (\ref{eq: lemma unif convergence - limit}) can be deduced from (\ref{eq: lemma unif convergence - closure of unrevisited}) by contradiction. Indeed, let us assume by contradiction that we can find $\varepsilon > 0$ and an increasing sequence of integers $(m_i)_{i \in \N}$ satisfying $m_i \to +\infty$ as $i \to +\infty$ such that
\[
d_g (y_{m_i},z) > \varepsilon, \quad \forall z \in W^\ru (K) \cap \overline{\Vcal}, \quad \forall i \in \N,
\]
where
\[
y_{m_i} \in \overline{\Vcal} \cap \varphi^{m_i} (\overline{\Vcal}), \quad d_g (y_{m_i},W^{\ru} (K) \cap \overline{\Vcal}) = \sup_{y \in \overline{\Vcal} \cap \varphi^{m_i} (\overline{\Vcal})} d_g \left(y, W^{\ru} (K) \cap \overline{\Vcal} \right), \quad \forall i \in \N.
\]
By compactness of $M$, one can extract a subsequence $(y_{\widetilde{m_i}})_i$ of $(y_{m_i})_i$ which converges to a point $y_{\infty}$ as $i$ tends to $+\infty$. Assuming (\ref{eq: lemma unif convergence - closure of unrevisited}), we~have on one hand that
\[
y_{\infty} \in \bigcap_{i \in \N} \overline{\Vcal} \cap \varphi^{m_i} (\overline{\Vcal}) = \bigcap_{m \in \N} \overline{\Vcal} \cap \varphi^{m} (\overline{\Vcal}) = W^\ru (K) \cap \overline{\Vcal}
\]
and on the other hand that
\[
d_g (y_{\infty}, W^\ru (K) \cap \overline{\Vcal}) \geq \varepsilon.
\]
It leads to the expected contradiction. Now, it remains to prove (\ref{eq: lemma unif convergence - closure of unrevisited}). We begin by proving the direct inclusion for the first equation in (\ref{eq: lemma unif convergence - closure of unrevisited}). To do so, consider a point $y$ which belongs to $\bigcap_{m \in \N} \overline{\Vcal} \cap \varphi^m (\overline{\Vcal})$ and let us show that $y$ lies in $W^\ru (K)$. Thanks to Lemma \ref{partition of the manifold in stable manifolds}, $M$ decomposes into the unstable sets of the basic sets and thus there exists a unique basic set $K'$ such that $ y \in W^\ru (K')$. We claim that $K' = K$. By~contradiction, if $K' \neq K$, then we choose $\eta > 0$ sufficiently small so that $\{x, \: d_g(x,K') < \eta\} \cap \overline{\Vcal} = \emptyset$ and we choose $m \in \N$ sufficient large so that $\varphi^{-m} (y) \in W_{\eta}^\ru (K') \subset M \setminus \overline{\Vcal}$ (which is possible by definition of the unstable set of $K$). But, by definition of $y$, we~have for all $m \geq 0$ that $\varphi^{-m} (y) \in \varphi^{-m} (\overline{\Vcal} \cap \varphi^m (\overline{\Vcal})) \subset \overline{\Vcal}$, which leads to the expected contradiction. Therefore, we~must have $K' = K$ and thus $y \in W^\ru (K) \cap \overline{\Vcal}$. The reverse inclusion for the first equation of (\ref{eq: lemma unif convergence - closure of unrevisited}) is a consequence of the fact that $\overline{\Vcal}$ is unrevisited together with the definition of $W^\ru (K)$.

Now, it remains to prove the second equation of (\ref{eq: lemma unif convergence - closure of unrevisited}). To do so, it is enough to remark that the inclusions
\[
W^\ru (K) \cap \Vcal \subseteq \Vcal \cap \bigcap_{m \in \N} \Vcal \cap \varphi^m (\Vcal) \subseteq \Vcal \cap \bigcap_{m \in \N} \overline{\Vcal} \cap \varphi^m (\overline{\Vcal}) =\Vcal \cap W^\ru (K) \cap \overline{\Vcal} = W^\ru (K) \cap \Vcal
\]
are actually equalities thanks to the first equation of (\ref{eq: lemma unif convergence - closure of unrevisited}). This ends the proof of the lemma.
\end{proof}

Putting together the properties \eqref{P1} and \eqref{P2} we deduce the next property.

\begin{enumerate}\renewcommand{\theenumi}{{P\arabic{enumi}}}\setcounter{enumi}{2}
\item\label{P3}
If $\Vcal$ is an unrevisited neighborhood of a basic set $K$ such that $\overline{\Vcal} \cap \Omega = K$, then the sequence $(\varphi^m (\Vcal) \cap \varphi^{-m} (\Vcal))_{m \in \N}$ tends to
\[
\left(W^\rs (K) \cap \overline{\Vcal} \right) \cap \left(W^\ru (K) \cap \overline{\Vcal} \right) = (W^\rs (K) \cap W^\ru (K))\cap \overline{\Vcal} = K
\]
in the sense that
\[
\sup_{y \in \varphi^m(\overline{\Vcal}) \cap \varphi^{-m} (\overline{\Vcal})} d_g \left(y, K \right) \underset{m \to + \infty}{\to} 0.
\]
\end{enumerate}

\begin{remark} The property \eqref{P3} rests on the relation $W^\rs (K) \cap W^\ru (K) = K$ which is a consequence of Theorem \ref{thm Smale}, but whose proof does not need the strong transversality assumption. Moreover, they are both related with the local maximality of the basic set~$K$. Indeed, from \eqref{P3} we deduce
\[
K \subseteq \bigcap_{t \in \R} \varphi^t (\Vcal) \subseteq \bigcap_{m \in \Z} \varphi^m (\Vcal) = \bigcap_{m \in \N} \varphi^m (\Vcal) \cap \varphi^{-m} (\Vcal) = K.
\]
\end{remark}

\section{Escape functions for Axiom~A flows}\label{section Escape function}

Following the strategy of Faure and Sjöstrand \cite{FS}, we~will construct some function, called an escape function, which will enable us to define anisotropic Sobolev spaces on which the Lie derivative operator $-\Lie_V$ has nice spectral properties. This function is related to the construction of energy functions (also called Lyapunov functions) whose existence will be stated in paragraph \ref{subsection: energy function for the hamiltonian flow} and proved in Sections \ref{section: Dynamical proofs} and \ref{section: Compactness result and energy functions for the Hamiltonian flow}.

According to the strategy of \cite{FS}, we~need to construct an energy function on the unitary cotangent bundle $S^*M$ which is increasing along the (projected) Hamiltonian flow $\widetilde{\Phi}^t$. We choose here to split its construction into that of two energy functions which are slightly easier to build independently: one on the base manifold $M$ and one on the fibers of $S^* M$.

\subsection{Energy functions on $M$}

As recalled in Proposition \ref{prop: existence of continuous Lyapunov functions}, it is known from the work of Conley \cite{conley} that any continuous flow on a compact manifold behaves like a gradient flow outside an invariant set called the \textit{chain recurrent set} (see for instance~\cite{palis-demelo} of \cite{FH} for a definition). For example, it is the case for gradient flows of Morse functions where the chain recurrent set equals the set of hyperbolic fixed points and an energy function is given by the Morse function itself. For Axiom~A flows satisfying the no-cycle property, the chain recurrent set equals the non-wandering set---see \cite[Th.\,5.3.47]{FH}. Later, Wilson \cite{wilson}, Fathi and Pageault \cite{fathi-pageault} explained how to regularized Conley's Lyapunov function. For the sake of completeness and since its proof will be very instructive for the construction of energy functions on $S^* M$, we~will provide a smooth version to Conley's result (recalled in Theorem \ref{prop: existence of continuous Lyapunov functions}) using filtrations and unrevisited neighborhoods. Furthermore, the next proposition has the advantage to give $\Ccal^{\infty}$ Lyapunov functions with arbitrary values on the basic sets of any Axiom~A flow satisfying the no-cycle property. The proof of next Proposition is given in Section \ref{section: Dynamical proofs} and the analysis on $S^*M$ will be made in Section \ref{section: Compactness result and energy functions for the Hamiltonian flow}.

\begin{prop}[Energy function for Axiom~A flows]\label{prop energy function for Axiom A flows}
Let $\varphi^t$ be an Axiom~A flow which satisfies the no-cycle property. For every $\varepsilon > 0$ and for every family of pairwise distinct real numbers $(\lambda_i)_{1\leq i \leq N}$ compatible with the graph structure in the sense that
\[
K_i \leq K_j \implies \lambda_i \leq \lambda_j,
\]
there exist an energy function $E \in \Ccal^{\infty}(M)$, $\varepsilon$-neighborhoods $ \Ncal_i$ of $K_i$ and a constant $\eta > 0$ such that:
\[
\Lie_V E \geq 0\;\text{ on } M, \quad \text{and}\quad \Lie_V E > \eta \;\text{ on } M \setminus \biggl(\bigcup_{i = 1}^N \Ncal_i \biggr).
\]
Moreover, for all $i\in\lcr1,N\rcr$, the map $E$ is close to $\lambda_i$ on each $\Ncal_i$ in the sense that
\[
E = \lambda_i \;\text{ on } K_i \quad\text{and}\quad \sup_{x \in \Ncal_i} |E(x) - \lambda_i | < \varepsilon.
\]
\end{prop}

\emph{From now on}, the vector field $V \in \Gamma (TM)$ will be considered to be Axiom~A and to satisfy the strong transversality assumption \eqref{Transversality assumption} and in particular the no-cycle property.

\subsection{Energy functions for the Hamiltonian flow}\label{subsection: energy function for the hamiltonian flow}

Let us define the following $\widetilde{\Phi}^t$-invariant subset of $S^*M$:
\begin{align*}
\Sigma_{\ruo} &:= \bigcup_{x \in M} \kappa \bigl(E_{\rso}^* (x) \setminus 0_M \bigr),& \Sigma_{\rs}&:= \bigcup_{x \in M} \kappa \bigl(E_\ru^* (x) \setminus 0_M \bigr),\\
\Sigma_{\ru} &:= \bigcup_{x \in M} \kappa \bigl(E_\rs^* (x) \setminus 0_M \bigr),& \Sigma_{\rso}&:= \bigcup_{x \in M} \kappa \bigl(E_{\ruo}^* (x) \setminus 0_M \bigr),
\end{align*}
where $\kappa$ denotes the projection on the unitary cotangent bundle defined in \eqref{eq: def kappa}.
They will be our basic ingredients to construct energy functions on $S^*M$. Indeed, following the ideas of Faure-Sjöstrand \cite{FS} and Dang-Rivière \cite{DR0, DRI}, we~will see that $(\Sigma_{\ruo},\Sigma_{\rs})$ and $(\Sigma_{\ru},\Sigma_{\rso})$ are both a pair of repelling and attracting compact invariant sets for the Hamiltonian flow $\widetilde{\Phi}^t$. It will be enough to construct an energy function on the fiber. First, let us recall that the \textit{strong transversality assumption} implies that
\[
\Sigma_{\ruo} \cap \Sigma_{\rs} = \emptyset = \Sigma_{\ru} \cap \Sigma_{\rso}.
\]

The following lemma proved in Section \ref{subsection: proof Lemma attracting and repelling sets for the Hamiltonian flow} tells us that they are indeed attracting and repelling sets for the Hamiltonian flow:

\begin{lemma}\label{lemma attractor and repeller on the phase space}
For every $(x,\xi) \in S^* M \setminus(\Sigma_{\rs} \cup \Sigma_{\ruo})$, we~have
\[
d_{S^* M} \bigl(\widetilde{\Phi}^{t} (x,\xi), \Sigma_{\rs} \bigr) \underset{t \to +\infty}{\to } 0 \quad\text{and}\quad
d_{S^* M} \bigl(\widetilde{\Phi}^{-t} (x,\xi), \Sigma_{\ruo} \bigr) \underset{t \to +\infty}{\to } 0.
\]
Similarly, for every $(x,\xi) \in S^* M \setminus \left(\Sigma_{\rso} \cup \Sigma_{\ru} \right)$, we~have
\[
d_{S^* M} \bigl(\widetilde{\Phi}^{t} (x,\xi), \Sigma_{\rso} \bigr) \underset{t \to +\infty}{\to } 0 \quad\text{and}\quad
d_{S^* M} \bigl(\widetilde{\Phi}^{-t} (x,\xi), \Sigma_{\ru} \bigr) \underset{t \to +\infty}{\to } 0.
\]
\end{lemma}

Moreover, contrary to Anosov flows for which it is rather immediate, we~need to make sure that these sets are compact sets. The next proposition is similar to the compactness result of Dang and Rivière \cite[Lem.\,3.7, p.\,15]{DR0}. Its proof is given in Section~\ref{subsection: proof compactness result}.

\begin{prop}[Compactness]\label{prop compactness}Let $\varphi^t$ be an Axiom~A flow which satisfies the strong transversality assumption \eqref{Transversality assumption}. Then, the subsets
$\Sigma_{\ru}$, $\Sigma_{\ruo}$, $\Sigma_{\rso}$, $\Sigma_{\rs}$ of $S^*M$ are $\widetilde{\Phi}^t$-invariant compact sets.
\end{prop}

To construct energy functions, we~also need the existence of arbitrarily small stable neighborhoods. The proof of next lemma is given in Section  \ref{subsection: Lemma invariant neighborhoods for the Hamiltonian flow}.

\begin{lemma}[Invariant neighborhoods]\label{lemma invariant neighborhoods}
For every $\varepsilon > 0$, there exist $\varepsilon$-neighborhoods $\Ucal^{\rs/\rso}$ (\resp $\Ucal^{\ru/\ruo}$) of $\Sigma_{\rs/\rso}$ (\resp $\Sigma_{\ru/\ruo}$) which are $\widetilde{\Phi}^{1}$-stable (\resp $\widetilde{\Phi}^{-1}$-stable).
\end{lemma}

As a consequence of these three results, we~obtain energy functions on the fiber of~$S^* M$. The proof of the following proposition is given in Section \ref{subsection: Proof energy functions for the Hamiltonian flow}.

\skpt
\begin{prop}[Energy functions for the Hamiltonian flow]\label{prop energy function on the unitary cotangent bundle}
Let $\varphi^t$ be an Axiom~A flow which satisfies the strong transversality assumption \eqref{Transversality assumption}. For every $\varepsilon > 0$, there exist energy functions $E_{\pm} \in \Ccal^{\infty}(S^* M; [0,1])$, $\varepsilon$\nobreakdash-neighborhoods $\Wcal^{\rs/\rso}$ of $\Sigma_{\rs/\rso}$, $\Wcal^{\ruo/\ru}$ of $\Sigma_{\ruo/\ru}$ and a constant $\eta > 0$ such that:
\begin{align*}
\begin{split}
\Lie_{X_H} E_{+} &\geq 0 \text{ on } S^* M \: \text{and } \Lie_{X_H} E_{+} > \eta \text{ on } S^* M \setminus (\Wcal^{\ruo} \cup \Wcal^{s}), \\
\Lie_{X_H} E_{-} &\geq 0 \text{ on } S^* M \: \text{and } \Lie_{X_H} E_{-} > \eta \text{ on } S^* M \setminus (\Wcal^{\ru} \cup \Wcal^{\rso}).
\end{split}
\end{align*}
Moreover, the map $E_{\pm}$ are constant on each $\Sigma_{*}$ and we have the estimate
\begin{align*}
\sup_{(x,\xi) \in \Wcal^{\ru/\ruo}} \left| E_{\pm}(x,\xi) - 0 \right| \leq \varepsilon \quad \text{and} \quad \sup_{(x,\xi) \in \Wcal^{\rs/\rso}} \left| E_{\pm}(x,\xi) - 1 \right| \leq \varepsilon.
\end{align*}
\end{prop}

\subsection{The escape functions}

Next, we~give a global escape function which extends the construction of Dyatlov and Guillarmou \cite{dyatlov-guillarmou_2016} to the whole manifold and coincide with the one of Dang and Rivière \cite{DR0} for Morse-Smale gradient flows and the one of Faure-Sjöstrand \cite{FS} for Anosov flows. The proof of next proposition is given in Section \ref{section: Proof of escape function}.

\begin{prop}[Escape function]\label{prop escape function}
Let $u,s,n_0 \in \R$ be such that $u < 0 \leq n_0 < s$. There exists a smooth function $m(x,\xi) \in C^{\infty} (T^* M)$ called an \emph{order function} and an \emph{escape function} $G_m \in \Ccal^{\infty} (T^*M)$ defined by:
\[
G_{m} (x, \xi) = m(x, \xi) \log \sqrt{1 + f(x,\xi)^2}
\]
where $f \in C^{\infty}(T^* M)$ is positive everywhere and homogeneous of degree $1$ in $\xi$ as soon as $|\xi| \geq 1$, and where $m$ is defined by $m(x,\xi) = \chi (|\xi|^2) \E(x,\sfrac{\xi}{|\xi|})$ with $\chi$ being a smooth cut-off function such that $\chi = 0$ on $]-\infty, 1/2]$, $\chi = 1$ on $[1,+\infty[$, $\chi \geq\nobreak 0$ everywhere as in Figure \ref{figure chi}, and $\E$ being a linear combination of previous energy functions:
\[
\E (x,\xi):= -E(x) + 2s + (2u - n_0) E_+ (x,\xi) + (n_0 - 2s) E_- (x,\xi).
\]
Moreover, we~have the following estimates:
\begin{enumerate}
\item\label{prop escape function1} There exist conical neighborhoods $\widetilde{\Ncal}^{\rs/\ro/\ru}$ of $\bigcup_{z \in M} E_{\rs/\ro/\ru}^* (z) \setminus 0_M$ such that $f = |\xi(V)|$ on $\widetilde{\Ncal}^\ro$ and for $|\xi| \geq 1$,
\[
\frac{1}{2} n_0 \leq m \leq 2n_0 \quad\text{on } \widetilde{\Ncal}^{\ro},
\]
\[
m \geq s \;\text{ on } \widetilde{\Ncal}^{\rs}, \quad \text{and} \quad m \leq u \;\text{ on } \widetilde{\Ncal}^{\ru}.
\]
Also, the open sets can be chosen arbitrarily close to the invariant distributions $E_\rs^*$, $E_\ro^*$ and $E_\ru^*$ as in Proposition \ref{prop energy function on the unitary cotangent bundle}.
\item\label{prop escape function2} The map $G_m$ is strictly decreasing along the flow $\Phi^t$ except at points $(x,\xi)$ where~$|\xi|$ is small or where~$x$ is in a small neighborhood of the non-wandering set and~$\xi$ is in a conical neighborhood of $E_\ro^*$: for all $i\in\lcr1,N\rcr$, there exist an open neighborhood~$\Ncal_i$ of~$K_i$ and a radius $R > 0$ such that for all $(x,\xi)\in \bigcup_{z\in \cup N_i} T_z^*M\setminus \widetilde{\Ncal}^\ro$ such that $ |\xi| \geq R$ \emph{and} for all $(x,\xi) \in \bigcup_{z\notin \oldbigcup_i \Ncal_i} T_z^*M$ such that $ |\xi| \geq R$
\[
X(G_{m}) (x,\xi) < -c \min (s,|u|) =: - C_m < 0,
\]
with $c > 0$ being independent of the constants $u, n_0, s$ and of the size of the conical neighborhoods.
\item\label{prop escape function3} More generally, for every $(x,\xi) \in T^* M$ such that $|\xi| \geq R$, we~have
\[
X(G_{m}) (x, \xi) \leq 0.
\]
\end{enumerate}
\end{prop}

\begin{figure}[!t]
\centering
\begin{tikzpicture}[line cap=round,line join=round,>=triangle 45,x=1.6cm,y=1.6cm]
\draw[->,color=black] (-0.5,0) -- (2,0);
\draw[->,color=black] (0,-0.5) -- (0,2);
\draw[shift={(1,0)},color=black] (0pt,2pt) -- (0pt,-2pt) node[below] {\footnotesize $1$};
\draw[shift={(0.5,0)},color=black] (0pt,2pt) -- (0pt,-2pt) node[below] {\footnotesize $\frac{1}{2}$};
\draw[shift={(0,1)},color=black] (2pt,0pt) -- (-2pt,0pt) node[left] {\footnotesize $1$};
\draw[color=black] (0pt,-10pt) node[left] {\footnotesize $0$};
\clip(-0.5,-0.5) rectangle (2,2);
\draw [domain=-0.5:0.5] plot(\x,{(-0-0*\x)/-0.5});
\draw [domain=1.0:2.0] plot(\x,{(--0.5-0*\x)/0.5});
\draw [shift={(0.5,0.2)}] plot[domain=4.69:5.32,variable=\t]({1*0.2*cos(\t r)+0*0.2*sin(\t r)},{0*0.2*cos(\t r)+1*0.2*sin(\t r)});
\draw [shift={(0.49,0.2)}] plot[domain=5.38:6.11,variable=\t]({1*0.21*cos(\t r)+0*0.21*sin(\t r)},{0*0.21*cos(\t r)+1*0.21*sin(\t r)});
\draw [shift={(1.01,0.8)}] plot[domain=2.23:2.97,variable=\t]({1*0.21*cos(\t r)+0*0.21*sin(\t r)},{0*0.21*cos(\t r)+1*0.21*sin(\t r)});
\draw [shift={(1,0.8)}] plot[domain=1.55:2.18,variable=\t]({1*0.2*cos(\t r)+0*0.2*sin(\t r)},{0*0.2*cos(\t r)+1*0.2*sin(\t r)});
\draw (0.7,0.17)-- (0.75,0.5);
\draw (0.8,0.83)-- (0.75,0.5);
\draw (0.42,0.73) node[anchor=north west] {$\chi $};
\end{tikzpicture}
\caption{Cut-off function $\chi$.} \label{figure chi}
\end{figure}

\section{Anisotropic Sobolev spaces}\label{section Anisotropic Sobolev spaces}

The purpose of this section is to construct some Hilbert spaces in which the operator $-\Lie_V^{(k)}$ acting on sections of $\Ecal_k:= \Lambda^k T^*M \otimes \C$ has good spectral properties. For $k = 0$ and for Anosov flows, Faure and Sjöstrand defined an anisotropic Sobolev space which can be roughly written as $\exp (G_m (x,-iD))^{-1} L^2 (M ; \C)$ using the escape function $G_m$ of Proposition \ref{prop escape function}. This construction can be extended to the fiber bundle $ \Lambda^k T^*M \otimes \C$ of differential forms with complex values for every $k\in\lcr0,n\rcr$ as explained in \cite{dyatlov-zworski}. So, we~construct a pseudodifferential operator acting on sections of $\Ecal_k$ which has for principal symbol\enlargethispage{\baselineskip}
\[
A_m^{(k)} (x,\xi) = \exp (G_m (x,\xi))\Id_{\Ecal_k (x)} \in \Ccal^{\infty} \bigl(T^*M, \Hom (\pi^* \Ecal_k, \pi^* \Ecal_k) \bigr),
\]
with $G_m$ being the escape function obtained in Proposition \ref{prop escape function} and were $\pi^* \Ecal_k \to T^* M$ denotes the pull-back bundle by the projection $\pi: T^* M \to M$. We denote this symmetric\footnote{This pseudodifferential operator is defined using Weyl quantization in a coordinate system. This quantization is known for mapping a real symbol to a symmetric operator. We refer the reader to the book \cite[Chap.\,4, Th.\,4.1]{zworski} for Weyl quantization on $\R^n$ and to the book \cite[App.\,E]{DZ-book} for Weyl quantization on a variety.} pseudodifferential operators by $\A_m^{(k)}: \Omega^k (M ; \C) \subset L^2 (M; \Ecal_k) \to \Omega^k (M ; \C) $ and we refer for instance to \cite[Chap.\,4, p.\,56]{zworski} for a definition of the Weyl quantization on $\R^n$ and to \cite[App.\,C.1, p.\,29]{dyatlov-zworski},\cite[\S 9.2 p.\,40]{DRI} for pseudodifferential operators on manifolds and vector bundles. Furthermore, the symbols $A_m^{(k)} (x,\xi) $ belong to a class of symbols with variable order whose properties are discussed in the appendix of \cite{FRS} and they are \emph{elliptic} in the sense of \cite[Def.\,8, p.\,40]{FRS}. It implies the existence of a smoothing operator $\hat{r}: \D'^{,k} (M;\C) \to \Omega^k (M; \C)$ such that $\widetilde{\A}{}_m^{(k)}:= \A_m^{(k)} + \hat{r}: \Omega^k (M; \C) \to \Omega^k (M; \C)$ is formally self-adjoint, elliptic and invertible---see \cite[Lem.\,12, p.\,42]{FRS}. Thus, we~choose $(\widetilde{\A}{}_m^{(k)})^{-1}$ to be a representative of the inverse of~$\A_m^{(k)}$ modulo smoothing operators by setting $(\widetilde{\A}{}_k^{(k)})^{-1}:= (\A_m^{(k)} + \hat{r})^{-1}$ and we define, following \cite{FRS} and \cite{FS}, the anisotropic Sobolev space
\[
\Hil_k^{m}:= (\widetilde{\A}{}_m^{(k)})^{-1} L^2 (M; \Ecal_k), \qquad \w{u,v}_{\Hil_k^m}:= \w{\widetilde{\A}{}_m^{(k)} u, \widetilde{\A}{}_m^{(k)} v}_{L^2 (M;\Ecal_k)}.
\]
The space $\Hil_k^m$ endowed with $\w{.,.}_{\Hil_k^m}$ is isometric to the space $L^2 (M; \Ecal_k)$. Moreover, it is isomorphic\footnote{The idea is that any current $u \in \D'^{,k}(M)$ writes in coordinates as a $k$-form with coefficients in $\D'(M)$, \ie $u = \sum u_{i_1,\dots, i_k} dx^{i_1} \wedge \cdots \wedge dx^{i_k}$ where $u_{i_1,\dots, i_k} \in \D' (M)$. A partition of unity argument gives the result.} to the space $\Hil_0^m \otimes \Omega^k (M; \C)$, and the following inclusions
\[
\Omega^k (M; \C) \subset \Hil_k^m \subset \D'^{,k} (M ;\C).
\]
are continuous, where $\D'^{,k} (M; \C)$ is endowed with the weak topology---see \cite{schwartz}.

\subsection{Spectral properties}

Adapting the proof of \cite[Th.\,1.4]{FS} to the case of vector bundles and using the properties\footnote{These are the only properties used in the proof of \cite{FS}.} of the escape function stated in Proposition \ref{prop escape function} (which is the exact analogue of Lemma $1.2$ in \cite{FS}), one can establish the existence of a discrete spectrum on these anisotropic Sobolev spaces:

\begin{thm}[Discrete spectrum]\label{thm: discrete sprectrum} Let $\varphi^t$ be an Axiom~A flow satisfying the strong transversality assumption (\ref{Transversality assumption - first}). Let $G_m$ be an escape function. For every $k\in\lcr0,n\rcr$, the operator $-\Lie_V^{(k)}$ defines a maximal closed unbounded operator on $\Hil_k^m$ with domain
\[
\D (-\Lie_V^{(k)}) = \{u \in \Hil_k^m: \: -\Lie_V^{(k)} u \in \Hil_k^m \} \: \text{ and } \: \Lie_V^{(k)}: \D (-\Lie_V^{(k)}) \subset \Hil_k^m \to \Hil_k^m.
\]
Moreover, there exists a constant $C_0 \in \R$ (which depends on the choice of the escape function $G_m$) such that
\[
-\Lie_V^{(k)} \text{ has empty spectrum on } \re (z) \geq C_0,
\]
and there exists a constant $C_1 > 0$ (which only depends on the vector field $V$ and the metric $g$) such that
\[
-\Lie_V^{(k)} \text{ has discrete spectrum on } \re (z) \geq -C_m + C_1,
\]
where $C_m>0$ is the constant given by Proposition \ref{prop escape function}.
\end{thm}

\begin{figure}[!t]
\begin{center}
\definecolor{uququq}{rgb}{0.25,0.25,0.25}
\definecolor{zzzzzz}{rgb}{0.6,0.6,0.6}
\definecolor{qqqqff}{rgb}{0,0,1}
\definecolor{ttttff}{rgb}{0.2,0.2,1}
\begin{tikzpicture}[line cap=round,line join=round,>=triangle 45,x=0.4cm,y=0.3cm]
\draw[->,color=black] (-14,0) -- (10,0);
\draw[->,color=black] (0,-8) -- (0,8);
\clip(-14,-8) rectangle (10,8);
\fill[line width=0pt,dash pattern=on 2pt off 2pt,color=zzzzzz,fill=zzzzzz,fill opacity=0.1] (-14,-8) -- (-14,8) -- (-8,8) -- (-8,-8) -- cycle;
\draw [dash pattern=on 1pt off 1pt] (6,-8) -- (6,8);
\draw [dash pattern=on 1pt off 1pt] (-8,-8) -- (-8,8);
\draw (-8.63,-0.08) node[anchor=north west] {$-C_m + C_1$};
\draw (5.74,0) node[anchor=north west] {$C_0$};
\draw (-14.20,-6.22) node[anchor=north west] {\scriptsize{$\text{Essential spectrum}$}};
\draw (-2.9,-6.22) node[anchor=north west] {\scriptsize{$\text{Discrete spectrum}$}};
\draw (5.8,-6.22) node[anchor=north west] {\scriptsize{$\text{No spectrum}$}};
\draw (0,0) node[anchor=north west] {$0$};
\begin{scriptsize}
\draw [color=black] (6,0)-- ++(-2.0pt,0 pt) -- ++(4.0pt,0 pt) ++(-2.0pt,-2.0pt) -- ++(0 pt,4.0pt);
\fill [color=qqqqff] (0,0) circle (1.5pt);
\fill [color=qqqqff] (1.48,-1.27) circle (1.5pt);
\fill [color=qqqqff] (4,-2) circle (1.5pt);
\fill [color=qqqqff] (5.04,-3.45) circle (1.5pt);
\fill [color=qqqqff] (3.4,-7.12) circle (1.5pt);
\fill [color=qqqqff] (-1.36,-3.13) circle (1.5pt);
\fill [color=qqqqff] (-1.99,-5.16) circle (1.5pt);
\fill [color=qqqqff] (-4.11,-2.97) circle (1.5pt);
\fill [color=qqqqff] (-6.56,-3.98) circle (1.5pt);
\fill [color=qqqqff] (-4.88,-6.32) circle (1.5pt);
\fill [color=qqqqff] (-6.56,3.98) circle (1.5pt);
\fill [color=qqqqff] (-4.11,2.97) circle (1.5pt);
\fill [color=qqqqff] (-1.36,3.13) circle (1.5pt);
\fill [color=qqqqff] (-4.88,6.32) circle (1.5pt);
\fill [color=qqqqff] (-1.99,5.16) circle (1.5pt);
\fill [color=qqqqff] (1.48,1.27) circle (1.5pt);
\fill [color=qqqqff] (4,2) circle (1.5pt);
\fill [color=qqqqff] (5.04,3.45) circle (1.5pt);
\fill [color=qqqqff] (3.4,7.12) circle (1.5pt);
\draw [color=uququq] (-8,0)-- ++(-2.0pt,0 pt) -- ++(4.0pt,0 pt) ++(-2.0pt,-2.0pt) -- ++(0 pt,4.0pt);
\end{scriptsize}
\end{tikzpicture}
\caption{Illustration of Pollicott-Ruelle resonances of $-\Lie_V^{(k)}$ on the anisotropic Sobolev space $\Hil_k^m$. The fact that $0$ is a resonance or not depends on $k$.}
\end{center}
\end{figure}

The eigenvalues of $-\Lie_V^{(k)}$ on the anisotropic Sobolev space are called the \textit{Pollicott-Ruelle resonances} of $-\Lie_V^{(k)}$. There are many (equivalent) definitions of resonances. In~particular, they can be viewed as the poles of the resolvent operator $(-\Lie_V^{(k)} - z)^{-1}: \Omega^k (M; \C) \to \D'^{,k} (M; \C)$. Let us make a few remarks about them:

\skpt
\begin{remark}
\begin{itemize}
\item The discrete spectrum of $-\Lie_V^{(k)}$ is intrinsic in the sense that it does not depend on the escape function, and the essential spectrum can be chosen as far as we want to the origin by taking $m$ such that $C_m \gg 1$. We refer to \cite[Th.\,1.5, p.\,134]{FS} for a proof in the case of Anosov vector fields.
\item The set of resonances is symmetric along the real axis since the vector field $V$ is real.
\item When $k = 0$, the resonances are included in the set $\{ \re (z) \leq 0 \}$ and the point $z = 0$ is a resonance since the constants are solutions of $\Lie_V u = 0$. This fact is not true in general for $k > 0$ and the optimal constant $h \in \R$ such that there is no spectrum in the set $\re (z) > h$ is related to the \textit{topological entropy} of the basic sets.
\item From the previous remark, we~can see that the resonance $0$ is somehow related to Morse inequalities:
\[
\dim (\Ker(-\Lie_V)) \geq b_0 = \dim \Hsc_0 (M),
\]
where $b_0(M)$ is the number of connected components of $M$ and where $\Ker (- \Lie_V)$ denotes the kernel of the Lie derivative viewed as an unbounded operator on $\Hil_0^m$.
\item From the first point, we~can deduce that the space $\Ker ((-\Lie_V^{(k)})^{\ell})$ does not depend on the space function $m$ for any $\ell \in \N$ (provided $m$ is chosen such that $-C_m+C_1<0$).
\end{itemize}
\end{remark}

Before going deeper into topological considerations, let us recall some useful properties of the operators $-\Lie_V^{(k)}$.

\skpt
\begin{remark}
\begin{itemize}
\item When proving the discrete spectrum theorem for Anosov vector fields, precisely in \cite[Lem.\,3.3, p.\,343]{FS}, Faure and Sjöstrand obtained a bound on the resolvent operator which remains true in our context. For every $z$ such that \hbox{$\re (z) \!>\! C_0$}, we~have
\[
\| (\Lie_V^{(k)} + z)^{-1} \|_{\Hil_k^m \to \Hil_k^m} \leq \frac{1}{\re (z) - C_0}.
\]
An application of the Hille-Yosida theorem \cite[Cor.\,3.6, p.\,76]{engel-nagel} yields that
\[
(\varphi^{-t})^*: \Hil_k^m \to \Hil_k^m, \quad \forall t \geq 0,
\]
generates a \emph{strongly continuous semi-group} whose norm is bounded by $e^{C_0 t}$. Therefore, for every $z$ such that $\re (z) > C_0$ we can write the resolvent as follows:
\[
(\Lie_V^{(k)} + z)^{-1} = \int_0^{+\infty} e^{-zt } (\varphi^{-t})^* dt: \Hil_k^m \to \Hil_k^m,
\]
where the integral converges absolutely. Note that it is convenient to use the convention $(\varphi^{-t})^* = e^{-t \Lie_V^{(k)}}: \Hil_k^m \to \Hil_k^m$.
\item If $z_0 $ is any resonance of $-\Lie_V^{(k)}$ such that $\re (z_0) > -C_m + C_1$, then we can define the Riesz projector
\[
\pi_{z_0}^{(k)}:= \frac{1}{2 i \pi} \int_{\gamma_{z_0}} (\Lie_V^{(k)} + z)^{-1} dz: \Hil_k^m \to \Hil_k^m,
\]
where the integral is over a positively oriented closed curved $\gamma_{z_0}$ which surrounds the resonance $z_0$ and no other resonances. Moreover, it commutes with $\Lie_V^{(k)}$ and it has finite rank. Note that this definition still makes sense when $z_0$ is not a resonance and in that case, $\pi_{z_0}^{(k)}$ is identically $0$. Note also that $\pi_{z_0}^{(k)}$ commutes with the exterior derivatives $d$ because $\Lie_V^{(k)}$ commutes with $d$ thanks to the Cartan formula.
\Changel
\item The resolvent operator writes as a Laurent series near $z_0$:
\[
(\Lie_V^{(k)} + z)^{-1} = \sum_{l = 1}^{m_k (z_0)} (-1)^{l-1} \frac{(\Lie_V^{(k)} + z)^{l-1} \pi_{z_0}^{(k)}}{(z - z_0)^{l}} + R_{z_0,k} (z),
\]
where $R_{z_0,k}$ is holomorphic near $z_0$.
\end{itemize}
\Changelback
\end{remark}

\section{Construction of the Morse-de Rham complex} \label{section topology}

Let us define $\Res_k (V)$ as the set of resonances $z \in \C$ of the operator $-\Lie_V^{(k)}$, \ie the set of points $z_0\in \C$ such that we can find an escape function $G_m$ with $\re (z_0) > -C_m+C_1$ and such that the algebraic multiplicity of $z_0$, denoted by $m_k (z_0)$, satisfies $m_k (z_0) \neq 0$. We can then define $C_V^k(z_0)$ as the range of the projector $\pi_{z_0}^{(k)}$ defined on the space of $k$-forms $\Omega^k (M;\C)$. Equivalently, we~have
\[
C_V^k (z_0) = \Ker \bigl((\Lie_V^{(k)} + z_0)^{m_k (z_0)} \bigr),
\]
and since $\pi_{z_0}^{(k)}$ has finite rank, the vector space $C_V^k (z_0)$ is \emph{finite dimensional}. Recall that this space is independent of the choice of the escape function used to define Hilbert spaces---see \cite[Th.\,1.5 p.\,334]{FS}. Also, the exterior derivative $d$ maps the space $C_V^k (z_0)$ into $C_V^{k+1} (z_0)$ because $d$ commutes with the spectral projector $\pi_{z_0}^{(k)}$, a property which follows from the fact that $d$ commutes with $\Lie_V^{(k)}$ thanks to Cartan's formula.

We will denote by $\Hg (C_V^* (0),d)$ the cohomology of the spectral complex associated with the eigenvalue $0$ and by $\Hg^k (M; \C)$ the complex $k$-th de~Rham cohomology:
\begin{align*}
\Hg^k (C_V^* (0),d) &= \faktor{ \Ker (\rest{d}{C_V^k (0)}) }{\Ran (\rest{d}{C_V^{k-1} (0)}) }, \\
\Hg^k (M; \C) &= \faktor{\Ker (\rest{d}{\Omega^k (M;\C)}) }{ \Ran (\rest{d}{\Omega^{k-1} (M;\C)})}.
\end{align*}

Theorem \ref{thm intro} of the introduction is a consequence of the next theorem.

\begin{thm}\label{thm isomorphism}
Suppose that the vector field V generates an Axiom~A flow which satisfies the strong transversality assumption \eqref{Transversality assumption - first}. For every integer $k\in\lcr0,n\rcr$, the map
\[
\pi_0^{(k)}: \Omega^k (M ;\C) \to C_V^k (0)
\]
induces an isomorphism between $\Hg^k (M; \C)$ and $\Hg^k (C_V^* (0),d)$.
\end{thm}

The proof of this theorem is given in \cite[Th.\,2.1]{DRIII}. It was actually proved in the case of Morse-Smale flows but the proof also works for Axiom~A flows since it rests on algebraic arguments. Its starting point is the commutation formula $\pi_0^{(k+1)} \circ d = d \circ \pi_0^{(k)}$ which is a consequence of the fact that the Lie derivative commutes with the exterior derivative, due to the Cartan formula. Then, the (quasi-)isomorphism of the statement can be deduced from the following de~Rham theorem.

\begin{thm}[de~Rham]\label{thm: de Rham} The following statements are true:
\begin{enumerate}
\item Let $u$ be an element in $\Hil_k^m$ satisfying $d u = 0$. There exists $\omega \in \Omega^k (M; \C)$ such that $u - \omega \in d\left(\Hil_{k-1}^{m+1} \right)$.
\item If $u = d v$ with $u \in \Omega^k (M; \C)$ and $v \in \D'^{,k-1} (M; \C)$, then there exists $\omega \in \Omega^{k-1} (M; \C)$ such that
\[
u = d \omega.
\]
\end{enumerate}
\end{thm}

This result can be found in \cite[p.\,355]{schwartz} or in \cite[p.\,16]{DRIII} for this version using anisotropic Sobolev spaces.

\label{section: proof of isomorphism theorem}

\section{Energy functions and application to Axiom A flows}\label{section: Dynamical proofs}

\subsubsection{A useful lemma}\label{subsubsection: a useful lemma} In this part, we~recall a general analysis introduced in \cite{FS} which will be applied to both flows $\varphi^t: M \to M$ and $\widetilde{\Phi}^t: S^* M \to S^* M$.

Let us consider some smooth vector field $v \in \Gamma (TX)$ on a compact Riemannian manifold $(X,g_0)$ and let us denote by $\exp (t\cdot v)$ the flow generated by $v$. A pair of compact sets $(K_+,K_-)$ is said to be \emph{attractor-repeller} for the flow $\exp (t\cdot v)$ on $X$ if it satisfies the two following conditions:
\begin{enumeratei}
\item $\forall x \in X \setminus (K_- \cup K_+)$, $d_{g_0}(\exp (t\cdot v) (x), K_{\pm}) \underset{t \to \pm \infty}{\dpl\to} 0$.
\item There exist open neighborhoods $\Vcal_{\pm}$ of $K_{\pm}$ stable by $\exp (\pm 1\cdot v)$ such that $\overline{\Vcal_{-}} \cap \overline{\Vcal_{+}} = \emptyset$.
\end{enumeratei}

\begin{figure}[!t]
\centering
\def\svgscale{0.7}
\includesvg{attractor-repeller}
\caption{Illustration of an attractor-repeller system.}
\label{fig: attractor-repeller model}
\end{figure}

\begin{lemma}[{Faure-Sjöstrand, \cite{FS}}]\label{lemma: energy function for attractor-repeller model}
Let $(K_-,K_+)$ be an attractor-repeller pair for the flow $\exp (t\cdot v)$ on the compact Riemannian manifold $(X,g_0)$ and fix $\varepsilon > 0$. There exist $\varepsilon$-neighborhoods $\Wcal_{\pm}$ of $K_{\pm}$, an energy function $m \in \Ccal^{\infty} (X; [0,1])$ and a constant $\eta > 0$, which depends on $\varepsilon$, such that $v (m) \geq 0$ everywhere and $v (m) > \eta$ outside $\Wcal_{-} \cup \Wcal_{+}$. Moreover, we~have $m > 1 - \varepsilon$ on $\Wcal_{+}$ and $m = 1$ on $K_+$. Similarly, we~have $m < \varepsilon$ on $\Wcal_{-}$ and $m = 0$ on $K_-$.
\end{lemma}

This lemma is proved in a way similar to Lemma $2.1$ of Faure-Sjöstrand \cite[p.\,336]{FS}.

\subsection{Energy functions for Axiom~A flows}

In this part, we~explain how to construct an energy function for the flow $\varphi^t$ from the previous lemma. More precisely, we~prove Proposition \ref{prop energy function for Axiom A flows}. The construction of an energy function presented here splits in three parts:
\begin{itemize}
\item First, for every total order relation in the sense of Section \ref{subsubsection total order relation} and for every $\varepsilon > 0$, we~define an $\varepsilon$-filtration from the unrevisited neighborhoods. It is stronger to the existence of a filtration given by Corollary \ref{cor: existence of filtrations} since it holds for every total order relation and not only for the one induced by the continuous Lyapunov function of Theorem \ref{prop: existence of continuous Lyapunov functions}. It will enable us to choose the value of the Lyapunov function on the basic sets, as in the statement of Proposition \ref{prop energy function for Axiom A flows}.
\item Then, we~prove that for any total order relation and for every $j\in\lcr2,N\rcr$, the pair $(\bigcup_{k \geq j}W^\ru (K_k), \bigcup_{i < j}W^\rs (K_i))$ is an attractor-repeller pair.
\item Finally, as a consequence of Lemma \ref{lemma: energy function for attractor-repeller model} we obtain a family of energy functions~$E_i$ and a linear combination of them gives the global energy function for $\varphi^t$.
\end{itemize}

\begin{lemma}[Invariant neighborhoods on the base]\label{lemma: attractor and repeller for the flow on the basis}
For every total order relation given in Section \ref{subsubsection total order relation} and for every $j\in\lcr2,N\rcr$, $\bigcup_{i < j}W^\rs (K_i)$ and $\bigcup_{k \geq j}W^\ru (K_k)$ are disjoint invariant compact sets such that:
\begin{equation}\label{eq: lemma invariant neigh. on the base - eq 1}
\forall x \notin \bigcup_{i < j}W^\rs (K_i), \quad d (\varphi^t (x), \bigcup_{k \geq j}W^\ru (K_k)) \underset{t \to+\infty}{\to} 0
\end{equation}
and
\begin{equation}\label{eq: lemma invariant neigh. on the base - eq 2}
\forall x \notin \bigcup_{k \geq j}W^\ru (K_k), \quad d (\varphi^{-t} (x),\bigcup_{i < j}W^\rs (K_i)) \underset{t \to+\infty}{\to} 0.
\end{equation}
\end{lemma}

\begin{proof}
The fact that these sets are disjoint and compact is a direct consequence of the order relation's properties. Now, consider $x \notin \bigcup_{i < j}W^\rs (K_i)$. From the decomposition (\ref{partition of the manifold in stable manifolds}) of $M$ into stable manifolds of the basic sets, there exists $k\in\lcr1,N\rcr$ such that $x \in W^\rs(K_k)$. Thanks to our choice of a total order relation in the sense of Section~\ref{subsubsection total order relation}, we~necessarily have $k \geq j$. This proves the convergence (\ref{eq: lemma invariant neigh. on the base - eq 1}). Up to replacing the flow $\varphi^t$ by $\varphi^{-t}$, we~also get the convergence (\ref{eq: lemma invariant neigh. on the base - eq 2}).
\end{proof}

Now, the key point is to prove the second property of an attractor-repeller. This is given by the next lemma which is slightly more precise, is a corollary of Theorem \ref{prop: existence of continuous Lyapunov functions}, and is inspired by the filtration lemma of Robbin \cite[7.9]{robbin}

\begin{lemma}\label{lemma: equivalence filtration and unrevisited}
For every total order relation (Section \ref{subsubsection total order relation}) and for every $\varepsilon >0$, there exists a filtration $(\Ocal_{\ell}^- (\varepsilon))_{1\leq \ell \leq N}$ on $M$ for $\varphi^{-1}$ which is within a distance $\varepsilon$ of the stable manifolds and there exists a filtration $(\Ocal_{\ell}^+ (\varepsilon))_{1\leq \ell \leq N}$ on $M$ for $\varphi^{1}$ which is within a distance $\varepsilon$ of the unstable manifolds, in the sense that $\forall j \in \lcr1,N\rcr$,
\[
\sup_{y \in \Ocal_j^- (\varepsilon)} d_g \Bigl(y, \bigcup_{i \leq j} W^\rs (K_i) \Bigr) < \varepsilon, \qquad \sup_{y \in \Ocal_{N-j+1}^+ (\varepsilon)} d_g \Bigl(y, \bigcup_{k \geq j} W^\ru (K_k)\Bigr) < \varepsilon.
\]
\end{lemma}

The proof rests on the next sub-lemma which enables us to construct the filtration by induction from a total order relation on the basic sets. This sub-lemma will also be used in the proof of Proposition \ref{prop energy function on the unitary cotangent bundle} in order to lift the filtration on $S^* M$.

\begin{sublemma}[Uniform traveling time for unstable annulus]\label{Sublemma: uniform traveling time for unstable annulus}
Let $\Vcal$ be an unrevisited neighborhood of a basic set $K$ such that $\Vcal \cap \Omega = K$. Let $O$ be an open set such that $\Vcal \subset M\setminus O$, $\varphi^1 (\overline{O}) \subset O$ and such that $O$ is a neighborhood of $\{K_j, \: K < K_j \}$. For every $m \in \N$, we~define the \textit{annulus} $\Acal (m)$ by
\[
\Acal (m):= \varphi^m (\overline{\Vcal}) \cap \overline{\Vcal} \setminus \varphi^{-1}(\Vcal).
\]
If $K$ is not an attractor, \ie $W^\ru (K) \setminus K \neq \emptyset$, then $O \neq \emptyset$, the compact set $\Acal (m)$ satisfies $\Acal (m) \neq \emptyset$, $\Acal(m) \cap W^\ru (K) = \Acal(0) \cap W^\ru (K)$ for every integer $m \geq 0$ and
there exist $k_0, m_0 \in \N$ such that for every $m\geq m_0$
\[
\varphi^{k_0} \left(\Acal (m) \right) \subset O.
\]
\end{sublemma}

\begin{proof}We split the proof in $4$ steps.

\Changel
\subsubsection*{Step 1} We show that for every $m \in \N$ we have $\Acal(m) \neq \emptyset$ and $W^\ru (K) \cap \Acal(m) = W^\ru (K) \cap \Acal(0)$. Since the relation $W^\ru (K) \cap \varphi^{m} (\overline{\Vcal}) \cap \overline{\Vcal} = W^\ru (K) \cap \overline{\Vcal}$ is satisfied, we~must have
\[
W^\ru (K)\cap \Acal (m) = W^\ru (K) \cap \Acal (0) = \{ x \in W^\ru (K) \cap \overline{\Vcal}, \: \varphi^{1} (x) \notin \Vcal \}
\]
for all $m \in \N$. Let us prove that the last set is non-empty. If $x$ belongs to \hbox{$W^\ru (K) \cap \overline{\Vcal} \setminus K$}, which is non-empty as $K$ is not an attractor, then one can find an integer $l_0 \geq 0$ such that $\varphi^{l_0}(x) \in W^\ru (K) \cap \overline{\Vcal}$ and $\varphi^{l_0 + 1}(x) \notin \overline{\Vcal}$ by fixing $l_0 = \sup \{ l \in \N, \: \varphi^{l} (x) \in \overline{\Vcal} \}$ (which is finite by definition of $x$).

\subsubsection*{Step 2} We find for every $x$ in $W^\ru (K) \cap \Acal(0)$ an integer $k (x) > 0$ such that we have $\varphi^{k (x)}(x) \in O$. Thanks to Lemma \ref{lemma: spectral decomposition of the manifold}, we~know that there exists a basic set~$K'$ such that $x \in W^\rs (K')$. Thus $x \in W^\ru (K) \cap W^\rs (K')$. By~definition of the Smale relation, it implies $K \leq K'$. Also, we~must have $K \neq K'$. Indeed, otherwise we would have $x \in W^\ru (K) \cap W^\rs (K) = K$, according to Theorem \ref{thm Smale}, which is in contradiction with the fact that $x \in \Acal (0) \subset M\setminus K$. Consequently, we~must have $K' \subset O$ and the claimed existence of the integer $k (x)$ follows from the definition of the stable manifold of $K'$ together with the inclusion $K' \subset O$.

\subsubsection*{Step 3} We prove by contradiction that there exists $k_0 \in \N$ such that
\[
\varphi^{k_0} (W^\ru (K) \cap \Acal(0)) \subset O.
\]
By contradiction, assume that for every integer $k \in \N^*$ there exists $x_k\!\in\! W^\ru (K) \cap\nobreak \Acal (0)$ such that $\varphi^{k} (x_k) \notin O$. By~compactness of $\overline{\Vcal}$, we~can extract a subsequence $(x_{k_l})_{l \in \N}$ which converges to some $x_{\infty} \in \overline{\Vcal}$. According to Remark \ref{lemma: clos. of stable man on unrevisited neigh.}, we~must have $x_{\infty} \in W^\ru (K) \cap \overline{\Vcal}$. Also, by definition of $\Acal (0)$, the elements $\varphi^1 (x_{k_l})$ belong to $M \setminus \Vcal$. So, letting $l$ tend to $+\infty$, we~deduce that $\varphi^1 (x_{\infty}) \in M \setminus \Vcal $. Therefore, we~have $x_{\infty} \in W^\ru (K) \cap \Acal (0)$ and step 2 implies that $\varphi^{k (x_{\infty})} (x_{\infty}) \in O$. By~continuity of $\varphi^{k (x_{\infty})}$ and since $O$ is open, there exists $l_0 \geq 0$ such that $\varphi^{k (x_{\infty})} (x_{k_l}) \in O$ for every $l \geq l_0$. By~stability of $O$, this leads to
\[
\varphi^{k} (x_{k_l}) \in O, \quad \forall k \geq k(x_{\infty}), \quad \forall l \geq l_0,
\]
which is in contradiction with the construction of $x_{k_l}$ for $l$ sufficiently large.

\Changelback

\subsubsection*{Step 4} We extend step 3 to $\Acal (m)$ for $m$ large enough. Thanks to Lemma \ref{lemme: cv nonrevisite vers variete stable et instable}, we~obtain that
\[
\sup_{x \in \Acal (m) } d_g \left(x, \Acal (0) \cap W^\ru(K) \right) \underset{m \to +\infty}{\to} 0.
\]
Therefore, we~deduce by continuity from step 3 that there exists an integer $m_0 \in \N$ such that
\begin{equation}\label{eq: lemme filtration cv uniform R(n)}
\varphi^{ k_0} (\Acal (m)) \in O, \qquad \forall m \geq m_0.
\end{equation}
This concludes the proof of the sub-lemma.
\end{proof}

\begin{proof}[Proof of Lemma \ref{lemma: equivalence filtration and unrevisited}] Consider a total order relation on the family of basic sets. Let us proceed by induction to construct a filtration on $M$ stable by $\varphi^{1}$. The following arguments adapt easily to construct a filtration $\varphi^{-1}$-stable if we change $\varphi^1$ by $\varphi^{-1}$.

\subsubsection*{Base case (construction of $\Ocal_1^+$)}
Fix $\varepsilon > 0$. Since the indices of $(K_i)_{1 \leq i \leq N}$ have been chosen compatible with the relation $\leq$, the basic set $K_{N}$ must be an attractor. So if we consider some unrevisited neighborhood $\Vcal_{N}$ of $K_N$ small enough so that $\overline{\Vcal_N} \subset W^\rs(K_N)$, then $\Vcal_N$ and $\overline{\Vcal_N}$ are $\varphi^1$-stable and we will simply define
\[
\Ocal_{1}^+:= \Vcal_N \cap \varphi^{k}(\Vcal_{N})
\]
for a large enough value of $k$. Indeed, thanks to Lemma \ref{lemme: cv nonrevisite vers variete stable et instable}, we~can choose $k$ so that $\sup_{y \in \Vcal_N \cap \varphi^{k}(\Vcal_{N})} d_g (y, W^\ru (K_{N}) \cap \Vcal_N) < \varepsilon $. Since $K_N$ is an attractor, the equality $W^\ru (K_N) = K_N$ holds and it implies
\[
\sup_{y \in \Ocal_{1}^+ (\varepsilon)} d_g (y, K_N) < \varepsilon.
\]

\subsubsection*{Induction step}
Fix $\varepsilon > 0$ and assume that the open sets $\Ocal_1^+ (\varepsilon) \subseteq \Ocal_2^+ (\varepsilon) \subseteq \cdots \subseteq \Ocal_{j-1}^+ (\varepsilon) \subseteq \Ocal_j^+ (\varepsilon)$ are constructed for some $j\in\lcr1,N\rcr$ so that they define a filtration for the family of basic sets $(K_{N-i})_{0 \leq i \leq j-1}$ which is $\varepsilon$-close to the unstable manifolds, in the sense that, for every $i\in\lcr1,j\rcr$,
\[
\sup_{y \in \Ocal_i^+ (\varepsilon)} d_g \Bigl(y, \bigcup_{k \geq N- i +1} W^\ru (K_{k}) \Bigr) < \varepsilon.
\]
We want to construct a $\varphi^1$-stable $\varepsilon$-neighborhood $\Ocal_{j+1}^+ (\varepsilon)$ of $\bigcup_{k \geq N-j} W^\ru (K_{k})$. To~lighten the proof, let us denote by $K$ the basic set $K_{N-j}$. If $K$ is an attractor, then we can proceed exactly as in the base case and take the union of this open set with~$\mathcal{O}_{j}^+$. So let us assume that $K$ is not a attractor, \ie $W^\ru (K) \neq K$, and consider some unrevisited neighborhood $\Vcal$ of $K$ such that $\overline{\Vcal} \cap \Omega = K$. Applying the sub-lemma~\ref{Sublemma: uniform traveling time for unstable annulus} to $K$, $\Vcal$, $O = \Ocal_j^+ (\varepsilon)$, we~obtain integers $k_0, m_0 \geq 0$ as in the statement. Therefore, for every $m \geq m_0$, we~define
\[
O (m):= \Ocal_j^+ (\varepsilon) \cup \bigcup_{k = 0}^{k_0 - 1} \varphi^{k} (\varphi^{m} (\Vcal) \cap \Vcal),
\]
which is $\varphi^1$-stable by construction of $\Ocal_j^+ (\varepsilon)$ if we remark that
\begin{align*}
\varphi^{k_0} (\varphi^{m} (\Vcal) \cap \Vcal) &\subset \varphi^{ k_0} \left(\varphi^{m} (\Vcal) \cap \varphi^{-1} (\Vcal)\right) \sqcup \varphi^{k_0}\left(\Acal (m)\right) \\
&\subset \varphi^{ k_0 - 1} \underbrace{\left(\varphi^{m+1} (\Vcal) \cap \Vcal \right)}_{\subset \varphi^{m} (\Vcal) \cap \Vcal} \sqcup \underbrace{\varphi^{k_0}\left(\Acal (m)\right)}_{\subset \Ocal_j^+ (\varepsilon)} \quad \subset O (m).
\end{align*}
Note that $\overline{O (m)}$ is also $\varphi^1$-stable. Now, if we choose~$m$ sufficiently large according to Lemma \ref{lemme: cv nonrevisite vers variete stable et instable} so that
\[
\sup_{y \in \oldbigcup_{k = 0}^{k_0 - 1} \varphi^{k} (\varphi^{m} (\overline{\Vcal}) \cap \overline{\Vcal})}d_g \Bigl(y, W^\ru (K) \cap \bigcup_{k = 0}^{k_0 - 1} \varphi^k (\overline{\Vcal}) \Bigr) < \varepsilon,
\]
then we get the claimed result by setting $\Ocal_{j+1}^+ (\varepsilon):= O (m)$.
\end{proof}

\begin{remark}\label{rk: epsilon neighborhood proof below filtration}
Thanks to the total order relation \eqref{eq: total order relation, indices compatible}, the compact sets $\bigcup_{i \leq j} W^\rs (K_i)$ and $\bigcup_{k \geq j}W^\ru (K_k)$ intersect on $K_j$, \ie
\begin{equation}\label{eq: remark epsilon unrevisited neighborhoods}
\bigcup_{i \leq j}W^\rs (K_i) \cap \bigcup_{k \geq j}W^\ru (K_k) = K_j.
\end{equation}
Therefore, from Remark \ref{rk: filtration}, the set
\[
\Vcal_{j} (m):= \varphi^{-m}\left(\Ocal_j^- (\varepsilon) \right)\cap \varphi^m \left(\Ocal_{N-j+1}^+ (\varepsilon)\right)
\]
defines a decreasing (for the inclusion) family of unrevisited neighborhoods of $K_j$. Moreover, if we choose $m$ sufficiently large, then $\Vcal_j (m)$ can be made arbitrarily close to $K_j$ in the sense that
\[
\sup_{x \in \overline{\Vcal_{j} (m)}} d_g \left(x, K_j \right) \underset{m \to +\infty}{\to} 0.
\]
This fact can be proved by contradiction (similarly to the proof of Lemma \ref{lemme: cv nonrevisite vers variete stable et instable}) by using the relations
\[
\bigcap_{m \in \N} \overline{\Vcal_{j} (m)} = \bigcap_{m \in \N} \varphi^{-m}\bigl(\overline{\Ocal_j^- (\varepsilon)} \bigr) \cap \bigcap_{m \in \N} \varphi^m \bigl(\overline{\Ocal_{N-j+1}^+ (\varepsilon)}\bigr) = \bigcup_{i \leq j}W^\rs (K_i) \cap \bigcup_{k \geq j}W^\ru (K_k).
\]
\end{remark}

A direct application of these lemmas gives what we were looking for. Namely:

\begin{prop}\label{prop: attractor reppeler couple on the base}
For every $j\in\lcr2,N\rcr$, $(\bigcup_{k \geq j} W^\ru (K_k), \bigcup_{i < j}W^\rs (K_i))$ defines an attractor-repeller pair.
\end{prop}

\begin{proof}It is a direct application of Lemmas \ref{lemma: attractor and repeller for the flow on the basis} and \ref{lemma: equivalence filtration and unrevisited} once we have chosen $\varepsilon\ll 1$ small enough to ensure that
\[
B_g \Bigl(\bigcup_{k \geq j} W^\ru (K_k), 2\varepsilon \Bigr) \bigcap B_g \Bigl(\bigcup_{i < j} W^\rs (K_i), 2\varepsilon \Bigr) = \emptyset, \quad \forall j,\ 1 < j \leq N,
\]
where $B_g(S, 2\varepsilon)$ denotes the geodesic ball at distance $2\varepsilon$ to a compact set $S$.
\end{proof}

Now, we~are ready to construct an energy function for $\varphi^t$.

\begin{proof}[Proof of Proposition \ref{prop energy function for Axiom A flows}]
Let $\varepsilon > 0$ as in the previous proof and fix a sequence of pairwise distinct real numbers $(\lambda_i)_{1 \leq i \leq N}$ compatible with the graph structure in the sense that $\lambda_i \leq \lambda_j \Leftrightarrow K_i \leq K_j$. Up to a permutation of the indices of the basic sets, we~can assume that $\lambda_1 < \lambda_2 < \cdots < \lambda_N$.
In order to find an energy function $E$ such that $E = \lambda_j$ on $K_j$, we~will apply Lemma \ref{lemma: energy function for attractor-repeller model} for each attractor-repeller given in Proposition \ref{prop: attractor reppeler couple on the base}. Thanks to Lemma \ref{lemma: equivalence filtration and unrevisited} together with Remark \ref{rk: epsilon neighborhood proof below filtration}, there exist filtrations $\Ocal_0^- (\varepsilon) \subset \Ocal_1^- (\varepsilon) \subset \cdots \subset \Ocal_{N}^- (\varepsilon)$ and $\Ocal_N^+ (\varepsilon) \supset \cdots \supset \Ocal_{1}^+ (\varepsilon) \supset \Ocal_0^+ (\varepsilon)$ which are respectively $\varphi^{-1}$-stable and $\varphi^{1}$-stable such that
\[
\forall j,\ 1 \leq j \leq N, \qquad \Ocal_{N-j+1}^+ (\varepsilon) \cap \Ocal_{j}^- (\varepsilon) \text{ is a } \varepsilon \text{-neighborhood of } K_j.
\]
Thanks to Lemma \ref{lemma: energy function for attractor-repeller model}, we~obtain for every $j\in\lcr2,N\rcr$ a smooth energy function $E_{j} \in \Ccal^{\infty} (M; [0,1])$, $\varepsilon$-neighborhoods $\Wcal_{j}^- \subset \Ocal_{j-1}^- (\varepsilon)$ and $\Wcal_{j}^+ \subset \Ocal_{N-j+1}^+ (\varepsilon)$ of $ \bigcup_{i < j} W^\rs (K_i)$ and $ \bigcup_{k \geq j} W^\ru (K_k)$ respectively, a constant $\eta_0 > 0$ (which only depends on $\varepsilon$) such that $\Lie_V (E_j) \geq 0$ on $M$ and
\begin{itemize}
\item $\Lie_V (E_j) > \eta_0$ on $M\setminus (\Wcal_j^- \cup \Wcal_j^+)$;
\item $E_j < \varepsilon$ on $\Wcal_j^-$ and $E_j = 0$ on $\bigcup_{i < j} W^\rs (K_i)$; in~particular, we~have $E_j = 0$ on $\bigcup_{i < j} K_i$;
\item $E_j > 1 - \varepsilon$ on $\Wcal_j^+$ and $E_j = 1$ on $\bigcup_{k \geq j} W^\ru (K_k)$; in~particular, we~have $E_j = 1$ on $\bigcup_{k \geq j} K_k$.
\end{itemize}
We define a global energy function $E \in \Ccal^{\infty} (M)$ as a linear combination of previous energy functions:
\[
E = \lambda_1 + \sum_{j = 2}^N (\lambda_{j} - \lambda_{j-1}) E_{j}.
\]
Thanks to the properties of the energy functions stated above, we~deduce that
\[
\Lie_V (E) = \sum_{j = 1}^N (\lambda_j - \lambda_{j-1}) \Lie_V (E_j) > \min_{1 < j \leq N}(\lambda_j - \lambda_{j-1}) \eta_0 =: \eta \: \text{ on } \: \biggl(\bigcap_{j =2}^N (\Wcal_j^- \cup \Wcal_j^+) \biggr)^c.
\]
It remains to prove that $\bigcap_{j = 2}^N (\Wcal_j^- \cup \Wcal_j^+)$ is an $\varepsilon$-neighborhood of the non-wandering set and that $E$ is close to $\lambda_i$ near $K_i$ with equality on $K_i$. One has:
\begin{align*}
\bigcap_{j = 2}^N (\Wcal_j^- \cup \Wcal_j^+) &= \bigcup_{\tau: \lcr1,N\rcr \to \{\pm\}} \bigcap_{j = 2}^N \Wcal_j^{\tau (j)} \\
&= \bigcup_{j = 1}^N \Ncal_j:= \bigcup_{j = 1}^N \left(\Wcal_{2}^+ \cap \cdots \cap \Wcal_{j}^+ \cap \Wcal_{j+1}^- \cap \cdots \cap \Wcal_{N}^- \right),
\end{align*}
using the convention $\Ncal_1 = \Wcal_{2}^- \cap \cdots \cap \Wcal_{N}^-$.
Note that the second equality is a consequence of the fact that the intersection $\Ocal_{j-1}^- (\varepsilon) \cap \Ocal_{N-j+1}^+ (\varepsilon)$ is empty for all $j\in\lcr2,N\rcr$ for $\varepsilon $ small enough, which implies:
\[
\text{if } 2 \leq i < j \text{ and } (\tau (i), \tau (j)) = (-,+) \text{ then } \bigcap_{j = 2}^N \Wcal_j^{\tau (j)} = \emptyset.
\]
\Changel
\noindent
Furthermore, thanks to the definition of the sets $W_l^{\pm}$ for $l\in\lcr2,N\rcr$ and since $(\Ocal_l^{\pm})_l$ is a $\varepsilon$-filtration for $\varphi^{\pm 1}$, the sets $\Ncal_j$ satisfy the inclusion
\[
\Ncal_j \subset \Ocal_{N-1}^+ (\varepsilon) \cap \cdots \cap \Ocal_{N-j+1}^+ (\varepsilon) \cap \Ocal_{j}^- (\varepsilon) \cap \cdots \cap \Ocal_{N-1}^- (\varepsilon) = \Ocal_{N-j+1}^+ (\varepsilon) \cap \Ocal_{j}^- (\varepsilon)
\]
and are therefore $\varepsilon$-neighborhoods of the $K_j$ according to remark \ref{rk: epsilon neighborhood proof below filtration}.
Moreover, we~have on each $K_i$:
\[
E = \lambda_1 + \sum_{j = 2}^N (\lambda_j - \lambda_{j-1}) E_j = \lambda_1 + \sum_{j = 2}^N (\lambda_j - \lambda_{j-1}) \delta_{j \leq i} = \lambda_i.
\]
\Changelback
\noindent
It remains to prove that $E$ is close to $\lambda_i$ on the neighborhood $\Ncal_i$ of $K_i$. To that aim, let us note that for every $x \in M$ and for every $i\in\lcr1,N\rcr$, we~have
\[
|E(x) - \lambda_i | \leq \sum_{j = 2}^N (\lambda_j - \lambda_{j-1}) |E_j (x) - \delta_{j \leq i} |.
\]
For any $j\in\lcr2,N\rcr$ and $x \in \Ncal_i$ we will bound $|E_j (x) - \delta_{j \leq i}|$ by a small quantity independent of $x$ in $\Ncal_i$. We have two cases to deal with:
\begin{itemize}
\item If $i < j$, then $\delta_{j \leq i} = 0$ and $|E_j (x) - 0| = E_j (x) < \varepsilon$ by definition of $E_j$.
\item If $i \geq j$, then $\delta_{j \leq i} = 1$ and $|E_j (x) - 1| = 1 - E_j (x) < \varepsilon$ again by definition of $E_j$.
\end{itemize}
Finally, we~obtain the upper bound
\[
\sup_{1 \leq i \leq N} \sup_{x \in \Ncal_i} |E (x) - \lambda_i| \leq \varepsilon \sum_{j = 2}^N (\lambda_j - \lambda_{j-1}) \leq (\lambda_N - \lambda_1) \varepsilon.
\]
Up to dividing $\varepsilon$ by $\max (1, \lambda_N- \lambda_1)$ everywhere, we~get the claimed result.
\end{proof}

Now, the idea will be to perform the same analysis for the (projected Hamiltonian) flow $\widetilde{\Phi}^t$ acting on $S^* M$. However, in that case, proving that one has an attractor-repeller structure for $(\Sigma_{\rs}, \Sigma_{\ruo})$ and $(\Sigma_{\rso}, \Sigma_{\ru})$ reveals to be more challenging because we need to understand what happens to the fiber part of $\widetilde{\Phi}^t (x,\xi)$ when the orbit of $(x,\xi)$ comes close to a basic set. The next part is devoted to the local analysis of the Hamiltonian near basic sets in view of applications to the proofs of Proposition \ref{prop energy function for Axiom A flows}, Lemmas \ref{lemma: attractor and repeller for the flow on the basis} and \ref{lemma: equivalence filtration and unrevisited} and finally the existence of the energy function on $S^*M$.

\section{Compactness result and energy functions for the Hamiltonian flow}\label{section: Compactness result and energy functions for the Hamiltonian flow}

On a basic set $K$, one can define conical neighborhoods of the unstable distributions $E_\ru^*$ and $E_{\ruo}^*$ which are stable under the Hamiltonian flow $\Phi^t$ as soon as $t > 0$. This characterization of hyperbolicity is sometimes referred to as the Alekseev cone field criterion \cite[Prop.\,5.1.7]{FH}. The goal of this part is to extend cones of the weak/stable/unstable distributions in an invariant fashion on a small neighborhood of the non-wandering set according to the global dynamics given by the Smale ordering of the basic sets.

\subsection{Adapted metric on a basic set}\label{subsubsection adapted metric}

From the fixed Riemannian metric $g$ on $M$, we~can define on any basic set $K$ not reduced to a fixed point the following new metric called \emph{adapted metric} for the flow $\varphi^t$. More precisely, for every $x \in K$ and every $v = v_s + v_u + v_o \in E_\rs (x) \oplus E_\ru (x) \oplus E_\ro (x) = T_x M$, we~set
\begin{align}\label{eq: def extended metric near basic set}
\quad\hat{g} (v,v) \!\!&\hphantom{:}= \hat{g}(v_s, v_s) + \hat{g}(v_u, v_u) + \hat{g}(v_o, v_o) \notag\\[-1pt]\\[-10pt]\notag
&:= \int_0^{+\infty} e^{\lambda t /2} (\varphi^{t*} g) (v_s,v_s) dt + \int_0^{+\infty} e^{\lambda t /2} (\varphi^{-t*} g) (v_u,v_u) dt + \alpha(x) (v_o)^2,
\end{align}
where $\lambda$ denotes the hyperbolic exponent on the basic set $K$ (see Appendix \ref{Appendix: hyperbolic set}) and~$\alpha$ denotes the Anosov $1$-form defined on $K$ by
\[
\Ker \alpha (x) = E_\ru (x) \oplus E_\rs (x), \qquad \alpha (x) (V(x)) = 1.
\]
In the case where the basic set is reduced to a fixed point, we~call adapted metric for the flow a metric which satisfies (\ref{eq: def extended metric near basic set}) without the term $\alpha (x) (v_0)^2$.

This new metric is well-defined thanks to hyperbolicity and to the invariance properties of the vector bundles on $K$. Recall also that the distributions $E_\rs, E_\ru$ are only Hölder continuous in general. Precisely, $\hat{g}$ will be seen as a continuous section of the vector bundle of metrics (\ie symmetric $(2,0)$-tensors) $\otimes_{\sym}^2 T M \to M$ defined on the compact set $K$. Let us denote by $|.|_{\hat{g}}$ the norm induced by $\hat{g}$, \ie $|v|_{\hat{g}}:= \sqrt{\hat{g}(v,v)}$. The metric is said to be \textit{adapted} due to the following hyperbolic estimates: for every $x \in K$,
\begin{equation}\label{hyperbolic estimates for the adapted metric}
\begin{split}
|D\varphi^t (x) v_s |_{\hat{g}} &\leq e^{-\lambda t /2} |v_s |_{\hat{g}}, \qquad \forall t \geq 0, \: \forall v_s \in E_\rs (x) \\
|D\varphi^{-t} (x) v_u |_{\hat{g}} &\leq e^{-\lambda t /2} |v_u |_{\hat{g}}, \qquad \forall t \geq 0, \: \forall v_u \in E_\ru (x) \\
|D\varphi^t (x) v_o |_{\hat{g}} &= |v_o |_{\hat{g}}, \qquad \forall t \in \R, \:\forall v_o \in E_\ro (x).
\end{split}
\end{equation}

\subsection{Extension of the invariant distributions near basic sets.}

In order to analyze the dynamics near basic sets, it will be convenient to extend the previous hyperbolic estimates in some neighborhood of $K$. To that aim, we~need to extend the distributions $E_\rs, E_\ru$, $E_\ro$ and $|.|_{\hat{g}}$ near each basic set. This can be achieved thanks to the following lemma:

\begin{lemma}[{Extension lemma, \cite[Lem.\,4.4 p.\,128]{HPPS}}]\label{Lemme geometrique prolongement des sections} Let $X$ be a smooth manifold and let $\pi: \Ecal \to X$ be some fiber bundle over $X$. If $\Gamma: K \to \Ecal$ denotes a continuous section defined on a compact set $K \subseteq M$, then $\Gamma$ extends as a continuous section $\overline{\Gamma}: \Ncal \to E$ on a neighborhood $\Ncal$ of $K$.
\end{lemma}

\subsubsection{Extension of the adapted metric near a basic set} Let us apply this lemma with
\[
X:= M \quad \text{and} \quad \Ecal = T^* M \otimes_{\sym} T^* M.
\]
The Riemannian metric $\hat{g}$ can be seen as a continuous section $\hat{g}: K \to \Ecal$. Since the basic set $K$ is a compact subset of $M$, the lemma applies and it enables to extend~$\hat{g}$ continuously on an open neighborhood $\Ncal_0$ of $K$. Up to considering smaller $\Ncal_0$, we~can assume that the extended metric remains Riemannian on $\Ncal_0$ as positivity and semi-definiteness are open conditions. Since we can do this extension near each basic set and since the metric $g$ is Riemannian, a partition of unity argument enables to prove:

\begin{lemma}
For any Axiom~A flow on a compact manifold, there exists a continuous Riemannian metric (globally defined) which is adapted to the dynamics on each basic set, in the sense that \eqref{hyperbolic estimates for the adapted metric} holds on each basic set.
\end{lemma}

In what follows, we~will always assume that $g$ is a continuous Riemannian metric adapted to the dynamics on each basic set and we will denote by $|.|$ its norm on the fibers of $TM$ and $T^* M$ to lighten notations.

\subsubsection{Extension of distributions} We now apply Lemma \ref{Lemme geometrique prolongement des sections} in order to extend the distributions $E_{\rs}^*$, $E_{\rso}^*$, $E_{\ru}^*$, $E_{\ruo}^*$ (defined on $K$) continuously on a neighborhood of~$K$. Note that we already explained that $E_\rs^*, E_\ru^*,\dots$ are well-defined all over $M$ using the partition into unstable manifolds. The point of this new extension based on the bundles on $K$ (and not on the global dynamics) is that we expect that these new bundles have good hyperbolic properties in the sense of (\ref{hyperbolic estimates for the adapted metric}). We will also need to make sure that our local analysis is related to the invariant distributions $E_{\ru/\ruo}^*$ and $E_{\rs/\rso}^*$ defined all over $M$.

Fix a basic set $K$ and recall that the dimension of the distributions $E_\rs, E_\ru, E_\ro$ are constant on $K$. We denote by $d_{\rs} = \dim E_\rs$ and $d_{\ru} = \dim E_\ru$ their dimension on $K$. According to Lemma \ref{lemma: closure of local stable manifold on a basic set}, the inclusions $\overline{W_{\varepsilon}^\rs (K)} \subseteq W_{\varepsilon_0}^\rs (K)$ and $\overline{W_{\varepsilon}^\ru (K)} \subseteq W_{\varepsilon_0}^\ru (K)$ are satisfied as soon as $\varepsilon < \varepsilon_0$ where $\varepsilon_0 > 0$ is given by the stable manifold theorem.

\begin{lemma}\label{lemma: continuity of extension of distributions}For every $0 < \varepsilon < \varepsilon_0$, with $\varepsilon_0 > 0$ given by the stable manifold theorem, the map
\[
\Gamma: \overline{W_{\varepsilon}^{\rs} (K)} \to G_{\rs}(M), \quad \Gamma(x) = E_\rs (x)
\]
with value in the Grassmann vector bundle of subspaces of dimension $d_{\rs}$, is continuous.
\end{lemma}

An important remark is that we aim to extend the stable distribution $E_\rs$, defined on the stable set of $K$, on a neighborhood of $K$ in a continuous fashion. This extension will be denoted by $\widetilde{E}_\rs$ to avoid confusion. The subtle point is the following fact: $E_\rs$ is not continuous in general on a neighborhood of $K$ but only on the stable set $W_{\varepsilon_0}^\rs (K)$ of $K$. Here, $\widetilde{E}_\rs$ will be equal to $E_\rs$ on the stable set of $K$ and will be continuous on a neighborhood of $K$ by construction.

\Changel
\begin{proof}
The result is a direct consequence of the stable manifold theorem when the basic set is reduced to a fixed point. So let us assume that the basic set $K$ is not reduced to a fixed point. In~this setting, the proof will be a consequence of the $\lambda$-lemma (also called inclination lemma) in a slightly more general version than usual which is recalled in Proposition \ref{Prop: lambda lemma for non periodic orbit} from Appendix \ref{Appendix: proof of Smale ordering}. Indeed, we~first consider a sequence $(z_n) \in K^{\N}$, a family of elements $x_n \in W_{\varepsilon}^\rs (z_n)$, $v_n \in E_\rs (x_n) = T_{x_n} W_{\varepsilon}^\rs (z_n)$ and we assume that the three sequences $(z_n),(x_n),(v_n)$ are converging respectively to $z \in K$, $x \in M$, $v \in T_x M$. Proceeding as in the proof of Lemma \ref{lemma: closure of local stable manifold on a basic set}, we~deduce that $x \in W_{\varepsilon_0}^\rs (z)$. Now, we~assume by contradiction that $v \notin E_\rs (x) = T_x W_{\varepsilon_0}^\rs (z)$ and we will use the $\lambda$-lemma to get the contradiction. Precisely, we~are going to prove first that $v \in E_{\rso} (x) = T_x W_{\varepsilon_0}^{\rso} (z)$ using the $\lambda$-lemma and then we will deduce that $v \in E_\rs (x)$. First, we~suppose that $v \notin E_{\rso}(x)$ and we consider a $\Ccal^{\infty}$ disk~$\D$ which intersects $W_{\varepsilon_0}^{\rso} (z)$ transversally, so that it satisfies the hypothesis of the $\lambda$\nobreakdash-lemma, such that \hbox{$\D \cap W_{\varepsilon_0}^{\rso} (z) = \{x\}$} and such that $v \in T_x \D$. The $\lambda$-lemma implies that $\varphi^k \left(\D \right)$ is arbitrarily close to the weak-unstable manifold $W_{\varepsilon_1}^{\ru} (\varphi^k (z))$ for the $\Ccal^1$ topology when $k \to +\infty$ up to considering a smaller disk $\D$ (of same dimension) containing~$x$, where~$\varepsilon_1$ is the constant given in Proposition \ref{Prop: lambda lemma for non periodic orbit}. In~particular, the vector $D \varphi^k (x)v \in D \varphi^k (x) \left(T_x \D \right) = T_{\varphi^k (x)} \varphi^k \left(\D \right)$ can be made arbitrarily close to $E_{\ru} (\varphi^k (z))$ as $k \to + \infty$ in the sense that $D \varphi^k (x)v$ belongs to any conical neighborhood $C_{\ru} (\varphi^k (x))$ defined as follows. For all $y \in U$ sufficiently small which is given by Lemma~\ref{Lemme geometrique prolongement des sections}, we~denote by
\[
C_{\ru} (y) = \left\{w = w_{\rso} + w_{\ru} \in T_y M =\widetilde{E}_{\rso} (y) \oplus \widetilde{E}_{\ru} (y), |w_{\ru}| > |w_{\rso}| \right\}
\]
the conical neighborhood of the extended weak-unstable distribution. Here, $T_y M = \widetilde{E}_{\rso} \oplus \widetilde{E}_{\ru}$ denotes a continuous extension of the hyperbolic splitting of the tangent space in the neighborhood $U$ of $K$ given by Lemma \ref{Lemme geometrique prolongement des sections}, which follows from the continuity of the stable/weak-unstable distributions on $K$ \cite[Prop.\,5.1.4]{FH}. Moreover, thanks to the stable manifold theorem, we~know that the stable manifold $W_{\varepsilon_0}^\rs (z')$ at each point \hbox{$z' \in K$} satisfies by continuity the inclusion $ T_y W_{\varepsilon_0}^{\rs} (z') \subset C_{\rs} (y)$ (which is defined similarly to the weak-unstable cones exchanging $u$ and $so$) as long as $y$ is close enough to $K$. Now, we~fix $k$ sufficiently large so that $D \varphi^k (x)v \in C_{\ru} (\varphi^k (x))$. Since, on one hand, $D \varphi^k (x_l)v_l$ converges to $D \varphi^k (x)v$, and since we proved that it belongs to $ T_{\varphi^k (x_l)} W_{\varepsilon_0}^{\rs} (\varphi^k (z_l)) \subset C_{\rs} (\varphi^k (x_l))$ on the other hand, we~must have by continuity $ D \varphi^k (x)v \in \overline{C_{\rs} (\varphi^k (x))}.$
Finally, putting all together, we~obtain $D \varphi^k (x)v \in C_{\ru} (\varphi^k (x))$ from the $\lambda$-lemma and we obtain $ D \varphi^k (x)v \in \overline{C_{\rs} (\varphi^k (x))}$ by our continuity argument. Therefore, since $\overline{C_{\rs} (\varphi^k (x))} \cap C_{\ru} (\varphi^k (x)) = \emptyset$, we~get that $v \in E_{\rso} (x)$. It remains to prove that $v \in E_\rs (x)$. To do so, we~write $v = v_\rs + v_\ro \in E_\rs (x) \oplus \R V(x)$ ($v_\ro \neq 0$) and we see that
\[
|D\varphi^t (x) v| = |D\varphi^t (x) v_\rs + D\varphi^t (x) v_\ro| \underset{t \to +\infty}{\sim} |D\varphi^t (x) v_\ro | = |v_\ro|,
\]
where the last equality follows from the choice of an adapted metric on $K$. Thus, there exists $t_0 \in \R$ such that
\[
\forall t \geq t_0, \quad D\varphi^t (x) v \in C_{\ruo}(\varphi^t (x)).
\]
Moreover, since, we~also have that $ D \varphi^k (x)v \in \overline{C_{\rs} (\varphi^k (x))}$ for any sufficiently large integer $k$, we~get that $D \varphi^k (x)v \in \overline{C_{\rs} (\varphi^k (x))} \cap C_{\ruo}(\varphi^k (x)) = \emptyset$ for such a $k$. This leads to the expected contradiction and it implies that $v \in E_\rs (x)$.
\end{proof}
\Changelback%

In the next paragraph, we~aim to extend the hyperbolic splitting of the tangent space on a neighborhood of the local stable manifold $W_{\varepsilon}^\rs (K)$ using the previous lemma together with Lemma \ref{Lemme geometrique prolongement des sections}. It will give us a hyperbolic splitting with some invariance properties in the stable manifold and the unstable manifold of $K$ which will be essential in the proof of Proposition \ref{prop compactness}.

If we apply Lemma \ref{Lemme geometrique prolongement des sections} to the continuous section defined in the previous lemma, then we obtain an open set $\Ncal_{\rs} \supset \overline{W_{\varepsilon}^\rs (K)} $ and an extension $\widetilde{E}_\rs$ of the distribution $\bigcup_{x \in \overline{W_{\varepsilon}^\rs (K)}}E_{\rs} (x)$ on $\Ncal_{\rs}$. If we replace $\rs$ by $\ru$ in the previous construction then we obtain similarly a continuous extension $\widetilde{E}_\ru$ of $ \bigcup_{x \in \overline{W_{\varepsilon}^\ru (K)}} E_{\ru} (x)$ on a neighborhood~$\Ncal_\ru$ of $\overline{W_{\varepsilon}^\ru (K)}$.
Next, we~define $ \widetilde{E}_{\rso}^*$ (\resp $\widetilde{E}_{\ruo}^*$) by taking the dual orthogonal of $\widetilde{E}_\rs$ (\resp $\widetilde{E}_\ru$). The notation $\widetilde{E}_{\rso}^*$ can seem a little ambiguous at first, because we extend first and then take the dual orthogonal. Yet, everything is consistent here since $\widetilde{E}_{\rso}^*$ also extends continuously the distribution $\bigcup_{x \in \overline{W_{\varepsilon}^\rs (K)}} E_{\rso}^* (x)$. A similar remark holds for $\widetilde{E}_{\ruo}^*$. Moreover, by setting $\widetilde{E}_\rs^*:= \widetilde{E}_{\rso}^* \cap \{\xi \in T_x^*M, \: \xi(V(x)) = 0\}$ we obtain a continuous extension of $\bigcup_{x \in \overline{W_{\varepsilon}^\rs (K)}}E_{\rs}^* (x)$. The different steps can be summarized in the next diagram (which of course also hold if we replace s by u):
\[
\bigcup_{x \in \overline{W_{\varepsilon}^\rs (K)}} E_\rs (x) \dra \widetilde{E}_\rs \dra \widetilde{E}_{\rso}^* \dra \widetilde{E}_\rs^*.
\]
Moreover, the distributions $\widetilde{E}_{\rso/\rs}^* $ extend $E_{\rso/\rs}^*$ on a neighborhood of the local stable manifold of $K$:
\[
\forall x \in W_{\varepsilon}^\rs (K), \quad \widetilde{E}_{\rso}^* (x) = E_{\rso}^* (x) \;\text{ and } \widetilde{E}_\rs^* (x) = E_\rs^* (x).
\]
Similarly, we~have
\[
\forall x \in W_{\varepsilon}^\ru (K), \quad \widetilde{E}_{\ruo}^* (x) = E_{\ruo}^* (x) \;\text{ and } \widetilde{E}_\ru^* (x) = E_\ru^* (x).
\]
These two last statement will be crucial in the proof of the compactness proposition~\ref{prop compactness}.

Finally, in order to extend continuously the neutral direction, we~define for all $x \in \Ncal =\Ncal_{\rs} \cap \Ncal_\ru$,
\[
\widetilde{E}_\ro^* (x):= \bigl\{ \xi \in T_x^* M, \: \xi \bigl(\widetilde{E}_\rs(x) + \widetilde{E}_\ru(x)\bigr) = 0 \bigr\}.
\]

\begin{remark}\label{remark: upper triangular block matrices}
It is important to note that for every $x \in W_{\varepsilon}^{\rs} (K) \cap \Ncal$, we~have $\widetilde{E}_\ro^* (x) \subseteq \widetilde{E}_{\rso}^* (x)$. Similarly, we~have for every $x \in W_{\varepsilon}^{\ru} (K) \cap \Ncal$, the inclusion $\widetilde{E}_\ro^* (x) \subseteq \widetilde{E}_{\ruo}^* (x)$.
\end{remark}

\begin{figure}[t]
\centering
\def\svgscale{0.5}
\includesvg{neighborhoods_extension_distribution}
\caption{Illustration of some neighborhoods used in the extension of distributions}
\label{fig: neighborhoods extension of distributions}
\end{figure}

Up to considering smaller $\Ncal$, we~can assume that
\begin{equation}\label{decomposition cotangent au voisinage d un compact basique hyperbolique}
\widetilde{E}_{\ru}^*(x) \oplus \widetilde{E}_{\rs}^*(x) \oplus \widetilde{E}_\ro^*(x) = T_x^* M, \qquad \forall x \in \Ncal.
\end{equation}
Note that this decomposition of the cotangent space is \emph{not invariant} by the flow~$\Phi^t$ in general. Despite that, hyperbolic estimates as well as stability of good conical neighborhoods of these bundles should extend by continuity on a neighborhood of $K$.

\subsubsection{Stability of conical neighborhoods near a basic set}

We set for all $\delta > 0$ and all $x \in \Ncal$,
\begin{equation}\label{definition voisinage conique tilde}
\begin{split}
\Ccal_{\ru}^{\delta} (x) &:= \bigl\{ \xi \in T_{x}^* M, \: \delta |\xi_\ru| > |\xi_\rs| + |\xi_\ro| \bigr\}, \\
\Ccal_{\ruo}^{\delta} (x) &:= \bigl\{ x \in T_{x}^* M, \: \delta (|\xi_\ru| + |\xi_\ro|) > |\xi_\rs| \bigr\},
\end{split}
\end{equation}
where we used the decomposition (\ref{decomposition cotangent au voisinage d un compact basique hyperbolique}) on the fibers, \ie $\xi = \xi_{\ru} + \xi_{\rs}+ \xi_{\ro} \in \widetilde{E}_{\ru}^*(x) \oplus \widetilde{E}_{\rs}^*(x) \oplus \widetilde{E}_\ro^*(x)$. If we replace $\rs$ by $\ru$ in (\ref{definition voisinage conique tilde}), then we can define similarly the stable and weak-stable conical neighborhoods $\Ccal_{\rs}^{\delta}$ and $\Ccal_{\rso}^{\delta}$ on $\Ncal$. The main technical statement of this section is:

\begin{lemma}\label{stabilite voisinage conique tilde pour Phi1} Let $K$ be a basic set and let $\Ncal$ be the open neighborhood appearing in~\eqref{decomposition cotangent au voisinage d un compact basique hyperbolique}. There exists $\varepsilon_1 >0$ such that, for every unrevisited neighborhood $\Vcal \subset \Ncal \cap\nobreak\varphi^{-1}(\Ncal)$ contained in an $\varepsilon_1$-neighborhood of $K$, the following hold:
\begin{enumeratei}
\item\label{stabilitei} For every $\delta_0\in(0,1]$, one can find $m_{\delta_0} \geq 0$ so that, for every $\delta\in[0,\delta_0]$, for every $m \geq m_{\delta_0}$ and for every $x\in\mathcal{V}\cap\varphi^{m}(\Vcal)$,
one has the inclusions
\begin{equation}\label{lemma stab vois conique tilde - Unstable inlusions}
\Phi^{1}\left(\Ccal_{\ru}^{\delta}(x)\right) \subseteq \Ccal_{\ru}^{\delta'}(\varphi^{1}(x)), \qquad
\Phi^{1}\left(\Ccal_{\ruo}^{\delta}(x) \right) \subseteq \Ccal_{\ruo}^{\delta'}(\varphi^{1}(x))
\end{equation}
where $\delta' = e^{-\lambda/3} \delta$ and with $\lambda$ being the constant appearing in the definition \eqref{eq: def extended metric near basic set}.
\item\label{stabiliteii} Moreover, there exists $\delta_1\in(0,1]$ such that for every $\delta \leq \delta_1$, one can find $m_{\delta} \geq 0$ so that, for every $m \geq m_\delta$, for every $x \in \Vcal \cap\varphi^{m}(\Vcal)$ and for every $\xi \in C^\delta_\ru(x)$, we~have
\[
|\Phi^1 (x,\xi)| \geq e^{\lambda /4} |\xi|.
\]
\end{enumeratei}
\end{lemma}

\begin{proof}
The first step of the proof consists in extending the hyperbolic estimates on the unstable manifold of $K$. Let us denote by $\pi_{\alpha}$ the continuous projector on $\widetilde{E}_{\alpha}^*$ for each $\alpha \in \{\rs,\ro,\ru\}$ and let us define
\[
\forall \alpha, \beta \in \{\rs,\ro,\ru\}, \quad A_{\alpha \beta}:= \pi_{\alpha} \circ \Phi^1 \circ \pi_{\beta} \text{ and } B_{\alpha \beta}:= \pi_{\alpha} \circ \Phi^{-1} \circ \pi_{\beta}.
\]
We will also use the notation $\xi_{\alpha}^1$ for $\pi_{\alpha} \circ \Phi^1 (\xi)$. For every $x \in \Ncal \cap \varphi^{-1}(\Ncal) $, we~have $\xi = \xi_\ru + \xi_\rs + \xi_\ro \in \widetilde{E}_\ru^* (x) \oplus \widetilde{E}_\rs^* (x) \oplus \widetilde{E}_\ro^* (x) = T_x^* M$ and $\Phi_x^1 (\xi) = \xi_\ru^1 + \xi_\rs^1 + \xi_\ro^1 \in \widetilde{E}_\ru^* (\varphi^1 (x)) \oplus \widetilde{E}_\rs^* (\varphi^1 (x)) \oplus \widetilde{E}_\ro^* (\varphi^1 (x))$. The dynamics of $\Phi^1$ and $\Phi^{-1}$ can be encode within the two following matrices of linear morphisms:
\[
\begin{pmatrix}\xi_\ru^1 \\ \xi_\ro^1 \\ \xi_\rs^1 \end{pmatrix}
= \begin{pmatrix}
A_{\ru\ru} & A_{\ro\ru} & A_{\rs\ru} \\
A_{\ruo} & A_{\ro\ro} & A_{\rso} \\
A_{\ru\rs} & A_{\ro\rs} & A_{\rs\rs}
\end{pmatrix}
\begin{pmatrix}\xi_\ru \\ \xi_\ro \\ \xi_\rs \end{pmatrix}
\qquad \text{and} \qquad \begin{pmatrix}\xi_\ru \\ \xi_\ro \\ \xi_\rs \end{pmatrix}
= \begin{pmatrix}
B_{\ru\ru} & B_{\ro\ru} & B_{\rs\ru} \\
B_{\ruo} & B_{\ro\ro} & B_{\rso} \\
B_{\ru\rs} & B_{\ro\rs} & B_{\rs\rs}
\end{pmatrix}
\begin{pmatrix}\xi_\ru^1 \\ \xi_\ro^1 \\ \xi_\rs^1 \end{pmatrix}.
\]
To be more precise, we~should have written $A_{\alpha \beta} (x)$ and $B_{\alpha \beta} (\varphi^1(x))$ to indicate that $\xi \in T_x^* M$. In~the particular case where $x \in W^\ru (K) \cap \Ncal \cap \varphi^{-1}(\Ncal)$, both matrices are upper triangular block matrices thanks to our definition of $\widetilde{E}_\ru^*$, $\widetilde{E}_\ro^*$ and $\widetilde{E}_\rs^*$, \ie
\[
A_{\ruo} = A_{\ru\rs} = A_{\ro\rs} = B_{\ruo} = B_{\ru\rs} = B_{\ro\rs} = 0 \quad\text{on }W^\ru (K) \cap \Ncal \cap \varphi^{-1}(\Ncal).
\]
However, on the whole open set $\Ncal \cap \varphi^{-1}(\Ncal)$ this is not the case anymore.\footnote{This follows from the fact that in the general case (where $\widetilde{E}_\rs^*$ and $\widetilde{E}_\ru$ are both non trivial) none of the extended distribution $\widetilde{E}_\rs^*$, $\widetilde{E}_\ru^*$ or $\widetilde{E}_\ro^*$ are invariant by $\Phi^t$ on $\Ncal$.} Now, we~extend the hyperbolic estimates (\ref{hyperbolic estimates for the adapted metric}) on the unstable manifold of $K$. For every $\varepsilon \geq 0$, we~define the set $\Ncal(\varepsilon)$ of points $y \in \Ncal \cap \varphi^{-1}(\Ncal)$ on which we have for all $\xi \in T_{y}^* M$,
\begin{gather}
|A_{\ru\ru} (\xi_\ru)| \geq (e^{\lambda/2 }-\varepsilon) |\xi_\ru|, \quad |B_{\rs\rs} (\xi_\rs^1)| \: \geq (e^{\lambda/2 } - \varepsilon) |\xi_\rs^1|,
\notag\\
\label{eq perturbe: lemme stabilite vois conique tilde - estimées matrices}
|A_{\ro\ro} (\xi_\ro)| \geq (1 - \varepsilon) |\xi_\ro|, \quad |B_{\ro\ro} (\xi_\ro^1)| \geq (1 - \varepsilon) |\xi_\ro^1|
\\
\notag
\| A_{\ro\ru} \| \leq \varepsilon, \quad \| A_{\rs\ru} \| \leq \varepsilon, \quad \| A_{\rso} \| \leq \varepsilon,\\
\notag
\| B_{\ro\ru} \| \leq \varepsilon, \quad \| B_{\rs\ru} \| \leq \varepsilon, \quad \| B_{\rso} \| \leq \varepsilon, \\
\notag
A_{\ruo} = A_{\ru\rs} = A_{\ro\rs} = B_{\ruo} = B_{\ru\rs} = B_{\ro\rs} = 0,
\end{gather}
where $\|. \|$ denotes the operator norm defined by:
\[
\forall \alpha, \beta \in \{\ru,\ro,\rs\}, \quad \| A_{\alpha \beta} \|:= \sup_{\xi_{\alpha} \in \widetilde{E}_{\alpha}^* (y) \setminus \{0 \}} \frac{|A_{\alpha \beta} (\xi_{\alpha})|}{ |\xi_{\alpha}|}.
\]
Note that the map $A_{\alpha \beta}$ and the norm $\|.\|$ depend on the point $y$, but we will not precise the point $y$ if everything is clear. For every $\varepsilon > 0$ and for every unrevisited neighborhood $\Vcal$ sufficiently close to $K$, we~have $\Vcal \cap W^\ru (K) \subset \Ncal (\varepsilon)$. Let first check that the inclusions (\ref{lemma stab vois conique tilde - Unstable inlusions}) hold on the unstable manifold of $K$ for some $\delta' \leq \delta e^{-\lambda/2}$ as soon as $\varepsilon$ is sufficiently small. The result on the whole neighborhood will then follow by continuity.

\subsubsection*{First inclusion, on unstable cones $\Ccal_\ru^{\delta}$} Let us fix $\delta\in(0,1]$, $x \in W^\ru (K) \cap \Ncal \cap \varphi^{-1} (\Ncal)$ and $\xi \in \Ccal_\ru^{\delta} (x)$. Our goal is to find a parameter $\delta' \leq \delta e^{-\lambda/2} $ such that $\xi^1:= \Phi_x^1 (\xi) \in \widetilde {\Ccal}_\ru^{\delta'} (\varphi^1 (x))$. The first step consists in computing $|\xi_\ru^1|^2$ in order to find a lower bound which depends on $|\xi_\rs|$ and $|\xi_\ro|$. Precisely, we~have
\begin{equation}\label{eq: proof perturb vois conique tilde - lower bound unstable}
\begin{split}
\delta |\xi_\ru^1| &= \delta |A_{\ru\ru} \xi_\ru + A_{\rs\ru} \xi_\rs + A_{\ro\ru} \xi_\ro | \geq \delta |A_{\ru\ru} \xi_\ru| - \delta |A_{\rs\ru} \xi_\rs + A_{\ro\ru} \xi_\ro | \\
&\qquad \geq (e^{\lambda/2} -\varepsilon) \delta | \xi_\ru| - \delta \varepsilon (|\xi_\rs| + |\xi_\ro|) \\
& \qquad > (e^{\lambda/2} -\varepsilon - \delta \varepsilon) (|\xi_\rs| + |\xi_\ro|) =: C_1 (\varepsilon) (|\xi_\rs| + |\xi_\ro|).
\end{split}
\end{equation}
The second inequality is obtained thanks to the estimate on $A_{\ru\ru}$ in (\ref{eq perturbe: lemme stabilite vois conique tilde - estimées matrices}) and the third one follows from our choice of $\xi \in \Ccal_\ru^{\delta} (x)$. Note that we implicitly show the following estimate which will imply the exponential estimate (as we will see later on) for $\delta \leq 1$:
\begin{equation}\label{inequality intermediate exponential bound}
|\xi_\ru^1| > C_1 (\varepsilon) |\xi_\ru| \geq \bigl(e^{\lambda/2} -2\varepsilon \bigr) |\xi_\ru|.
\end{equation}

Now, let us do a similar computation for the matrix $B$.
Again using Cauchy-Schwarz inequality on $J$ and using again estimates (\ref{eq perturbe: lemme stabilite vois conique tilde - estimées matrices}), we~deduce
\begin{equation}\label{eq: proof stability of conical neigh. - inequality 2 unstable}
\begin{split}
|\xi_\rs| + |\xi_\ro| &= | B_{\rs\rs} \xi_\rs^1 | + |B_{\rso} \xi_\rs^1 + B_{\ro\ro} \xi_\ro^1 | \geq |B_{\rs\rs} \xi_\rs^1| + |B_{\ro\ro} \xi_\ro^1| - |B_{\rso} \xi_\rs^1| \\
&\qquad \geq (e^{\lambda/2} -\varepsilon) |\xi_\rs^1| + (1 - \varepsilon)|\xi_\ro^1| - \varepsilon |\xi_\rs^1| \\
&\qquad \geq (1 - \varepsilon)(|\xi_\rs^1| + |\xi_\ro^1|) =: C_2 (\varepsilon) (|\xi_\rs^1| + |\xi_\ro^1|).
\end{split}
\end{equation}
where the last inequality holds for $\varepsilon$ small enough such that $e^{\lambda/2} - 2 \varepsilon \geq 1 - \varepsilon$.
Putting together equations \eqref{eq: proof perturb vois conique tilde - lower bound unstable} and \eqref{eq: proof stability of conical neigh. - inequality 2 unstable}, we~deduce
\begin{align*}
\delta|\xi_\ru^1| > C_1 (\varepsilon) C_2 (\varepsilon) (|\xi_\rs^1| + |\xi_\ro^1|).
\end{align*}
Finally, we~obtain $\xi^1 = \Phi_x^1 (\xi) \in \Ccal_\ru^{\delta'}(\varphi^1 (x))$ for
\[
\delta'(\varepsilon) = \frac{\delta }{C_1 (\varepsilon) C_2 (\varepsilon)},
\]
where the polynomial functions $C_1$ and $C_2$ (with respect to the variable $\varepsilon$) satisfy $C_1 (\varepsilon) \leq e^{\lambda}$ for every $\varepsilon >0$ and
\[
C_1 (\varepsilon) \underset{\varepsilon \to 0}{\to} e^{\lambda/2} \quad\text{and}\quad C_2 (\varepsilon) \underset{\varepsilon\to 0}{\to} 1.
\]
If we fix $\varepsilon$ sufficiently small so that
\begin{equation}\label{eq: proof big lemma - estimates C1 C2}
C_1 (\varepsilon) C_2 (\varepsilon) \geq e^{\lambda/3}, \quad C_1 (\varepsilon) \geq e^{\lambda/3},
\end{equation}
then we get for every $\delta\in(0,1]$ and for every $x \in \Ncal (\varepsilon)\cap W^\ru(K)$ the inclusion
\[
\Phi^{1}\left(\Ccal_{\ru}^{\delta}(x)\right) \subseteq \Ccal_{\ru}^{\delta'}(\varphi^{1}(x))
\]
for $\delta'(\varepsilon) \leq \delta e^{-\lambda/3}$. Moreover, we~deduce from (\ref{inequality intermediate exponential bound}) the bound
\begin{equation}\label{inequality intermediate exponential bound 2}
|\xi_\ru^{1}| \geq e^{\lambda/3} |\xi_\ru|
\end{equation}

\subsubsection*{Second inclusion, on weak unstable cones $\Ccal_{\ruo}^{\delta}$} Let us fix $x \in W^\ru (K) \cap \Ncal \cap \varphi^{-1}(\Ncal) $ and $\xi \in \Ccal_{\ruo}^{\delta} (x)$. By~definition of the conical neighborhood, we~have $\delta (|\xi_\ru| + |\xi_\ro|) > |\xi_\rs| $. The idea of the proof is very similar to the proof of the inclusion for the unstable conical neighborhood $\Ccal_{\ru}^{\delta}$. However, we~won't have an exponential estimate because of the neutral direction $\widetilde{E}_\ro^*$. Let us check the computations for this case. Using Cauchy-Schwarz inequality, we~have
\begin{align*}
|\xi_\ru^1| + |\xi_\ro^1| &= |A_{\ru\ru} \xi_\ru + A_{\rs\ru} \xi_\rs + A_{\ro\ru} \xi_\ro | + |A_{\rso} \xi_\rs + A_{\ro\ro} \xi_\ro | \\
&\geq |A_{\ru\ru} \xi_\ru| - |A_{\rs\ru} \xi_\rs| - |A_{\ro\ru} \xi_\ro | + | A_{\ro\ro} \xi_\ro | - | A_{\rso} \xi_\rs |\\
&\geq (e^{\lambda/2} - \varepsilon)|\xi_\ru| + (1 - \varepsilon) |\xi_\ro| - 2\varepsilon|\xi_\rs| \\
&\geq (1 - \varepsilon) (|\xi_\ru| + |\xi_\ro|) - 2\varepsilon |\xi_\rs|.
\end{align*}
Now, multiplying by $\delta \leq 1$ on both side of the previous inequalities and using the fact that $\xi \in \Ccal_{\ruo}^{\delta} (x)$, we~get
\begin{equation}\label{eq: lemme vois conique tilde - C3 2nd inclusion}
\begin{split}
\delta (|\xi_\ru^1| + |\xi_\ro^1|) &\geq (1 - 3\varepsilon) |\xi_\rs| =: C_3 (\varepsilon) |\xi_\rs|.
\end{split}
\end{equation}
It remains to find a lower bound for $|\xi_\rs|$. This case is much simpler since we have
\begin{align}\label{eq: lemme vois conique tilde - 3rd step, 2nd inclusion}
|\xi_\rs| = |B_{\rs\rs} \xi_\rs^1| \geq (e^{\lambda/2} -\varepsilon) |\xi_\rs^1| =: C_4 (\varepsilon) |\xi_\rs^1|.
\end{align}

Putting together equations \eqref{eq: lemme vois conique tilde - C3 2nd inclusion} and \eqref{eq: lemme vois conique tilde - 3rd step, 2nd inclusion}, we~deduce that
\begin{align*}
\delta(|\xi_\ru^1| + |\xi_\ro^1|) \geq C_3 (\varepsilon) C_4 (\varepsilon) |\xi_\rs^1|.
\end{align*}

Therefore, we~have $\xi^1 = \Phi_x^1 (\xi) \in \Ccal_\ru^{\delta'}(\varphi^1 (x))$ for
\[
\delta'(\varepsilon) = \frac{\delta }{C_3 (\varepsilon) C_4 (\varepsilon)},
\]
where the polynomial functions $C_3$ and $C_4$ (w.r.t the variable $\varepsilon$) satisfy $|C_3 (\varepsilon)| \leq 1$ for all $\varepsilon\in(0,1/3]$, $ C_3 (\varepsilon) \underset{\varepsilon \to 0}{\to} 1$ and $C_4 (\varepsilon) \underset{\varepsilon \to 0}{\to} e^{\lambda}$. For $\varepsilon$ sufficiently small so that
\begin{equation}\label{eq: proof big lemma - estimates C3 C4}
C_3 (\varepsilon) C_4 (\varepsilon) \geq e^{\lambda/3},
\end{equation}
we get for every $\delta\in[0,1]$ and for every $x \in \Ncal (\varepsilon) \cap W^\ru(K)$ the inclusion
\[
\Phi^{1}\left(\Ccal_{\ruo}^{\delta}(x)\right) \subseteq \Ccal_{\ruo}^{\delta'}(\varphi^{1}(x)),
\]
with $\delta' (\varepsilon) \leq \delta e^{-\lambda/3}$. It ends the proof of the inclusions along the unstable manifolds.

\medskip
Now, fix a value of $\varepsilon > 0$ sufficiently small so that (\ref{eq: proof big lemma - estimates C1 C2}) and (\ref{eq: proof big lemma - estimates C3 C4}) are satisfied. There exists $\varepsilon_1 > 0$ such that $\Vcal \cap W^\ru (K) \subset \Ncal (\varepsilon)$ hold for every unrevisited neighborhood $\Vcal$ contained in a $\varepsilon_1$-neighborhood of $K$. Let $\Vcal$ be an unrevisited neighborhood contained in a $\varepsilon_1$-neighborhood of $K$. Thanks to our choice of $\varepsilon$, the inclusions (\ref{lemma stab vois conique tilde - Unstable inlusions}) are satisfied on $\Vcal \cap W^\ru (K)$ and we would like to extend them to $\Vcal \cap \varphi^m(\Vcal)$ as stated in (\ref{lemma stab vois conique tilde - Unstable inlusions}).

Recall that each point of $\Vcal \cap \varphi^m (\Vcal)$ converges to the compact set $\overline{\Vcal} \cap W^\ru (K)$ as $m \to + \infty$ according to Lemma \ref{lemme: cv nonrevisite vers variete stable et instable}. The remaining of the proof consists in extending the inclusions in (\ref{lemma stab vois conique tilde - Unstable inlusions}) by continuity on $\Vcal \cap \varphi^m (\Vcal)$ for $m$ large enough. Fix $\delta_0\in(0,1]$. By~continuity of the maps
\begin{align*}
F (x,\xi) = \Bigl(\frac{|\xi_\rs| + |\xi_\ro|}{|\xi_\ru|}\Bigr)^{-1} \Bigl(\frac{|\xi_\rs^1| + |\xi_\ro^1|}{|\xi_\ru^1|} \Bigr), \quad G (x,\xi) = \frac{|\xi_\rs^1| + |\xi_\ro^1|}{|\xi_\ru^1|}, \\
\widetilde{F} (x,\xi) = \Bigl(\frac{|\xi_\rs|}{|\xi_\ru| + |\xi_\ro|}\Bigr)^{-1} \Bigl(\frac{|\xi_\rs^1|}{|\xi_\ru^1| + |\xi_\ro^1|} \Bigr), \quad \widetilde{G} (x,\xi) = \frac{|\xi_\rs^1|}{|\xi_\ru^1| + |\xi_\ro^1|},
\end{align*}
there exists an integer $m_{\delta_0} \in \N$ such that the claimed inclusions are satisfied on $\Vcal \cap \varphi^{m_{\delta_0}} (\Vcal)$ for the parameter $\delta' = \delta e^{-\lambda/4}$. This ends the proof of \eqref{stabilitei}.

\subsubsection*{Exponential estimate (proof of \eqref{stabiliteii})} Applying \eqref{stabilitei} to $\delta=\delta_0\in(0,1]$, we~get the existence of $m_{\delta} \in \N$ such that the inclusions (\ref{lemma stab vois conique tilde - Unstable inlusions}) are satisfied on $\Vcal \cap \varphi^{m_{\delta}} (\Vcal)$. Extending~(\ref{inequality intermediate exponential bound 2}) by continuity, we~can find $m_0 \in \N$ such that for every $x \in \Vcal \cap \varphi^{m_0} (\Vcal)$ and for every $\xi \in \Ccal_{\ru}^{\delta} (x)$, we~have
\[
|\xi_u^1| \geq e^{2\lambda / 7 } |\xi_u|.
\]
Up to considering a larger $m_{\delta}$, we~can assume that $m_{\delta} \geq m_0$. The exponential estimate follows by equivalence of the norms $|\xi |_u:= |\xi_u|$ and $|.|$ on the conical neighborhoods $\Ccal_\ru^{\delta}(x)$ and $\Ccal_\ru^{\delta}(\varphi^1(x))$. Indeed, we~have $\xi^1 \in \Ccal_\ru^{\delta} (\varphi^1 (x))$ thanks to our previous analysis and therefore
\[
|\Phi^1 (x,\xi)| = |\xi^1| \geq \frac{e^{2/7\lambda}}{1 + \delta}\, |\xi| \geq e^{\lambda /4} |\xi|.
\]
for every $\delta \leq \delta_1:= \min (1, e^{2/7\lambda} - 1)$ and every $x \in \Vcal \cap \varphi^{m_{\delta}} (\Vcal)$.
\end{proof}

The following corollary states in a quantitative manner that, if the trajectory of a point $(x,\xi)$ stays for a long time near a basic set, then the fiber part of $\Phi^t (x,\xi)$ gets attracted to the $\widetilde{E}_\ru^*$ or $\widetilde{E}_{\ruo}^*$ distribution:

\begin{corollary}\label{corollary big lemma} Let $K$ be a basic set and let $\Ncal$ be the open neighborhood appearing in \eqref{decomposition cotangent au voisinage d un compact basique hyperbolique}. There exists $\varepsilon_1 >0$ such that, for every unrevisited neighborhood $\Vcal\subset \varphi^1(\Ncal) \cap\nobreak\varphi^{-1}(\Ncal)$ contained in an $\varepsilon_1$-neighborhood of $K$, the following hold:
\begin{itemize}
\item For every $\delta,\delta'\in(0,1]$ satisfying $\delta' \leq \delta$, one can find $m_{0}\geq 0$ such that, for every $m \geq m_{0}$ and for every $x \in \varphi^{-m}(\Vcal)\cap \Vcal$, one has
\[
\xi \notin \Ccal_{\rso/\rs}^{\delta} (x) \implies \Phi^{ m} (x,\xi) \in \Ccal_{\ru/\ruo}^{\delta'} (\varphi^{m} (x)).
\]
\item For every $\delta\in(0,1]$, there exist $C > 0$, $m_1 \in \N$ such that, for every $m \geq m_1$, for every $x \in \varphi^{-m}(\Vcal)\cap \Vcal$ and for every $\xi \in T_x^* M$ such that $ \xi \notin \Ccal_{\rso}^{\delta} (x) $, we~have the inequality
\begin{equation}\label{eq: corollary big lemma - statement exponential estimates}
|\Phi^{ m} (\xi)| \geq C e^{m \lambda /4} |\xi|,
\end{equation}
with $\lambda$ being the constant appearing in the definition \eqref{eq: def extended metric near basic set}.
\end{itemize}
\end{corollary}

The idea of the proof is the following. Thanks to Lemma \ref{stabilite voisinage conique tilde pour Phi1}, we~know that the conical neighborhoods are $\Phi^1$-stable near the unstable manifold of the basic set, so~our goal is to use the different unrevisited neighborhoods represented in Figure \ref{fig: unrevisited neighborhoods near a basic set} to show that we can iterate the lemma.

\begin{proof}[Proof of the corollary]
Fix $\delta,\delta'\in(0,1]$ satisfying $\delta' \leq \delta$ and let $\Vcal$ be some unrevisited $\varepsilon_1$\nobreakdash-neighborhood of $K$ contained in $\varphi^1(\Ncal) \cap \varphi^{-1} (\Ncal)$ so that we can apply Lemma \ref{stabilite voisinage conique tilde pour Phi1} for the flows $\varphi^t$ and $\varphi^{-t}$ (with $\varepsilon_1$ given by the lemma). Doing so for the flows~$\varphi^t$ and~$\varphi^{-t}$, we~obtain the existence of an integer $m_{\delta'}$ as in the statement of the lemma (we take $\delta_0:= \delta'$ and we assume that $m_{\delta'}$ is the same integer obtained for~$\varphi^t$ and~$\varphi^{-t}$). Now, fix an even integer $m_0 \geq 2m_{\delta'}$ which will be chosen sufficiently large later on. For all $m \geq m_{0}$ and all $x \in \Vcal \cap \varphi^{-m} (\Vcal)$, we~have the inclusions
\[
\forall k \in \lcr 0, m_0 /2 \rcr, \quad \varphi^{k} (x) \in \varphi^k (\Vcal) \cap \varphi^{-m + k} (\Vcal) \subset \Vcal \cap \varphi^{- m_{0}/2} (\Vcal)
\]
and
\[
\forall k \in \lcr m_0/2, m \rcr, \quad \varphi^{k} (x) \in \varphi^k (\Vcal) \cap \varphi^{-m + k} (\Vcal) \subset \varphi^{m_{0}/2} (\Vcal) \cap \Vcal,
\]
see Figure \ref{fig: unrevisited neighborhoods near a basic set}. The remaining of the proof consists in interacting $m$ times the lemma for $(x,\xi) \in T^* M$ such that $x \in \Vcal \cap \varphi^{-m} (\Vcal)$ and $\xi \notin \Ccal_{\rso/\rs}^{\delta} (x)$: $m_0/2$ times for the backward flow $\varphi^{-t}$ when the orbit of $x$ belongs to $\Vcal \cap \varphi^{- m_0/2} (\Vcal)$ and $m - m_0/2$ times for $\varphi^t$ when the orbit of $x$ goes in $ \varphi^{m_{0}/2} (\Vcal) \cap \Vcal$. For $m$ sufficiently large, we~will get the result. Let us consider $(x,\xi)$ as above. Two analyses are needed: first for $0 \leq k \leq m_0/2$ and then for $m_0/2 \leq k \leq m$.
\begin{itemize}
\item First, since $\varphi^k (x) \in \Vcal \cap \varphi^{- m_{0}/2} (\Vcal) \subset \Vcal \cap \varphi^{- m_{\delta'}} (\Vcal)$ for all $k \in \lcr 0, m_0 /2-1 \rcr$, a~straight application of the lemma for the flow $\varphi^{-t}$ instead of $\varphi^t$ yields
\[
\xi \notin \overline{\Ccal_{\rso/\rs}^{\delta} (x)} \implies \Phi^{m_0/2} (x,\xi) \notin \overline{\Ccal_{\rso/\rs}^{\delta_{m_0/2}} (\varphi^{m_0/2}(x))}
\]
where $\delta_{m_0/2} = \min (e^{\lambda m_0 /8} \delta, 1)$. In~particular, if we choose $m_0 \in \N$ large enough so that
\[
e^{-m_0 \lambda /8} \leq \delta' \leq \delta,
\]
then we get $\delta_{m_0/2} = 1$,
\[
\Phi^{m_0/2} (x,\xi) \notin \overline{\Ccal_{\rso/\rs}^{1} (\varphi^{m_0/2} (x))} \quad\text{and thus } \Phi^{m_0/2} (x,\xi) \in \Ccal_{\ru/\ruo}^{1} (\varphi^{m_0/2} (x)).
\]
\item For every $k\in\lcr m_0/2,m-1\rcr$, we~can apply the lemma\footnote{Note that we use here the uniformity in $\delta \in [\delta', 1]$ stated in Lemma \ref{stabilite voisinage conique tilde pour Phi1}.} to the flow $\varphi^t$ and it yields
\begin{multline*}
\varphi^k (x) \in \Vcal \cap \varphi^{m_{\delta'}} (\Vcal) \quad \text{and} \quad \Phi^k (x,\xi) \in \Ccal_{\ru/\ruo}^{\mu^{k-m_0/2}} (\varphi^k (x)) \\
\implies \Phi^{k+1}(x,\xi) \in \Ccal_{\ru/\ruo}^{\mu^{k+1-m_0/2}} (\varphi^{k+1} (x)),
\end{multline*}
with $\mu = e^{-\lambda/4}$. Finally, we~obtain
$ \Phi^{m}(x,\xi) \in \Ccal_{\ru/\ruo}^{\mu^{m - m_0/2}} (\varphi^{m} (x)).$ Thanks to our choice of $m$ and $m_0$, we~have $\mu^{m - m_0/2} \leq \mu^{m_0/2} = e^{-m_0 \lambda /8} \leq \delta'$. This gives the first point.
\end{itemize}

To prove the exponential estimate (\ref{eq: corollary big lemma - statement exponential estimates}), we~first need to introduce the constant $\delta_1\in(0,1]$ given by  Lemma \ref{stabilite voisinage conique tilde pour Phi1}\eqref{stabiliteii}. From the above analysis, there exists an integer $m_1 \geq m_0$ such that for every $m \geq m_1$ and for every $k\in\lcr m_1,m\rcr$, we~have $ \Phi^k (x,\xi) \in \Ccal_{\ru/\ruo}^{\delta_1} (\varphi^k (x))$. The result is a direct application of Lemma \ref{stabilite voisinage conique tilde pour Phi1}\eqref{stabiliteii} and the constant $C$ is obtained by continuity of $\Phi^{m_1}$.
\end{proof}

\subsection{Proof of proposition \ref{prop compactness}: compactness.}
\label{subsection: proof compactness result}

We are now in position to prove the compactness of $\Sigma_{\ru/\ruo}$. Compactness of $\Sigma_{\rs/\rso}$ can be proved similarly if we exchange~$\varphi^1$ with~$\varphi^{-1}$. Let $((x_m, \xi_m))_{m \in \N}$ be a sequence of elements in $\Sigma_{\ru/\ruo}$. Up to extraction of a subsequence, we~can assume that there exists an integer $j\in\lcr1,N\rcr$ such that every point $x_m$ belongs to $W^\rs (K_j)$. Up to another extraction, we~can suppose that $((x_m,\xi_m))_m$ has a limit in $S^*M$ that we denote by $(x_{\infty},\xi_{\infty}) \in S^*M$. Our goal will be to prove that $(x_{\infty}, \xi_{\infty})$ belongs to $\Sigma_{\ru/\ruo}$. First of all, we~know from our assumption on~$x_m$ that the limit $x_{\infty}$ must lie in the closure of the set $ W^\rs (K_j)$, \ie $x_{\infty} \in \overline{W^\rs (K_j)}$. Since $\overline{W^\rs (K_j)} = \bigcup_{j_0, K_{j_0} \leq K_j} W^\rs (K_{j_0})$, we~can find some integer $j_0 \leq j$ such that $x_{\infty} \in W^\rs (K_{j_0})$. Now, let us assume by contradiction that $(x_{\infty},\xi_{\infty}) \notin \Sigma_{\ru/\ruo}$.

To obtain a contradiction, we~split the analysis in four steps. First, we~deal with the case where $j = j_0$. In~steps 2--4, we~treat the case $j \neq j_0$. Precisely, in the second step, we~construct by induction a family of integers $j_0 < j_1 < \cdots < j_{\ell} = j$ such that for each $k\in\lcr1,\ell\rcr$ one can find an element $(x_{\infty}^{(k)}, \xi_{\infty}^{(k)})$ with $x_{\infty}^{(k)} \in W^\rs (K_{j_k}) \cap W^\ru (K_{j_{k-1}})$ and a sequence $(x_m^{(k)},\xi_m^{(k)}) = \widetilde{\Phi}^{\tau_{m,k}} (x_m,\xi_m)$, for some parameter $ \tau_{m,k} \geq 0$ satisfying $\tau_{m,k +1} - \tau_{m,k} \underset{m\to+\infty}{\to} + \infty$, which accumulates to $(x_{\infty}^{(k)},\xi_{\infty}^{(k)})$ as $m$ tends to infinity. In~a third step, we~apply our previous analysis near the basic sets of Corollary \ref{corollary big lemma} to prove that if $(x_{\infty}^{(k)},\xi_{\infty}^{(k)}) \notin \Sigma_{\ru/\ruo}$, then $(x_{\infty}^{(k+1)}, \xi_{\infty}^{(k+1)})$ must belong to $\Sigma_{\rso/\rs}$. In~the last step, we~deduce the expected contradiction.

\begin{proof}[Step 1: we assume that $j = j_0$]
The contradiction follows from an application of Lemma \ref{lemma: continuity of extension of distributions} and more precisely from the continuity of $\Sigma_{\ru/\ruo}$ on $W_{\varepsilon}^\rs (K_j)$ for $\varepsilon < \varepsilon_0$ as in the statement of Lemma \ref{lemma: continuity of extension of distributions} (for example $\varepsilon = \varepsilon_0/2$), where $\varepsilon_0$ is given by the stable manifold theorem. Let $\Vcal$ be an unrevisited neighborhood of $K_j$ such that $\Vcal \subset B_g (K_j, \varepsilon') = \{z \in M,\, d_g (z, K_j) < \varepsilon' \}$ for $\varepsilon' > 0$ which will be chosen sufficiently small. If $x_{\infty} \in W^\rs (K_j)$, then there exists $T(\varepsilon') \geq 0$ such that $\varphi^{T (\varepsilon')} (x_{\infty}) \in W^\rs (K_j) \cap \Vcal$ (by definition of the stable set).
Since $x_m \in W^\rs (K_j)$ and since the sequence $(x_m)_{m \in \N}$ converges to $x_{\infty}$ as $m$ tends to $+\infty$, we~can find an integer $m_0 (\varepsilon') \geq 0$ such that $\varphi^{T(\varepsilon')} (x_m) \in W^\rs (K_j) \cap \Vcal$ for every $m \geq m_0 (\varepsilon')$. Since~$\Vcal$ is unrevisited, we~get that
\[
\forall m \in \lcr m_0 (\varepsilon'), +\infty \rcr,\;\forall k \geq 0, \quad \varphi^{T(\varepsilon') + k} (x_m) \in W^\rs (K_j) \cap \overline{\Vcal}.
\]
In particular, we~get that
\[
\forall m \in \lcr m_0 (\varepsilon'), +\infty \rcr,\; \forall k \geq 0, \quad d_g \bigl(\varphi^{T(\varepsilon') + k} (x_m), K_j \bigr) \leq \varepsilon'.
\]

As an application of the shadowing lemma \cite[Th.\,5.3.3]{FH} similar to the one used in the proof of the In-Phase theorem \cite[Th.\,5.3.25]{FH}, we~get that $\varphi^{T(\varepsilon')} (x_m) \in W_{\varepsilon}^\rs (K_j)$ for all $m \in \lcr m_0 (\varepsilon'), +\infty \rcr$ as soon as $\varepsilon'$ is sufficiently small (depending on $\varepsilon$).

From the continuity of $\Sigma_{\ru/\ruo}$ on $W_{\varepsilon}^\rs (K_j)$, we~deduce that $(x_{\infty}, \xi_{\infty})$ must belong to $\Sigma_{\ru/\ruo}$. Therefore, we~get
\[
(x_{\infty}, \xi_{\infty}) \in \Sigma_{\ru/\ruo},
\]
which is in contradiction with the assumption on $(x_{\infty}, \xi_{\infty})$. \emph{From now on}, we~assume that
$j_0 < j$.

\subsubsection*{Step 2: construction of the sequence by induction}
\begin{figure}[t]
\centering
\def\svgscale{0.6}
\includesvg{compactness}
\caption{Illustration of some element used in the proof}
\end{figure}
Since $j_0$ has already been defined, we~just have to set $(x_{\infty}^{(0)}, \xi_{\infty}^{(0)}):= (x_{\infty}, \xi_{\infty})$ to end the base case. Now, let us assume the integers $j_1 < j_2 < \cdots < j_k (<j)$, the elements of the unitary cotangent bundle $(x_{\infty}^{(1)}, \xi_{\infty}^{(1)}),\dots, (x_{\infty}^{(k)}, \xi_{\infty}^{(k)})$ and the parameters $\tau_{m,1},\dots, \tau_{m,k}$ are constructed as in the above discussion. Our goal is to define a sequence $(x_{m}^{(k+1)}, \xi_{m}^{(k+1)}):= \widetilde{\Phi}^{\tau_{m,k+1}} (x_m,\xi_m)$, the family $\tau_{m,k+1} \geq 0$ such that $\tau_{m,k+1} - \tau_{m,k}$ tends to $+\infty$ as $m$ goes to $+\infty$ and to exhibit a new accumulation point $(x_{\infty}^{(k+1)}, \xi_{\infty}^{(k+1)})$ with $x_{\infty}^{(k+1)} \in W^\ru(K_{j_k})\cap W^\rs(K_{j_{k+1}})$.

\subsubsection*{Definition of $\tau_{m,k+1}$} Let us consider an unrevisited neighborhood $\Vcal$ of $K_{j_k}$ so that $\overline{\Vcal} \cap \Omega = K_{j_k}$ and sufficiently small to be in the range of application of Corollary~\ref{corollary big lemma}. Recall that $x_{\infty}^{(k)}$ belongs to $W^\rs (K_{j_k})$, so there exists $T > 0$ such that $\varphi^T (x_{\infty}^{(k)}) \in W^\rs (K_{j_k}) \cap \Vcal$. Since every $x_{1}^{(k)}, x_{2}^{(k)},\dots, x_{m}^{(k)},\dots$ are in $W^\rs (K_{j})$ and since the sequence converges to $x_{\infty}^{(k)}$ as $m \to + \infty$, there exists an integer $m_0 \geq 0$ such that for every $m \geq m_0$ we have $\varphi^T (x_m^{(k)}) \in \Vcal$.
To lighten notations, let us denote by $y_m$ the point $\varphi^T (x_{m}^{(k)})$ and by $\eta_m$ the cotangent vector $\kappa ((D \varphi^{T} (x_m^{(k)})^{-1})^{\top} (\xi_m^{(k)})) \in S_{y_m}^* M$ for every $m \in \N \cup \{ \infty \}$.
For every $m \geq m_0$, we~can define the exit time of the unrevisited set~$\Vcal$ for the point~$y_m$ as follows:
\[
\tau_m:= \inf \{ p \in \N, \: \varphi^p (y_m) \notin \Vcal \} -1.
\]
It is well-defined by definition of $y_m$ and it depends on $k$. Note that $\tau_m$ is finite for all $m$ since $y_m \notin W^\rs (K_{j_k})$ (because $y_m \in W^\rs (K_j)$ with $j > j_k$) and since $\bigcap_{p \in \N} \varphi^{-p} (\Vcal) = W^\rs (K_{j_k}) \cap \Vcal$. Also, the convergence $y_m \underset{m \to +\infty}{\to} y_{\infty} \in W^\rs (K_{j_k})$ implies that $\tau_m$ goes to $+\infty$ when $m$ tends to $+\infty$. Now, we~define $y_m^{(1)} = \varphi^{\tau_m} (y_m)$ and $\eta_m^{(1)} = \kappa ((D\varphi^{\tau_m} (y_m)^{-1})^{ \top} (\eta_m))$. Up to extraction, we~can assume that $(y_m^{(1)}, \eta_m^{(1)})$ converges to an element $(y_{\infty}^{(1)},\eta_{\infty}^{(1)}) \in S^*M$ by compactness of $S^* M$. Let us define
\[
(x_m^{(k+1)}, \xi_m^{(k+1)}):= (y_m^{(1)}, \eta_m^{(1)}) = \widetilde{\Phi}^{\tau_{m,k} + \tau_m + T} (x_m, \xi_m)
\]
for every $m \in \N \cup \{ \infty \}$ and also
\[
\tau_{m, k+ 1}:= \tau_{m,k} + \tau_m + T \geq 0.
\]
Thanks to Lemma \ref{lemme: cv nonrevisite vers variete stable et instable} and to Remark \ref{lemma: clos. of stable man on unrevisited neigh.}, the point $y_{\infty}^{(1)}$ belongs to $\overline{W^\ru (K_{j_k})\cap \Vcal} = W^\ru (K_{j_k})\cap \overline{\Vcal}$. Indeed, it follows from the convergence $y_m^{(1)} \underset{m \to +\infty}{\to} y_{\infty}^{(1)}$ together with the fact that $y_m^{(1)} \in \varphi^{\tau_m} (\Vcal) \cap \Vcal$ and $\tau_m \underset{m \to +\infty}{\to} +\infty$. Moreover, we~have $\varphi^{1} (y_{\infty}^{(1)}) \notin \Vcal$ and thus $y_{\infty}^{(1)} \notin K_{j_k}$ because $\varphi^1 (y_m^{(1)}) \notin \Vcal$ by definition. Furthermore, there exists an integer $ j_{k +1}\in\lcr1,N\rcr$ such that $y_{\infty}^{(1)} \in W^\rs (K_{j_{k+1}})$. Now, using the properties of the Smale order relation given in Theorem \ref{thm Smale} together with the total order on the indices of Section~\ref{subsubsection total order relation}, we~deduce from $y_{\infty}^{(1)} \in W^\ru (K_{j_{k}}) \cap W^\rs (K_{j_{k+1}})$ and $y_{\infty}^{(1)} \in \overline{W^\rs (K_j)}$ (given by construction) that $j_k < j_{k+1} \leq j$.
This ends the induction step. Note that the algorithm stops once we have reached $K_j$, in which case we have defined $j_0 < j_1 < \cdots < j_{\ell} = j$.

\subsubsection*{Step 3: local analysis near a basic set}
For every $k\in\lcr1,\ell-1\rcr$, we~will prove the implication
\begin{equation}\label{eq: compactness - proof implication}
(x_{\infty}^{(k)}, \xi_{\infty}^{(k)}) \!\notin\! \Sigma_{\ru/\ruo} \text{ and } \tau_{m,k +1} - \tau_{m,k} \!\underset{m\to+\infty}{\to}\! + \infty \implies (x_{\infty}^{(k+1)}, \xi_{\infty}^{(k+1)}) \!\in\! \Sigma_{\rso/\rs}.
\end{equation}
Let us fix $k\in\lcr1,\ell-1\rcr$. To simplify, we~will use the same notations as the one used in the previous induction. By~definition, we~have $(x_\infty^{(k+1)},\xi_\infty^{(k+1)}) \in \Sigma_{\rso/\rs}$ if and only if $\eta_{\infty}^{(1)} \in E_{\ru/\ruo}^* (y_{\infty}^{(1)})$. Since $y_{\infty} \in W^\rs (K_{j_k}) \cap \overline{\Vcal}$, since $x_{\infty}^{(k+1)} = y_{\infty}^{(1)} \in W^\ru (K_{j_k}) \cap \overline{\Vcal} \setminus \varphi^{-1} (\Vcal)$ and since the conical neighborhoods $\Ccal_{*}^{\delta}$ are well-defined on $\overline{\Vcal}$, we~have by construction
\[
E_{\rso/\rs}^* (y_{\infty}) = \widetilde{E}_{\rso/\rs}^* (y_{\infty}) \quad \text{and} \quad E_{\ru/\ruo}^* (y_{\infty}^{(1)}) = \widetilde{E}_{\ru/\ruo}^* (y_{\infty}^{(1)}).
\]
Assume that $\eta_{\infty} \notin E_{\rso/\rs}^* (y_{\infty})$ and let us prove that $\eta_{\infty}^{(1)}$ belongs to $E_{\ru/\ruo}^* (y_{\infty}^{(1)})$ or, equivalently, that for any $\delta' > 0$ we have
\begin{equation}\label{eq: compactness - phase at the exit}
\eta_{\infty}^{(1)} \in \Ccal_{\ru/\ruo}^{\delta'} (y_{\infty}^{(1)}).
\end{equation}
Fix $\delta' > 0$. By~hypothesis, we~can find a small constant $\delta > 0$ such that
\[
\eta_{\infty} \notin \Ccal_{\rso/\rs}^{\delta} (y_{\infty}).
\]
Since $\tau_{m,k +1} - \tau_{m,k} \underset{m\to+\infty}{\to} + \infty$, we~can apply Corollary \ref{corollary big lemma} which gives \eqref{eq: compactness - phase at the exit} and thus $(x_{\infty}^{(k+1)}, \xi_{\infty}^{(k+1)}) \in \Sigma_{\rso/\rs}$.

\subsubsection*{Step 4: at the end of the algorithm}
According to step 2, we~have defined $(x_{\infty}^{(\ell)}, \xi_{\infty}^{(\ell)}) \in S^*M$ with $x_{\infty}^{(\ell)} \in W^\rs (K_j)$ as well as $x_{m}^{(\ell)} \in W^\rs (K_j)$ for every $m \in \N$. Furthermore, since we have assumed at the beginning $(x_{\infty},\xi_{\infty}) \notin \Sigma_{\ru/\ruo}$, we~deduce from the \eqref{eq: compactness - proof implication} of step 3 that\vspace*{-5pt}
\begin{multline*}
(x_{\infty},\xi_{\infty}) \notin \Sigma_{\ru/\ruo} \implies (x_{\infty}^{(1)}, \xi_{\infty}^{(1)}) \in \Sigma_{\rso/\rs} \implies (x_{\infty}^{(1)}, \xi_{\infty}^{(1)}) \notin \Sigma_{\ru/\ruo} \implies \cdots \\
\implies (x_{\infty}^{(\ell)}, \xi_{\infty}^{(\ell)}) \in \Sigma_{\rso/\rs},
\end{multline*}
where the strong transversality assumption is used extensively through the relation $\Sigma_{\ru/\ruo} \cap \Sigma_{\rso/\rs} = \emptyset$. However, proceeding similarly to the first step, we~also have the convergence\vspace*{-3pt}\enlargethispage{.5\baselineskip}
\[
\Sigma_{\ru/\ruo} \ni (x_m^{(\ell)},\xi_m^{(\ell)}) = \widetilde{\Phi}^{\tau_{m,\ell}} \underbrace{(x_m,\xi_m)}_{\in \Sigma_{\ru/\ruo}} \underset{m \to +\infty}{\to} (x_{\infty}^{(\ell)}, \xi_{\infty}^{(\ell)}) \in \Sigma_{\rso/\rs},
\]
which leads by continuity of $\Sigma_{\ru/\ruo}$ on $W_{\varepsilon}^\rs (K_j)$ for $\varepsilon \ll 1$ (according to Lemma \ref{lemma: continuity of extension of distributions}) to $(x_{\infty}^{(\ell)}, \xi_{\infty}^{(\ell)}) \in \Sigma_{\ru/\ruo}$ and thus gives the expected contradiction.

\subsubsection*{Conclusion}
In all cases we obtained a contradiction and therefore we must have $(x_{\infty},\xi_{\infty}) \in \Sigma_{\ru/\ruo}$.
\end{proof}

Let us state a corollary which will be used to construct the map $f$ which appears in the definition of the escape function---see Proposition \ref{prop escape function}.

\begin{corollary}\label{Corollary compactness proposition}
For every $\varepsilon > 0$, there exists $\varepsilon' > 0$ such that, for every point $x$ which is $\varepsilon'$-close to $K$, every element of $\kappa \bigl(E_{\ru/\ruo}^* (x) \setminus 0_M\bigr)$ is $\varepsilon$-close to the compact set $\bigcup_{z \in K} \kappa\bigl(E_{\ru/\ruo}^* (z)\setminus 0_M \bigr)$ for the distance $d_{S^*M}$ associated with the Sasaki metric on $S^*M$.
\end{corollary}

\begin{proof}
By contradiction, assume that there exists $\varepsilon > 0$ such that for every $m \in \N^*$, there exist $x_m$ and $\xi_m \in \kappa \bigl(E_{\ru/\ruo}^* (x_m) \setminus 0_M \bigr)$ which satisfy\vspace*{-3pt}
\begin{equation}\label{eq: corollary compactness}
d_g (x_m, K) \leq 1/m \quad \text{and} \quad d_{S^* M} \Bigl((x_m,\xi_m), \bigcup_{z \in K} \kappa\bigl(E_{\ru/\ruo}^* (z)\setminus 0_M \bigr) \Bigr) \geq \varepsilon.
\end{equation}
By compactness of $S^* M$, we~can extract a converging subsequence $((x_{m_k},\xi_{m_k}))_k$ which converges to an element $(x_{\infty},\xi_{\infty}) \in S^* M$ as $k \to + \infty$. Taking the limit in the inequalities (\ref{eq: corollary compactness}) implies that $x_{\infty} \in K$ and $\xi_{\infty} \notin \kappa \bigl(E_{\ru/\ruo}^* (x_{\infty}) \setminus 0_M \bigr)$. However, the set $\Sigma_{\rs/\rso} = \kappa \bigl(E_{\ru/\ruo}^* \setminus 0_M \bigr)$ is a compact set thanks to the compactness proposition \ref{prop compactness}. Therefore, we~must have $(x_{\infty},\xi_{\infty}) \in \Sigma_{\rs/\rso}$ or, equivalently, $\xi_{\infty} \in \kappa \bigl(E_{\ru/\ruo}^* (x_{\infty}) \setminus 0_M \bigr)$, which gives the contradiction.
\end{proof}

\Subsection{The compact sets are attracting and repelling sets for the Hamiltonian flow}
\label{subsection: proof Lemma attracting and repelling sets for the Hamiltonian flow}

\begin{figure}[!t]
\centering
\def\svgscale{1.1}
\includesvg{cvdualdistrib}
\caption{Convergence of the dual unstable distribution. The point~$z$ denotes an hyperbolic fixed point and the point $x$ belongs to the stable manifold of $z$.}
\label{fig: cv of dual distribution for Morse-Smale gradient flows}
\end{figure}

Now that we have proved the compactness of $\Sigma_{\ruo/\ru}$ and $\Sigma_{\rs/\rso}$, it remains to show that $(\Sigma_{\rs/\rso}, \Sigma_{\ruo/\ru})$ defines an attractor-repeller pair. Equivalently, we~have to prove Lemmas \ref{lemma attractor and repeller on the phase space} and \ref{lemma invariant neighborhoods}. In~the upcoming argument, we~generalize the convergence presented in Figure \ref{fig: cv of dual distribution for Morse-Smale gradient flows} to Axiom~A flows satisfying the strong transversality assumption~(\ref{Transversality assumption - first}). The proof is very similar to that given in \cite[Lem.\,2.10\,(4), p.\,13]{dyatlov-guillarmou_2016}, except we allow our phase point $\xi$ to have a neutral component, \ie $\xi (V)$ can be non zero.

\begin{proof}[Proof of Lemma \ref{lemma attractor and repeller on the phase space}] Let us take $(x,\xi) \in S^* M$ such that $(x,\xi) \notin \Sigma_{\ruo/\ru}$ and let us prove that $\widetilde{\Phi}^t (x, \xi)$ tends to $\Sigma_{\rs/\rso}$ as $t \to +\infty$. The case $t\to-\infty$ is obtained by applying the result to the vector field $-V$. From the spectral decomposition of~$M$ (Lemma \ref{lemma: spectral decomposition of the manifold}), we~have $M = \bigsqcup_{i = 1}^N W^\rs (K_i)$ and consequently there exists a unique integer $i\in\lcr1,N\rcr$ such that $x \in W^\rs (K_i)$. Let $\Vcal$ be an unrevisited neighborhood of $K_i$ sufficiently small to be in the range of application of Corollary \ref{corollary big lemma}. Also, there exists $T > 0$ such that $\varphi^t (x) \in W^\rs (K_i) \cap \Vcal$ for every $t \geq T$. Now, let us recall that for every $y \in W^\rs (K_i) \cap \Vcal$ we have
\[
\widetilde{E}_\rs^* (y) = E_\rs^* (y) \quad \text{and} \quad \widetilde{E}_{\rso}^* (y) = E_{\rso}^* (y)
\]
by construction. Since $\Sigma_{\ruo/\ru}$ is $\widetilde{\Phi}^t$-invariant, the condition $(x,\xi) \notin \Sigma_{\ruo/\ru} $ implies that $\widetilde{\Phi}^t (x,\xi) \notin \Sigma_{\ruo/\ru} = \kappa(E_{\rso/\rs}^*)$ for every $t \geq 0$. Therefore, there exists $\delta > 0$ such that $ \widetilde{\Phi}^T (x,\xi) \notin \Ccal_{\rso/\rs}^{\delta}$ and Corollary \ref{corollary big lemma} implies that
\[
d_{S^* M} \Bigl(\widetilde{\Phi}^t (x,\xi), \bigcup_{z \in K_i} \kappa \bigl(\Ccal_{\ru/\ruo}^{\delta'} (z) \bigr) \Bigl) \underset{t \to +\infty}{\to} 0, \qquad \forall \delta' > 0,
\]
or, equivalently, that
\[
d_{S^* M} \Bigl(\widetilde{\Phi}^t (x,\xi), \bigcup_{z \in K_i } \kappa \bigl(E_{\ru/\ruo}^* (z)\bigr) \Bigr) \underset{t \to +\infty}{\to} 0.
\]
This ends the proof.
\end{proof}

\Subsection{Proof of Lemma \ref{lemma invariant neighborhoods}: invariant neighborhoods for the Hamiltonian flow}\label{subsection: Lemma invariant neighborhoods for the Hamiltonian flow}

Now, we~need to find invariant neighborhoods of $\Sigma_{\rs/\rso}$ and $\Sigma_{\ruo/\ru}$, \ie we have to prove Lemma \ref{lemma invariant neighborhoods}. The idea of the proof follows from two remarks:

\begin{itemize}
\item In the case where $K$ is an attractor, we~can construct a $\widetilde{\Phi}^1$-stable neighborhood of $\{(x,\xi)\in\Sigma_{\rs/\rso}: x\in K\}$ as follows: if $\Vcal$ is an unrevisited neighborhood of $K$ small in the sense that the conical neighborhood $\Ccal_*^{\delta}$ are well-defined on $\Vcal$, then for all $\delta\in(0,1]$ the sets
\[
\bigcup_{z \in \Vcal} \kappa \left(\Ccal_{\ru}^{\delta}(z) \right) \quad \text{and} \quad \bigcup_{z \in \Vcal} \kappa \left(\Ccal_{\ruo}^{\delta}(z) \right)
\]
are $\widetilde{\Phi}^1$-stable neighborhoods of $\{(x,\xi)\in\Sigma_{\rs}: x\in K\}$ and $\{(x,\xi)\in\Sigma_{\rso}: x\in K\}$.
\item In the general case, thanks to Lemma \ref{stabilite voisinage conique tilde pour Phi1} we can note that if $\Vcal$ is an unrevisited neighborhood of $K$ (to be in the setting of this lemma), then the set
\[
\bigcup_{z \in \varphi^m (\Vcal)\cap \Vcal} \kappa \bigl(\Ccal_{\ru/\ruo}^{\delta}(z) \bigr)
\]
defines an unrevisited neighborhood containing $\bigcup_{x \in W^\ru (K) \cap \Vcal } \kappa\bigl(E_{\ru/\ruo}^* (x) \bigr)$ for every $\delta\in(0,1]$ and for every $m \gg 1$.
\end{itemize}

Our strategy will be to construct invariant neighborhoods $\Sigma_{\rs/\rso}$ and $\Sigma_{\ruo/\ruo}$ by induction from the unrevisited neighborhoods exhibited above. In~the upcoming proof, we~construct a filtration\footnote{Even if the definition of filtration was given for the flow $\varphi^t$, it can be adapted for the Hamiltonian without too much effort, see \cite{shub} or \cite{smale}.} by open sets for the Hamiltonian flow $\widetilde{\Phi}^t$ starting from the filtration on the base.

\begin{proof}[Proof of Lemma \ref{lemma invariant neighborhoods}]
Let us only treat the case of $\Sigma_{\rs/\rso}$, as the other cases can again be treating similarly by considering $\Phi^{-1}$ instead of $\Phi^1$. Consider a filtration $(\Ocal_j^+)_j$ of the manifold $M$ for $\varphi^t$ which is arbitrarily close to the unstable manifolds. The distance of the filtration to the unstable manifolds will be chosen sufficiently small in order to have the claimed quantitative result for the lifted filtration. First, we~want to construct by induction a filtration for the diffeomorphism $\widetilde{\Phi}^1$
\begin{equation}\label{eq: proof lemma invariant neighborhood HR}
\Ucal_j^{\rs/\rso} \: \: \varepsilon\text{-close to } \bigcup_{z \in \oldbigcup_{\ell > N-j} W^\ru (K_\ell)} \Sigma_{\rs/\rso} (z)
\end{equation}
For $j = 1$, we~define $\Ucal_1^{\rs/\rso}$ as a neighborhood of the attractor $\bigcup_{z \in K_N} \Sigma_{\rs/\rso} (z)$ by fixing
\[
\Ucal_1^{\rs/\rso}:= \bigcup_{z \in \varphi^m (\Vcal_N) \cap \Vcal_N} \kappa \bigl(\Ccal_{\ru/\ruo}^{\delta} (x) \bigr),
\]
where $\Vcal_N$ denotes an unrevisited neighborhood of the attractor $K_N$ sufficiently small to be in the range of application of Corollary \ref{corollary big lemma} and where $m$ is an integer. By~choosing $\delta$ small enough and then $m$ sufficiently large, \ie $m \geq m_{\delta}$, the conical neighborhoods and $\Ucal_1^{\rs/\rso}$ are $\widetilde{\Phi}^1$-stable and we can assume that
\[
\max_{(x,\xi) \in \Ucal_1^{\rs/\rso}} d_{S^* M} \Bigl((x,\xi), \bigcup_{z \in K_N} \Sigma_{\rs/\rso}(z) \Bigr) < \varepsilon.
\]
Note that this construction only uses the fact that $K_N$ is an attractor.

Now, let us deal with the induction step and assume that we can find for every $\varepsilon > 0$ some open sets $\Ucal_{j}^{\rs/\rso}$ satisfying \eqref{eq: proof lemma invariant neighborhood HR} for any $j < i$. Fix $\varepsilon > 0$. We want to construct $\Ucal_{i}^{\rs/\rso}$. But before doing that, we~consider an unrevisited neighborhood~$\Vcal$ of $K:= K_{N-i+1}$ small enough in the range of application of Corollary \ref{corollary big lemma}. Also, we~consider the family of unstable annuli $\Acal (m)$ of Sub-lemma \ref{Sublemma: uniform traveling time for unstable annulus} and we choose~$\Vcal$ sufficiently small as in the sub-lemma. According to Lemmas \ref{lemma: attractor and repeller for the flow on the basis} and \ref{lemma attractor and repeller on the phase space} and thanks to the fact that $\Sigma_{\rs/\rso}\cap \Sigma_{\ruo/\ru} = 0_M$ (given by the strong transversality assumption), we~have for all $m \geq 0$
\[
\forall (x,\xi) \in \bigcup_{z \in \Acal (m)} \Sigma_{\rs/\rso} (z), \quad d_{S^* M}\biggl(\widetilde{\Phi}^k (x,\xi), \bigcup_{z \in \oldbigcup_{\ell > N-i +1} W^\ru (K_\ell)} \Sigma_{\rs/\rso} (z) \biggr) \underset{k \to +\infty}{\to}0.
\]
Furthermore, from the compactness proposition \ref{prop compactness}, thanks to Lemma \ref{lemma attractor and repeller on the phase space} and by construction of $\Ucal_{i-1}^{\rs/\rso}$, there exist integers $k_0, m_0 \geq 0$ (potentially larger than the one of the sub-lemma) such that for all $m \geq m_0$,
\begin{equation}\label{eq: lemma invariant neighborhood - uniform cv induction}
\forall (x,\xi) \in \bigcup_{z \in \Acal (m)} \Sigma_{\rs/\rso} (z), \qquad \widetilde{\Phi}^{k_0} (x,\xi) \in \Ucal_{i-1}^{\rs/\rso}.
\end{equation}
By continuity and since $\Ucal_{i-1}^{\rs/\rso}$ is an open set, we~can extend (\ref{eq: lemma invariant neighborhood - uniform cv induction}) on small conical neighborhoods: there exists $\delta_0 > 0$ such that
\[
\forall (x,\xi) \in \bigcup_{z \in \Acal (m)} \kappa (\Ccal_{\ru/\ruo}^{\delta_0}(z)), \qquad \widetilde{\Phi}^{k_0} (x,\xi) \in \Ucal_{i-1}^{\rs/\rso}.
\]
Now, for all $m \geq 0$ and for all $\delta\in(0,1]$, we~define the set
\[
\Wcal (m,\delta):= \bigcup_{z \in \Vcal \cap \varphi^m (\Vcal)} \kappa (\Ccal_{\ru/\ruo}^{\delta}(z)).
\]

According to Lemma \ref{stabilite voisinage conique tilde pour Phi1}, for every $\delta > 0$ there exists an integer $m(\delta) > 0$ such that for all $m \geq m (\delta)$ the open set $\Wcal(m,\delta)$ is an arbitrarily small (as $m \to + \infty$ and $\delta \to 0$) unrevisited neighborhood containing $\bigcup_{z \in \Vcal\cap W^\ru (K)} \kappa \bigl(E_{\ru/\ruo}^* (z) \bigr)$. Actually, we~have something better. For every $m \geq m(\delta)$, the conical neighborhood $\Ccal_{\ru/\ruo}^{\delta}$ is $\Phi^1$-stable on $\Vcal \cap \varphi^m (\Vcal)$ in the sense that, for all $\ell \in \N$,
\begin{equation}\label{eq: lemma invariant neighborhoods - stability}
(x,\xi) \in \Wcal (m,\delta), \: \varphi^{\ell} (x) \in \Vcal \cap \varphi^{m} (\Vcal) \implies \: \forall k \in \lcr0, \ell \rcr, \: \widetilde{\Phi}^k (x,\xi) \in \Wcal (m,\delta).
\end{equation}
Therefore, we~define for all $\delta \leq \delta_0$ and all $m \geq m (\delta)$ the set
\[
\Ucal_i^{\rs/\rso} (m,\delta):= \Ucal_{i-1}^{\rs/\rso} \cup \bigcup_{k = 0}^{k_0 - 1} \widetilde{\Phi}^k (\Wcal (m,\delta)),
\]
which is $\widetilde{\Phi}^1$-stable because $\Ucal_{i-1}^{\rs/\rso}$ is. Also, $\widetilde{\Phi}^{k_0} (\Wcal (m,\delta))$ is equal to
\begin{align*}
& \widetilde{\Phi}^{k_0 } \biggl(\bigcup_{z \in \varphi^{m} (\Vcal) \cap \varphi^{-1} (\Vcal)} \kappa (\Ccal_{\ru/\ruo}^{\delta}(z)) \biggr) \cup \widetilde{\Phi}^{k_0} \biggl(\bigcup_{z \in \varphi^m (\Vcal) \cap \Vcal \cap \Acal (m)} \kappa (\Ccal_{\ru/\ruo}^{\delta}(z)) \biggr) \\
&\subset \widetilde{\Phi}^{k_0 -1} \left(\Wcal (m,\delta) \right) \cup \widetilde{\Phi}^{k_0} \biggl(\bigcup_{z \in \Acal (m)} \kappa (\Ccal_{\ru/\ruo}^{\delta}(z)) \biggr) \subset \Ucal_i^{\rs/\rso} (m,\delta),
\end{align*}
where we used in the first inclusion the fact that $\varphi^{1} (\varphi^m (\Vcal) \cap \varphi^{-1} (\Vcal)) \subset \varphi^m (\Vcal) \cap \Vcal$ (since $\Vcal$ is unrevisited) together with (\ref{eq: lemma invariant neighborhoods - stability}). For the second inclusion, we~used (\ref{eq: lemma invariant neighborhood - uniform cv induction}).

Finally, choosing $\delta$ small enough and then $m$ large enough, we~can assume the set $\Ucal_i^{\rs/\rso} (m,\delta)$ to be a $\varepsilon$-neighborhood of $\bigcup_{z \in \oldbigcup_{j > N-i} W^\ru (K_j)} \Sigma_{\rs/\rso} (z)$. This ends the induction and the proof.
\end{proof}

\begin{remark}
We can note that the set $\Ucal_N^{\rs/\rso}$ constructed in the previous proof defines an arbitrarily small neighborhood of $\Sigma_{\rs/\rso}$ which is stable by $\widetilde{\Phi}^1$.
\end{remark}

\Subsection{Proof of Proposition \ref{prop energy function on the unitary cotangent bundle}: energy function for the Hamiltonian flow}\label{subsection: Proof energy functions for the Hamiltonian flow}

Thanks to Proposition \ref{prop compactness}, Lemmas \ref{lemma attractor and repeller on the phase space} and \ref{lemma invariant neighborhoods}, we~deduce that $(\Sigma_{\rs},\Sigma_{\ruo})$ and $(\Sigma_{\rso},\Sigma_{\ru})$ define attractor-repeller pairs. Therefore, a straight application of Lemma~\ref{lemma: energy function for attractor-repeller model} gives the result.

\section{Construction of the escape function: proof of Proposition \ref{prop escape function}}\label{section: Proof of escape function}

Let us begin with some candidate for the escape function that depends on two auxiliary functions. We will see how the properties of the maps $\E$ and $f$ will be related to those of the escape function. First of all, let us assume that the maps~$\E$ and~$f$ have been defined properly and let us recall the expression of the escape function:
\[
G_m (x,\xi) = m(x,\xi) \log \sqrt{1 + f(x,\xi)^2},
\]
with $m(x, \xi) = \E \left(x, \sfrac{\xi}{|\xi|} \right) \chi (|\xi|^2)$. To lighten notations, we~will use the Japanese bracket $\w{r}:= \sqrt{1 + r^2}$ and the shortcut $\tilde{\xi}$ for $\sfrac{\xi}{|\xi|}$. Now, let us compute $\Lie_{X_H} (G_m)$, for $\|\xi\|\geq 1$
\[
\Lie_{X_H} \left(G_m \right) (x,\xi) = \Lie_{\widetilde{X}_H} (\E) (x, \tilde{\xi}) \log \w{f(x,\xi)} + \E (x, \tilde{\xi}) \Lie_{X_H}(\log \w{f}).
\]
Our goal is to make sure that this quantity is non-positive and is negative outside a conical neighborhood of $E_\ro^*$.

\subsubsection*{Definition of $\E$}
Fix $\varepsilon > 0$ such that
\[
\varepsilon \leq \min \Bigl(\frac{|u|}{2s - 2u}; \: \frac{7}{4}\frac{s}{2s - 2u}; \: \frac{n_0}{2s - n_0} ; \: \frac{1}{4} \frac{n_0}{n_0 - 2u}\Bigr).
\]
We define $E \in \Ccal^{\infty} (S^*M)$ as
\[
\E(x,\xi):= -E(x) + 2s + (2u - n_0) E_+ (x,\xi) + (n_0 - 2s) E_- (x,\xi),
\]
where $E_+$, $E_-$ are given by Proposition \ref{prop energy function on the unitary cotangent bundle} for the constant $\varepsilon$, namely there exist smooth energy functions $E_{\pm} \in \Ccal^{\infty}(S^*M,[0,1])$, $\varepsilon$-neighborhoods $\Wcal^{\rs/\rso}$ of $\Sigma_{\rs/\rso}$, $\Wcal^{\ruo/\ru}$ of $\Sigma_{\ruo/\ru}$ in $S^* M$ and a constant $\eta > 0$ such that:
\begin{align*}
\begin{split}
\Lie_{\widetilde{X}_H} E_{+} &\geq 0 \quad\text{ on } S^* M \quad\text{ and }\quad \Lie_{\widetilde{X}_H} E_{+} > \eta \quad\text{ on } S^* M \setminus (\Wcal^{\ruo} \cup \Wcal^{s}), \\
\Lie_{\widetilde{X}_H} E_{-} &\geq 0 \quad\text{ on } S^* M \quad\text{ and }\quad \Lie_{\widetilde{X}_H} E_{-} > \eta \quad\text{ on } S^* M \setminus (\Wcal^{\ru} \cup \Wcal^{\rso}).
\end{split}
\end{align*}
\begin{figure}[t]
\begin{minipage}[c]{.46\linewidth}
\centering
\def\svgscale{0.7}
\includesvg{Hamiltonianfirst}
\end{minipage}
\begin{minipage}[c]{.46\linewidth}
\centering
\def\svgscale{0.7}
\includesvg{Hamiltoniansecond}
\end{minipage}
\caption{Illustration of some elements used in the proof. Inspired from the picture of \cite[Fig.\,6]{FS}.}
\end{figure}
We also have the estimates $E_{+} \geq 1 - \varepsilon$ on $\Wcal^{s}$, $E_{+} \leq \varepsilon$ on $\Wcal^{\ruo}$, $E_{-} \geq 1 - \varepsilon$ on $\Wcal^{\rso}$ and $E_{-} \leq \varepsilon$ on $\Wcal^{\ru}$. In~order to use these estimates together, let us introduce new open sets in $S^*M$:
\[
\Ncal^{\rs}:= \Wcal^{\rs} \cap \Wcal^{\rso}, \quad \Ncal^{\ro}:= \Wcal^{\rso} \cap \Wcal^{\ruo}, \quad \Ncal^{\ru}:= \Wcal^{\ruo} \cap \Wcal^{\ru}.
\]
Moreover, $E \in \Ccal^{\infty} (M)$ denotes the energy function on the basis given by Proposition~\ref{prop energy function for Axiom A flows} for some parameter $\varepsilon' > 0$ and
\[
\lambda_1 = 0,\dots, \lambda_j = \frac{n_0 (j-1)}{4(N-1)},\dots, \lambda_N = \frac{n_0}{4}.
\]
Thus, the map $E(x)$ has value in the interval $[0, \sfrac{n_0}{4}]$ and there exists a family~$\Ncal_i$ of $\varepsilon'$-neighborhoods of $K_i$ and a constant $\eta_0 > 0$ (which only depends on $\varepsilon'$) such that\footnote{We make the assumption that $N \geq 2$. The case $N=1$ corresponds to the Anosov case.}
\[
- \Lie_V (E) < - \eta_0, \quad \text{on } M \setminus \bigcup_{i = 1}^N \Ncal_i.
\]

In addition, we~choose $\varepsilon$, $\varepsilon'$ sufficiently small so that, for every $i\in\lcr1,N\rcr$,
\begin{align*}
\left\{ z, \: d_g (z,K_i) < \varepsilon' \right\} &\subset \varphi^m (\Vcal_i) \cap \varphi^{-m} (\Vcal_i), \\
\bigcup_{d_g (z,K_i) < \varepsilon'} \bigl\{ (z,\xi) \in \Ncal^{\rs/\ru} \bigr\} &\subset \bigcup_{d_g (z,K_i) < \varepsilon'} \kappa \bigl(\Ccal_{\ru/\rs}^1 (z) \bigr),
\end{align*}
for some given unrevisited neighborhood $\Vcal_i$ of $K_i$ as in the statement of Corollary \ref{corollary big lemma} and for a fixed integer $m$ large enough ($\geq m_1$ of the corollary) such that $C e^{m\lambda/4} \geq 3$ (where $C$, $\lambda$ are the constant given in the corollary). These two inclusions will be used in the construction of $f$. But first, let us briefly explain why these inclusions are satisfied for $\varepsilon, \varepsilon' \ll 1$. Since $m$ is fixed and since $\varphi^m (\Vcal_i) \cap \varphi^{-m} (\Vcal_i)$ is a neighborhood of~$K_i$, the first inclusion is true for $\varepsilon'$ sufficiently small. The second inclusion is given implicitly in the proof of Lemma \ref{lemma invariant neighborhoods}. To be more precise, we~apply Corollary \ref{Corollary compactness proposition} which states that, up to a choice of a smaller $\varepsilon'$,
\[
\kappa \bigl(E_{\ru/\rs}^* (x) \setminus 0_M \bigr) \text{ is arbitrarily close to} \bigcup_{z \in K_i} \kappa \bigl(E_{\ru/\rs}^* (z) \setminus 0_M \bigr)
\]
for all $x$ satisfying $d_g (x,K_i) < \varepsilon'$. If we note that
\[
\bigcup_{d_g (z,K_i) < \varepsilon'} \kappa \bigl(\Ccal_{\rs/\ru}^1 (z) \bigr) \text{ is an open neighborhood of } \bigcup_{z \in K_i} \kappa \bigl(E_{\ru/\rs}^* (z) \setminus 0_M \bigr),
\]
and if we recall that $\Ncal^{\rs/\ru}$ is an $\varepsilon$-neighborhood of $\Sigma_{\rs/\ru} = \bigcup_{z \in M} \kappa \bigl(E_{\ru/\rs}^* (z) \setminus 0_M \bigr)$, then we get that
\[
\bigcup_{ d_g (z,K_i) < \varepsilon'}\bigl\{ (z,\xi) \in \Ncal^{\rs/\ru} \bigr\} \text{ is contained in } \bigcup_{d_g (z,K_i) < \varepsilon'} \kappa \bigl(\Ccal_{\rs/\ru}^1 (z) \bigr),
\]
as soon as $\varepsilon$ is small enough. Another consequence of these two inclusions is that
\begin{equation}\label{eq: proof escape function - inclusion of neighborhoods}
\bigcup_{z \in \Ncal_i} \bigl\{ (z,\xi) \in \Ncal^{\rs/\ru} \bigr\} \subset \bigcup_{\varphi^m (\Vcal_i) \cap \varphi^{-m} (\Vcal_i)} \kappa \bigl(\Ccal_{\rs/\ru}^1 (z) \bigr).
\end{equation}
This last inclusion together with Corollary \ref{corollary big lemma} (second point) will be essential for the definition of $f$.

Let us introduce another notation which will be useful: for any open set $U$ in $S^*M$, we~denote by $\Cone(U)$ the conical neighborhood\footnote{defined by $\Cone(U) = \left\{ (x,\lambda \xi) \in T^* M, \: \lambda \in \R^*, \: (x,\xi) \in U \right\}$.} of $U$ in $T^*M \setminus 0_M$. In~particular, the sets $\Cone(\Ncal^\rs)$, $\Cone(\Ncal^\ro)$ and $\Cone(\Ncal^\ru)$ define conical neighborhoods of $\bigcup_{x \in M} E_\ru^* (x)$, $\bigcup_{x \in M} E_\ro^* (x)$ and $\bigcup_{x \in M} E_\rs^* (x)$ respectively (outside the null section). Now that the setting is done, we~need to check that the escape function satisfies the desired properties.
\begin{itemize}
\item On $\Ncal^\rs$, we~have
\begin{align*}
\E(x,\xi) &\leq \underbrace{-E (x)}_{\leq 0} + 2s + (2u - n_0) (1 - \varepsilon_1) + (n_0 - 2s) (1 - \varepsilon_1) \\
&\leq 2s \varepsilon + 2u (1 - \varepsilon) \leq u.
\end{align*}
\item On $\Ncal^\ru$, we~obtain
\begin{align*}
\E(x,\xi) &\geq \underbrace{-E (x)}_{\geq - s/4} + 2s + (2u - n_0) \varepsilon + (n_0 - 2s) \varepsilon \\
&\geq 2s (1 - \varepsilon - \frac{1}{8}) + 2u \varepsilon \geq s.
\end{align*}
\item On $\Ncal^\ro$, we~obtain
\begin{align*}
\frac{n_0}{2} \leq - \frac{n_0}{4} + n_0 (1 - \varepsilon) + 2u \varepsilon \leq \E(x,\xi) \leq n_0 (1 - \varepsilon) + 2s \varepsilon \leq 2n_0.
\end{align*}
\end{itemize}

We now explain how to construct the function $f$ following the strategy of \cite{FS,dyatlov-zworski}. We want to construct a function $f \in \Ccal^{\infty} (T^*M)$ which is a homogeneous polynomial of degree $1$ for $|\xi| \geq 1$, so that $\Lie_{X_H} (\log \w{f})$ is a bounded function. Since we want $\Lie_{X_H}(G_m)$ to be negative everywhere, we~need to make sure that our definition of $f$ gives the right sign in the term $\E \cdot\Lie_{X_H}(\log \w{f})$. We will also seek $f$ such that $\Lie_{X_H} (f)$ vanishes near $E_\ro^*$ and such that $f(x,\xi) = \chi (|\xi|^2) \tilde{f} (x,\xi)$, with $\chi$ as before and with $\tilde{f} \in \Ccal^{\infty}(T^* M \setminus 0_M)$ which is $1$-homogeneous with respect to the variable $\xi$. We can already check that $\Lie_{X_H} (\log \w{f})$ is a bounded function on $T^* M$ for $|\xi| \geq 1$. Take an arbitrary basic set $K_i$.
\begin{itemize}
\item On $\widetilde{\Ncal}_i^\rs:= \Cone \left(\bigcup_{z \in \Ncal_i} \left\{(z,\xi) \in \Ncal^\rs \right\}\right)$, we~define the map $\tilde{f}$ by
\[
\widetilde{f} (x,\xi):= \int_0^m |\Phi^t (x,\xi)| dt.
\]
Thanks to Corollary \ref{corollary big lemma}, thanks to our choice of $m$ (such that $C e^{m \lambda /4} \geq 3)$ and thanks to the inclusion (\ref{eq: proof escape function - inclusion of neighborhoods}), we~get that
\[
\Lie_{X_H} (\tilde{f})(x,\xi) = |\Phi^m (x,\xi)| - |\xi| \geq Ce^{m\lambda/4} |\xi| - |\xi| \geq 2 |\xi|, \quad \forall (x,\xi) \in \widetilde{\Ncal}_i^\rs.
\]

Also, $f$ is a homogeneous polynomial of degree $1$ in the variable $\xi$, so we can find a constant $c > 0$, which does not depend on $x \in K_i$, such that $0 < c^{-1} |\xi| < \widetilde{f} (x,\xi) < c |\xi|$ on $\widetilde{\Ncal}^\rs \cap \{|\xi| \geq 1\}$. For $|\xi| \geq 1$, we~get
\[
\Lie_{X_H} (f) = \chi (|\xi|^2) \Lie_{X_H} (\tilde{f}) \geq \frac{2}{c}\cdot c|\xi| \geq \frac{2}{c}\, \tilde{f} = \frac{2}{c}\, f.
\]
Therefore, there exists a universal constant $\gamma > 0$ such that
\begin{equation}\label{eq: proof escape function - estimate 1 on f}
\Lie_{X_H} (\log \w{f}) = \frac{f}{\w{f}^2}\, \Lie_{X_H} (f) \geq \frac{2}{c}\, \frac{f^2}{\w{f}^2} \geq \gamma > 0.
\end{equation}
\item On $\widetilde{\Ncal}_i^\ru:= \Cone \left(\bigcup_{z \in \Ncal_i} \{(z,\xi) \in \Ncal^\ru \} \right)$, we~define similarly
\[
\tilde{f}(x,\xi):= \int_0^m |\Phi^{-t} (x,\xi)| dt.
\]

With the same remark, up to reducing the constant $\gamma$, we~obtain
\begin{equation}\label{eq: proof escape function - estimate 2 on f}
\Lie_{X_H} (\log \w{f}) \leq -\gamma < 0,
\end{equation}
as soon as $|\xi| \geq 1$.
\item On $\widetilde{\Ncal}_i^\ro:= \Cone \left(\bigcup_{z \in \Ncal_i} \left\{ (z,\xi) \in \Ncal^\ro \right\} \right)$, we~fix $\tilde{f} (x,\xi) = |\xi (V (x))|$. This ensures that
\[
\Lie_{X_H} (\log \w{f}) = 0.
\]
\item On $ \bigcup_{z \in \oldbigcup_i \Ncal_i} \left\{ (z,\xi) \in S^* M \setminus \left(\Ncal^\rs \sqcup \Ncal^\ro \sqcup \Ncal^\ru \right) \right\}$ and on $\bigcup_{z \notin \oldbigcup_{i} \Ncal_i} S_z^* M$, we~let $\tilde{f}$ take arbitrary positive values on $S^*M$.
\end{itemize}

It now remains to show that $G_m$ has the expected decaying properties (namely points \eqref{prop escape function2} and \eqref{prop escape function3}) of Proposition \ref{prop escape function}.

\subsubsection*{Decaying estimates}
To compute the derivative of the map $E$ along the flow $\widetilde{\Phi}^t$, we~will need at some point the next relation:
\[
\Lie_{\widetilde{X}_H} \E = - \Lie_V (E) + (2u - n_0) \Lie_{\widetilde{X}_H} E_+ + (n_0 -2s) \Lie_{\widetilde{X}_H} E_-
\]
and we can already see that it is non-positive everywhere. For $|\xi| \geq 1$, we~can estimate the quantity $\Lie_{X_H} G_m$ in different directions.
\begin{itemize}
\item On $\widetilde{\Ncal}^\rs:= \Cone \bigl(\bigcup_{z \in \oldbigcup_i \Ncal_i} \left\{ (z,\xi) \in \Ncal^\rs \right\} \bigr)$, we~get
\begin{align*}
\Lie_{X_H} G_m &= \underbrace{\Lie_{\widetilde{X}_H} (\E)}_{\leq 0} \underbrace{\log \w{f}}_{> 0} + \underbrace{\E (x,\tilde{\xi})}_{\leq u } \underbrace{\Lie_{X_H} (\log \w{f})}_{\geq \gamma} \\
&\leq - \gamma |u|.
\end{align*}
\item On $\widetilde{\Ncal}^\ru:= \Cone \bigl(\bigcup_{z \in \oldbigcup_i \Ncal_i} \left\{ (z,\xi) \in \Ncal^\ru \right\} \bigr)$, we~get
\begin{align*}
\Lie_{X_H} G_m &= \underbrace{\Lie_{\widetilde{X}_H} (\E)}_{\leq 0} \underbrace{\log \w{f}}_{> 0} + \underbrace{\E (x,\tilde{\xi})}_{\geq s} \underbrace{\Lie_{X_H} (\log \w{f})}_{\leq - \gamma} \\
&\leq - \gamma s.
\end{align*}
\item On $\widetilde{\Ncal}^\ro:= \Cone \left(\bigcup_{z \in \Ncal_i} \left\{ (z,\xi) \in \Ncal^\ro \right\}\right)$, we~obtain
\[
\Lie_{X_H} G_m = \Lie_{\widetilde{X}_H} (\E) \log \w{f} \leq 0.
\]
\item On $\Cone \bigl(\bigcup_{z \in \oldbigcup_i \Ncal_i} \left\{ (z,\xi) \in S^* M \setminus \left(\Ncal^\rs \sqcup \Ncal^\ro \sqcup \Ncal^\ru \right) \right\} \bigr)$. Or equivalently, for $x \in \bigcup_{1 \leq i \leq N} \Ncal_i$ with $(x,\xi) \notin \Cone (\Wcal^{\rso}\cup \Wcal^{\ru})$ or $(x,\xi) \notin \Cone (\Wcal^{s}\cup \Wcal^{\ruo})$. Since $\rest{\E}{S^* M}$ and $\Lie_{X_H} (\log \w{f})$ are bounded and since $f$ is $1$-homogeneous with respect to $\xi$ for $|\xi| \geq 1$, there exist constants $C_1, C_2 > 0$ such that for all $|\xi| \geq 1$,
\[
\Lie_{X_H} (G_m) = \Lie_{\widetilde{X}_H} (\E) \log \w{ |\xi| C_1} + C_2.
\]
Moreover, according to the construction of $\E$, one can find a positive constant $\eta > 0$ such that $ \Lie_{\widetilde{X}_H} (\E) (x,\tilde{\xi}) < - \eta < 0$. Therefore, there exists a positive radius $R > 0$ such that, for every $(x,\xi)$ in the set $T^*M \setminus \left(\Cone (\Ncal^\rs) \cup \Cone (\Ncal^\ro) \cup \Cone (\Ncal^\ru) \right)$ with $x \in \bigcup_{1 \leq i \leq N} \Ncal_i$ and $|\xi| \geq R$, we~have
\begin{equation}\label{eq: proof of escape function - bound on derivative of the escape function}
\Lie_{X_H} (G_m) \leq - \gamma \min(s, |u|).
\end{equation}
Therefore, we~define $C_m:= \gamma \min(s, |u|) $.
\item Outside a small neighborhood of the non-wandering set, \ie for $x \!\in\! M\!\setminus\! \bigl(\bigcup_{i = 1}^N \Ncal_i \bigr)$. In this case, we have $ \Lie_{\widetilde{X}_H} (\E) \leq - \Psfrac{n_0}{4 (N-1)}\eta_0 < 0$. So, using a similar argument as in the previous point, we~deduce that equation \eqref{eq: proof of escape function - bound on derivative of the escape function} still holds far away from the null section. \qed
\end{itemize}

\begin{appendix}

\section{Hyperbolic sets}\label{Appendix: hyperbolic set}

In this appendix, we~recall the definition of a hyperbolic set.

\begin{definition}A $\varphi^t$-invariant compact set $K$ is said to be \emph{hyperbolic} for the flow~$\varphi^t$ on $M$ if $K$ is the union of isolated fixed points and compact sets on which the induced vector field $V$ never vanishes, where
\begin{itemize}
\item for each $x \in K$, we~have the following decomposition
\begin{equation}\label{decompositionTM}
T_x M = E_\ru (x) \oplus E_\rs(x) \oplus E_\ro (x),
\end{equation}
where $E_\ro (x) = \R V(x)$ and $E_\rs$ (\resp $E_\ru$) is called the stable (\resp unstable) distribution. Moreover, this decomposition is continuous with respect to $x \in K$;
\item the decomposition (\ref{decompositionTM}) is invariant by the flow $\varphi^t$:
\[
\forall x \in K, \quad (D_x \varphi^t)(E_\ru(x)) = E_\ru(\varphi^t (x)) \quad \text{and} \quad (D_x \varphi^t)(E_\rs(x)) = E_\rs(\varphi^t (x))\,;
\]
\item there are constants $C > 0$ and $\lambda > 0$ such that for every $x \in K$ and for every $t \geq 0$, the following inequalities are satisfied:
\begin{equation}\label{eq: hyperbolic inegality on an hyperbolic set}
\begin{split}
| D_x \varphi^t (v_\rs) |_g &\leq C e^{-\lambda t} |v_\rs|_g, \quad \quad \forall v_\rs \in E_\rs(x), \\
| D_x \varphi^{-t} (v_\ru) |_g &\leq C e^{-\lambda t} |v_\ru|_g, \quad \quad \forall v_\ru \in E_\ru(x).
\end{split}
\end{equation}
\end{itemize}
\end{definition}

When $K$ is reduced to a singleton $\{z\}$, we~must have $V(z) \!=\! 0$ and thus \hbox{$E_\ro (x) \!=\! \{0\}$}. In~that case, we~say that $z$ is a hyperbolic fixed point. When the whole manifold is a hyperbolic compact set on which the vector field $V$ never vanishes, the flow is said to be Anosov. It was first introduced by Anosov in \cite{anosov} and this formal definition of hyperbolicity was motivated by the properties of the geodesic flow on negatively curved manifolds. Another famous example of Anosov flow is given by the suspension of an Anosov diffeomorphism. The notion of hyperbolicity was later extended by Smale who defined the notion of an Axiom~A flow which is at the heart of this article. There are many other examples of hyperbolic sets and we refer the reader to \cite{katok-hasselblatt} and \cite{palis-demelo} for a comprehensive study of hyperbolic dynamics. We also refer to \cite{dyatlov}, \cite{FH} for the case of flows.

\skpt
\begin{remark}
\begin{itemize}
\item The definition of hyperbolicity does not depend on the continuous metric~$g$ on~$M$. Indeed, if $g'$ denotes another smooth metric then, by compactness of $M$, $g'$ is equivalent to $g$ and \eqref{eq: hyperbolic inegality on an hyperbolic set} still holds with some constant $C'$ instead of $C$.
\item The distributions $E_\ru$ and $E_\rs$ are only Hölder-continuous in general. Let $d_\ru (x)$ and $d_\rs (x)$ be the dimensions of $E_\ru (x)$ and $E_\rs (x)$ at any point $x \in K$. The maps $d_\rs (x)$ and $d_\ru (x)$ are locally constants on $K$ and $E_\ru$ (\resp $E_\rs$) define a Hölder-continuous section of the Grassmann bundle $G_{d_\ru, n}$ (\resp $G_{d_\ru, n}$) of vector subspaces of dimension~$d_\ru$ (\resp $d_\rs$).
\end{itemize}
\end{remark}

\section{About Smale ordering: proof of Theorem \ref{thm Smale}}

\label{Appendix: proof of Smale ordering}
In this section, we~give a proof of Theorem \ref{thm Smale}. Although it seems to be well known to specialists---see \cite[p.\,158]{robinson} or the proof of Lemma 9.1.11 in \cite[p.\,536]{FH}, we~were not able to locate a precise reference for the proof in the Axiom~A case. In~the particular case of Morse-Smale flows, a proof was given by Smale in \cite{smale-morse} and a detailed proof can for instance be found in \cite[App.\,B]{DRI}. The most technical part of the proof consists in proving the transitivity of Smale's relation. This step is an application of a classical lemma of dynamical systems called the $\lambda$-lemma, which can be found in \cite[Prop.\,6.1.10, p.\,335]{FH}. However, in the usual $\lambda$-lemma for hyperbolic flows (as~referenced), the statement only involves hyperbolic periodic orbit and its proof relies on the analysis of diffeomorphisms near hyperbolic fixed points. Since we aim at dealing with hyperbolic orbit which are not only periodic, we~begin by presenting a version of the $\lambda$-lemma for non periodic hyperbolic orbits which will allow us to proceed similarly to \cite[App.\,B]{DRI}.

\subsection{Proof of Smale's theorem}

In this part, we~prove Theorem \ref{thm Smale} by using the following dynamical lemma whose proof will be given in paragraph \ref{Subsection appx: proof of generalized lambda lemma}.

\begin{prop}[Generalized $\lambda$-lemma]\label{Prop: lambda lemma for non periodic orbit} Let $z$ be a point of a topologically transitive hyperbolic set $K$. Let $\D$ be an embedded disk intersecting the local stable manifold $W_{\varepsilon_0}^{\rso} (z)$ (for some small $\varepsilon_0 > 0$ given by the stable manifold theorem) transversally at some point $q \in W_{\varepsilon_0}^\rs (z)$ such that $\dim (\D) = \dim(E_\ru)$, where $\dim (E_\ru)$ denotes the dimension of the unstable distribution on the topologically transitive hyperbolic set $K$. Then for any $\varepsilon > 0$, there exists $k_1 \in \N$ such that
for each $k \geq k_1$ there is a non-empty embedded disk $\D_k \subset \D$ (of the same dimension) containing $q$ such that $\varphi^k (\D_k)$ is $\Ccal^1$ $\varepsilon$-close\footnote{It means that $\varphi^k (\D_k)$ is the graph of a $\Ccal^1$ map over $W_{\varepsilon_1}^{\ru} (\varphi^k(z))$ whose distance for the $\Ccal^1$ norm is lower than $\varepsilon$. Particularly, it implies that $\varphi^k (\D_k)$ cannot be too small.} to $W_{\varepsilon_1}^{\ru} (\varphi^k(z))$, for some $\varepsilon_1\in(0,\varepsilon_0]$ which only depends on $\varphi^t$, $K$.
\end{prop}

\begin{lemma}[Technical lemma]\label{lemma: technical lemma appendix}
Let $K$ be a basic set and fix $\varepsilon_0 > 0$ given by the stable manifold theorem. Let $x$ be in $W_{\varepsilon_0}^\rs (x_+)$ with $x_+ \in K$ and let $y$ be in $W_{\varepsilon_0}^\ru (y_-)$ with $y_- \in K$, where $\varepsilon_0$ is given by the stable manifold theorem. If $\D_-$, $\D_+$ denote disks which respectively contain $y$, $x$ and intersect $W_{\varepsilon_0}^{\ruo} (y_-)$, $W_{\varepsilon_0}^{\rso} (x_+)$ transversally at a point, then, for all $T > 0$, there exists $t > T$ and there exist disks $\widetilde{\D}_-$, $\widetilde{\D}_+$ contained respectively in $\D_-$, $\D_+$, of same dimension, such that $y \in \widetilde{\D}_-$, $x \in \widetilde{\D}_+$ and
\[
\varphi^t (\widetilde{\D}_{+}) \text{ and } \widetilde{\D}_{-} \text{ intersect transversally}.
\]
\end{lemma}

Before giving the proof of this technical lemma, we~present some applications.

\begin{prop}\label{Appendix: prop intersection stable unstable sets of basic sets}
Let $K$ be a basic set for the Axiom~A flow $\varphi^t$ (we do not assume the strong transversality nor the no-cycle property here). Then, we~have $W^\rs (K) \cap\nobreak W^\ru (K) = K$.
\end{prop}

\begin{proof}
Assume by contradiction that the inclusion $K \subset W^\rs (K) \cap W^\ru (K)$ is strict, \ie there exists an element $x \in M \setminus K$ such that $x \in W^\rs (K) \cap W^\ru (K)$. By~definition of $W^\rs (K)$ and $W^\ru (K)$, there exist elements $x_-, x_+ \in K$ such that $x \in W^\rs (x_+) \cap W^\ru (x_-)$. Since $x \in W^\rs (x_{+})$, we~also know that for all $\varepsilon > 0$ there exists $t_{\varepsilon} > 0$ such that $\varphi^{t_{\varepsilon}} (x)$ is $\varepsilon$-close of $\varphi^{t_{\varepsilon}} (x_+)$. For $\varepsilon$ small enough, the point $x$ cannot belong to the non-wandering set, because otherwise $\varphi^{t_{\varepsilon}} (x)$ would belong to $\Omega$ (by invariance) and would also be $\varepsilon$-close to $K$ and thus belong to $K$ (since the basic sets are isolated). Then, our goal is to deduce from the fact that $K$ is basic that $x$ belongs to $\Omega$. It will lead to the contradiction. To prove that $x \in \Omega$, we~come back to the definition and consider an open set $O$ which contains $x$. Fix $T > 0$. It is sufficient to exhibit a $t > T$ such that $\varphi^t (O)$ intersects $O$. To do so, we~consider small disks $\D_{-}$, $\D_{+}$ contained in~$O$ which contain $x$ and intersect respectively the unstable manifold of $x_-$, the stable manifold of $x_+$ in the sense that $\varphi^{t_{\varepsilon_0}} (\D_+)$ intersects transversally $W_{\varepsilon_0}^{\rso} (\varphi^{t_{\varepsilon_0}} (x_+))$ to a point\footnote{This means that the disk $\D_+$ has the same dimension that $W_{\varepsilon_0}^\ru (\varphi^{t_{\varepsilon_0}} (x_+))$.} and $\varphi^{-t_{\varepsilon_0}} (\D_-)$ intersects transversally $W_{\varepsilon_0}^{\ruo} (\varphi^{-t_{\varepsilon_0}} (x_-))$ to a point\footnote{Similarly, the disk $\D_-$ has the same dimension that $W_{\varepsilon_0}^{\rs} (\varphi^{-t_{\varepsilon_0}} (x_-))$.}, where $\varepsilon_0$ is given by the stable manifold theorem. Now, the result follows from an application of the technical Lemma \ref{lemma: technical lemma appendix} which gives the existence of a $t > T$ and of disks $\widetilde{\D}_{+}$, $\widetilde{\D}_{+}$ contained respectively in $\D_{+}$, $\D_{+}$, such that $x \in \widetilde{\D}_+ \cap \widetilde{\D}_-$ and
\[
\varphi^t (\widetilde{\D}_{+}) \cap \widetilde{\D}_{-} \neq \emptyset.
\]
Therefore $x \in \Omega$, which leads to the expected contradiction.
\end{proof}

\begin{corollary}[Transitivity of Smale relation]\label{cor: transitivity of Smale relation}
Let us assume that the flow $\varphi^t$ satisfies the strong transversality assumption \eqref{Transversality assumption}. Let $K_i, K_j, K_k$ be three basic sets such that $K_i \leq K_j$ and $K_j \leq K_k$, then we have transitivity of the Smale relation in the sense that $K_i \leq K_k$.
\end{corollary}

\begin{proof}[Proof of Corollary \ref{cor: transitivity of Smale relation}]
Let $x$ be in $W^\ru (x_-) \cap W^\rs (x_+)$ with $(x_-,x_+) \in K_i \times K_j$ and let $y$ be in $W^\ru (y_-) \cap W^\rs (y_+)$ with $(y_-,y_+) \in K_j \times K_k$. Up to replacing $(x,x_-,x_+)$ by $(\varphi^{t_{\varepsilon_0}}(x), \varphi^{t_{\varepsilon_0}} (x_-), \varphi^{t_{\varepsilon_0}}(x_+))$ and $(y,y_-,y_+)$ by $(\varphi^{-t_{\varepsilon_0}}(y), \varphi^{-t_{\varepsilon_0}} (y_-), \varphi^{-t_{\varepsilon_0}}(y_+))$, where $\varepsilon_0$ is given by the stable manifold theorem and where $t_{\varepsilon_0} > 0$ is sufficiently large so that $\varphi^{t_{\varepsilon_0}}(x) \in W_{\varepsilon_0}^\rs (\varphi^{t_{\varepsilon_0}}(x_+))$ and $\varphi^{-t_{\varepsilon_0}}(y) \in W_{\varepsilon_0}^\rs (\varphi^{-t_{\varepsilon_0}}(y_-))$, we~can assume that $x \in W^\ru (x_-) \cap W_{\varepsilon_0}^\rs (x_+)$ and $y \in W_{\varepsilon_0}^\ru (y_-) \cap W^\rs (y_+)$. Since the intersections $W^{\ru} (x_-) \cap W^{\rso} (x_+)$, $W^{\ruo} (y_-) \cap W^{\rs} (y_+)$ are transverse, then, for any disk $\D_{+}$ of dimension $n - \dim W^{\rso} (x_+)$ transverse to $W^{\rso} (x_+)$ which contains $x$ and which is contained in $W^{\ru} (x_-)$ and for any disk $\D_{-}$ of dimension $n - \dim W^{\ruo} (y_-)$ transverse to $W^{\ruo} (y_-)$ which contains $y$ and which is contained in $W^{\rs} (y_+)$, Lemma \ref{lemma: technical lemma appendix} ensures that for every $T \geq 0$ there exists a $t > T$ and there exist smaller disks $\widetilde{\D}_-$, $\widetilde{\D}_+$ (as in the statement of the lemma) such that
\[
\varphi^t (\widetilde{\D}_{+}) \cap \widetilde{\D}_{-} \neq \emptyset.
\]
In particular, there exists $z$ such that $z \in W^{\ru} (K_i) \cap W^{\rs}(K_k)$.
\end{proof}

In the next corollary of Lemma \ref{lemma: technical lemma appendix}, we~give a proof of the claimed relations (\ref{eq: thm Smale - equivalence order relation}) and (\ref{eq: thm Smale - closure of stable and unstable sets}).

\begin{corollary}\label{cor: corollary technical lemma - closure of unstable sets} The following statements are satisfied.
\begin{enumerate}
\item\label{cor: corollary technical lemma - closure of unstable sets1} If $\varphi^t$ satisfies the no-cycle property, then for every total order relation on the basic sets (in the sense of Section \ref{subsubsection total order relation}) and for every $i\in\lcr1,N\rcr$ the set $\bigcup_{i \leq j}\!W^\ru (K_j)$ is compact.
\item\label{cor: corollary technical lemma - closure of unstable sets2} If $\varphi^t$ satisfies the strong transversality assumption, then for every basic set $K$, the set $\bigcup_{K', \: K \leq K'} W^\ru (K')$ is compact and is equal to $\overline{W^\ru (K)}$. Moreover, we~have $\overline{W^\ru (K)} \supseteq W^\ru (K')$ if and only if $K \leq K'$.
\end{enumerate}
\end{corollary}

\begin{proof}[Proof of Corollary \ref{cor: corollary technical lemma - closure of unstable sets}\eqref{cor: corollary technical lemma - closure of unstable sets1}]
This step only rests on the existence of filtrations given in Lemma \ref{lemma: equivalence filtration and unrevisited} for any total order relation (Section \ref{subsubsection total order relation}). Indeed, according to Lemma \ref{lemma: equivalence filtration and unrevisited}, for every $i\in\lcr1,N\rcr$, there exists an open set $\Ocal_{N-i + 1}^+$ which contains $\bigcup_{i \leq j} W^\ru (K_j)$, whose closure is $\varphi^1$ stable and does not intersect any $K_\ell$ with $\ell < i$. Now, we~fix some $i\in\lcr1,N\rcr$. From the previous fact, the closure of $\bigcup_{i \leq j} W^\ru (K_j)$ must be contained in $\overline{\Ocal_{N-i + 1}^+}$. Compactness follows from the equality
\[
\bigcup_{i \leq j} W^\ru (K_j) = \bigcap_{n \geq 0} \varphi^n \bigl(\overline{\Ocal_{N-i+1}^+} \bigr).
\]

\subsubsection*{Proof of Corollary \ref{cor: corollary technical lemma - closure of unstable sets}\eqref{cor: corollary technical lemma - closure of unstable sets2}} Point \eqref{cor: corollary technical lemma - closure of unstable sets2} is a consequence of \eqref{cor: corollary technical lemma - closure of unstable sets1} and of Lemma \ref{lemma: technical lemma appendix}, as we shall explain. First, we~fix some basic set $K$. As a consequence of Corollary~\ref{cor: transitivity of Smale relation}, we~can relabel the basic sets so that the set $\{K' \text{ basic}, \: K \leq K' \}$, of cardinal $N-i$ (for some $i\in\lcr0,N-1\rcr$), is sent bijectively to $\{K_j,\: i \leq j \leq N\}$ (with $K = K_i$) and so that the new label on $(K_\ell)_{1 \leq \ell \leq N}$ defines a total order relation in the sense of Section \ref{subsubsection total order relation}. Now, using (1) for this new total order relation, we~deduce that $\bigcup_{i \leq j} W^\ru (K_j) = \bigcup_{K', \: K \leq K'} W^\ru (K')$ is a compact set.

In order to prove the claimed statement, we~apply Lemma \ref{lemma: technical lemma appendix} to show that this last set is included in $\overline{W^\ru (K)}$, or equivalently that for every basic set $K'$ such that $K \leq K'$ we have $W^\ru (K') \subset \overline{W^\ru (K)}$. Let $K'$ be a basic set such that $K \leq K'$, fix~$y \in W^\ru (y_-)$ with $y_- \in K'$ and choose $x \in W^\ru (x_-) \cap W^\rs (x_+)$ with $(x_-, x_+) \in K \times K'$. As in the proof of Corollary \ref{cor: transitivity of Smale relation}, we~can assume that $x \in W^\ru (x_-) \cap W_{\varepsilon_0}^\rs (x_+)$ and $y \in W_{\varepsilon_0}^\ru (y_-)$. Moreover, for every $\varepsilon > 0$, if we choose some disks $\D_+$ and $\D_-$ such that
\[
x \in \D_+ \subset B_g (x, \varepsilon), \qquad y \in \D_- \subset B_g (y, \varepsilon),
\]
as in the statement of Lemma \ref{lemma: technical lemma appendix}, there exists an element $z \in \varphi^t (\D_{+}) \cap \D_{-}$ which satisfies by construction
\[
d_g \left(z, y \right) < \varepsilon \quad\text{ and }\quad z \in \varphi^t (\D_{+}) \subset W^\ru (K).
\]
Since this last result holds for every $\varepsilon > 0$, we~deduce that every $y \in W^\ru (K')$ belongs to $\overline{W^\ru (K)}$. Finally, for every basic set $K'$ such that $K\leq K'$ we have $\overline{W^\ru (K)} \supseteq W^\ru (K')$ and thus
\[
\overline{W^\ru (K)} \supseteq \bigcup_{K', \: K \leq K'} W^\ru (K').\qedhere
\]
\end{proof}

\begin{proof}[Proof of the technical lemma \ref{lemma: technical lemma appendix}]
The result is an application of the $\lambda$-lemma. Consider a disk $\D_{+}$ which is transverse to $W_{\varepsilon_0}^{\rso} (x_+)$ and consider a disk $\D_{-}$ which is transverse to $W_{\varepsilon_0}^{\ruo} (y_-)$ such that
\[
\D_{+} \cap W_{\varepsilon_0}^{\rso} (x_+) = \{q_+ \}, \qquad \D_{-} \cap W_{\varepsilon_0}^{\ruo} (y_-) = \{q_- \}.
\]
Up to replacing $x_+$ and $y_-$ by respectively $\varphi^t (x_+)$ and $\varphi^{t'}(y_-)$ for some $t, t' \in \R$, we~can assume that $q_+ \in W_{\varepsilon_0}^\rs (x_+)$ and $q_- \in W_{\varepsilon_0}^\ru (y_-)$.
Note that the disks $\D_{+}$, $\D_-$ are respectively of dimension $d_{\ru} = \dim \rest{E_{\ru}}{K}$, $d_{\rs} = \dim \rest{E_{\rs}}{K} = n - 1 - d_{\ru}$ (which are constant since $K$ is topologically transitive). In~this proof, we~are going to apply the $\lambda$-lemma (Lemma \ref{Prop: lambda lemma for non periodic orbit}) three times and thus consider smaller disks (as in the statement of the lemma) at each application. To lighten notations, we~still denote them by $\D_{+}$, $\D_-$. Now, applying the $\lambda$-lemma of Proposition \ref{Prop: lambda lemma for non periodic orbit}, we~deduce that, for every $\varepsilon > 0$, there exists $t (\varepsilon) > 0$ such that
\[
\varphi^{t (\varepsilon)} (\D_{+}) \text{ is } \Ccal^1 \: \varepsilon\text{-close to } W_{\varepsilon_1}^{\ru} (\varphi^{ t (\varepsilon)} (x_+))
\]
up to considering a smaller disk $\D_+$ whose size depends on $\varepsilon$, where $\varepsilon_1 > 0$ is the constant appearing in the $\lambda$-lemma.
If we apply the $\lambda$-lemma for the flow $\varphi^{-t}$ instead of $\varphi^t$ then we obtain similarly
\begin{equation}\label{eq: proof technical lemma - reverse lambda lemma}
\varphi^{-t (\varepsilon)} (\D_{-}) \text{ is } \Ccal^1 \: \varepsilon\text{-close to } W_{\varepsilon_1}^{s} (\varphi^{ - t (\varepsilon)} (y_-)),
\end{equation}
up to taking a larger $t (\varepsilon)$ and up to considering a smaller disk $\D_-$ (whose size also depends on $\varepsilon$). Now, we~choose $\varepsilon$ sufficiently small so that it ensures the uniform transversality of $\Ccal^1$ $\varepsilon$-perturbations of the (weak) stable and unstable manifolds on~$K$. Such a choice is always possible because transversality of submanifolds is an open condition for the $\Ccal^1$ topology---we refer to \cite[\S 5.3.3, p.\,88]{laudenbach-book} for this result\footnote{Although the referenced result deals with the transverse intersection of a graph of a $\Ccal^1$ map with a fixed submanifold, it can be applied to the case of transverse intersection of two graphs of $\Ccal^1$ maps. Indeed, since both graphs are submanifolds of $M$ (of dimension $d_{\ruo}$ and $d_\rs$ with $n = d_{\ruo} + d_\rs$ for example), we~can see them as the images of two injective maps $i_1: \R^{d_{\ruo}} \to M$ and $i_2: \R^{d_\rs} \to M$. Moreover, it is not hard to see that the two manifolds are transverse if and only if the map $(i_1,i_2)$ is transverse to the diagonal $\Delta_M:= \{ (x,x), x \in M\}$}---and the uniformity comes from the compactness of the basic set $K$ and  Hölder regularity of the stable and unstable foliations on $K$.

The main idea consists in connecting the points $\varphi^{ t (\varepsilon)} (x_+)$ and $\varphi^{ - t (\varepsilon)} (y_-)$ using periodic orbits whose distance to both points will be chosen sufficiently small. From the assumption that $K$ is topologically transitive, for every $\epsilon > 0$, there exists a point $z_{\epsilon} \in K$ whose positive orbit is dense in $K$ and a time $T (\epsilon) > 0$ such that
\[
d_g (\varphi^{ t (\varepsilon)} (x_+), z_{\epsilon}) < \frac{\epsilon}{2}, \qquad d_g (\varphi^{- t (\varepsilon)} (y_-), \varphi^{T (\epsilon)} (z_{\epsilon})) < \frac{\epsilon}{2}.
\]
By continuity of $\varphi^{T (\epsilon)}$ and since $K$ is the closure of its periodic orbits, we~deduce that there exists a periodic point $p_{\epsilon}$ (which depends on $\epsilon$ and $\varepsilon$) of period $T_0 (\epsilon) > 0$ such that
\[
d_g (\varphi^{ t (\varepsilon)} (x_+), p_{\epsilon}) < \epsilon, \qquad d_g (\varphi^{- t (\varepsilon)} (y_-), \varphi^{T (\epsilon)} (p_{\epsilon})) < \epsilon.
\]
Now, we~are going to use the Bowen brackets and their continuity, namely the continuity of the well-defined maps $[.,.]_{\flat / \sharp}: \{ (z_1,z_2) \in K \times K, \: d_g (z_1,z_2) \leq \epsilon\} \to K$ given by the transverse intersections
\[
W_{\epsilon}^{\ru} (z_1) \cap W_{\epsilon}^{\rso} (z_2) = \{[z_1,z_2]_{\flat} \}, \qquad W_{\epsilon}^{\ruo} (z_1) \cap W_{\epsilon}^{s} (z_2) = \{[z_1,z_2]_{\sharp} \}
\]
for all $z_1,z_2 \in K$ and for all sufficiently small $\epsilon > 0$ (which only depends on $K$ and~$\varphi^t$) as recalled\footnote{Note that the Bowen bracket has values in $K$ means that $K$ has a local product structure. This local product structure follows from the fact that the basic set $K$ is locally maximal (by definition). Indeed, it is not hard to see that $[z_1,z_2]_{\flat}, [z_1,z_2]_{\sharp} \in K$ as soon as $\epsilon$ is chosen small enough so that $\{x, \: d_g (x,K) < \epsilon \} \subset O$, for $O$ open set such that $\bigcap_{t \in \R} \varphi^t (O) = K$. Moreover, the continuity of the Bowen bracket follows from the continuity of the stable and unstable leaves.} in \cite[Prop.\,6.2.2]{FH}. Now, we~fix such a $\epsilon > 0$ and we choose $\epsilon$ sufficiently small so that $\epsilon < \varepsilon_1 /2$. If we note that $ W^{\ru} ([\varphi^{t (\varepsilon)} (x_+),p_{\epsilon}]_{\flat}) = W^{\ru} (\varphi^{t (\varepsilon)} (x_+))$ and that $ W^{\rso} ([\varphi^{t (\varepsilon)} (x_+),p_{\epsilon}]_{\flat}) = W^{\rso} (p_{\epsilon}) $ and if we consider the adapted chart at the point $[\varphi^{t (\varepsilon)} (x_+),p_{\epsilon}]_{\flat}$ which sends $ W_{\epsilon}^{\ru} ([\varphi^{t (\varepsilon)} (x_+),p_{\epsilon}]_{\flat})$ to $\R^{d_{\ru}} \times \{0\}^{d_\rs + 1}$ and $ W_{\epsilon}^{\rso} ([\varphi^{t (\varepsilon)} (x_+),p_{\epsilon}]) $ to $\{0\}^{d_{\ru}} \times \R^{d_\rs+ 1}$, then we get that
\[
\text{the disk $\varphi^{t (\varepsilon)} (\D_{+})$ intersects $W_{\epsilon}^{\rso} ([\varphi^{t (\varepsilon)} (x_+),p_{\epsilon}]_{\flat})$ transversally at a point.}
\]
Furthermore, since $W_{\varepsilon_1}^{\rso} (p_{\epsilon})$ contains $W_{\epsilon}^{\rso} ([\varphi^{t (\varepsilon)} (x_+),p_{\epsilon}]_{\flat})$ thanks to the choice of $\epsilon < \varepsilon_1 / 2$, we~deduce that
\[
\varphi^{t (\varepsilon)} (\D_{+}) \text{ intersects } W_{\varepsilon_1}^{\rso} (p_{\epsilon}) \text{ transversally at the same point}.
\]
Transporting this intersection along the diffeomorphism $\varphi^{T (\epsilon)}$ implies that
\[
\varphi^{t (\varepsilon)+ T (\epsilon)} (\D_{+}) \text{ intersects } \varphi^{T (\epsilon)} \left(W_{\varepsilon_1}^{\rso} (p_{\epsilon}) \right) \subset W_{\varepsilon_1}^{\rso} (\varphi^{T (\epsilon)} (p_{\epsilon})) \text{ transversally}.
\]
Now, if we apply again the $\lambda$-lemma for the disk $\varphi^{t (\varepsilon)+ T (\epsilon)} (\D_{+})$ which intersects transversally the stable manifold of the periodic point $\varphi^{T (\epsilon)} (p_{\epsilon})$, then there exists a disk $D \subset \varphi^{t (\varepsilon)+ T (\epsilon)} (\D_{+})$ of same dimension than $\D_+$ which contains $\varphi^{T (\epsilon)} (p_{\epsilon})$ and there exists an integer $k$ sufficiently large so that
\begin{equation}\label{eq: proof technical lemma - final transversality pushing forward}
\varphi^{k T_0 (\epsilon)} (D) \text{ is } \Ccal^1 \: \varepsilon \text{-close to } W_{\varepsilon_1}^{\ru} (\varphi^{T (\epsilon)} (p_{\epsilon})).
\end{equation}
On the other hand, let us go back to the disk $\varphi^{-t (\varepsilon)}(\D_{-})$ that we moved backward along the flow. We consider an adapted chart near the point $[\varphi^{T (\epsilon)} (p_{\epsilon}), \varphi^{- t (\varepsilon)} (y_-)]_{\sharp}$ such that $W_{\epsilon}^{\ruo}([\varphi^{T (\epsilon)} (p_{\epsilon}), \varphi^{- t (\varepsilon)} (y_-)]_{\sharp}) $ is sent to $\R^{d_\ru + 1} \times \{0\}^{d_\rs}$ and such that $W_{\epsilon}^{s} ([\varphi^{T (\epsilon)} (p_{\varepsilon}), \varphi^{- t (\varepsilon)} (y_-)]_{\sharp}) $ is sent to $\{0\}^{d_\ru + 1 } \times \R^{d_\rs}$. If we note that
\begin{align*}
W^{\ruo}([\varphi^{T (\epsilon)} (p_{\epsilon}), \varphi^{- t (\varepsilon)} (y_-)]_{\sharp}) &= W^{\ruo}(\varphi^{T (\epsilon)} (p_{\epsilon}))\\
\tag*{and}
W^{\rs} ([\varphi^{T (\epsilon)} (p_{\varepsilon}), \varphi^{- t (\varepsilon)} (y_-)]_{\sharp}) &= W^{\rs} (\varphi^{- t (\varepsilon)} (y_-)),
\end{align*}
then we deduce from (\ref{eq: proof technical lemma - reverse lambda lemma}) that
\[
\varphi^{-t (\varepsilon)} (\D_{-}) \text{ is } \Ccal^1 \: \varepsilon\text{-close to } W_{\epsilon}^{s} ([\varphi^{T (\epsilon)} (p_{\epsilon}), \varphi^{- t (\varepsilon)} (y_-)]),
\]
up to considering a smaller disk $\D_-$ (of same dimension and whose size depends on~$\varepsilon,\epsilon$). Finally, thanks to (\ref{eq: proof technical lemma - final transversality pushing forward}) and thanks to the choice of $\varepsilon$, we~deduce that
\[
\bigcup_{|t| < \varepsilon_1} \varphi^t \bigl(\varphi^{ k T_0 (\epsilon)} (D) \bigr) \text{ and }\varphi^{-t (\varepsilon)}(\D_{-}) \text{ intersect each other transversally.}\qedhere
\]
\end{proof}

\begin{figure}[h]
\centering
\def\svgscale{0.7}
\includesvg{proof-smale-order}
\caption{Illustration of some elements used in the proof of Lemma \ref{lemma: technical lemma appendix}.}
\end{figure}

\subsection{Proof of the generalized $\lambda$-lemma}\label{Subsection appx: proof of generalized lambda lemma}

When the orbit $z \in K$ along the flow is periodic (of period $T > 0$), the analysis takes place on a chart near $z$ on which the diffeomorphism $\varphi^T$ acts naturally. In~the case of non periodic orbit, we~need to consider a moving chart defined at each point of the orbit $(\varphi^k (z))_{k \geq 0}$. We in fact use $\Ccal^{1}$-adapted charts for this trajectory used and constructed in \cite[Def.\,4.8, p.\,36]{dyatlov}.

\begin{definition}[Adapted charts]\label{def: appendix adaped chart} Let $x \in K$. A $\Ccal^{1}$ diffeomorphism
\[
\zeta_x: \Ucal_x \to \Vcal_x, \quad x \in \Ucal_x \subset M, \quad 0 \in \Vcal_x \subset \R^n, \quad \zeta_x (x) = 0,
\]
is called an adapted chart for $\varphi^t$ centered at $x$ if:
\begin{enumerate}
\item For each $y \in \Ucal_x$, we~have $D \zeta_x (y) V(y) = \partial_{x_1} = (1,0,\dots,0)^{\top} \in T_{\zeta_x (y)} \R^n = \R^n$.
\item $D \zeta_x (x) E_\ru (x) = E_\ru (0_{\R^n}):= \{0\}^{d_\rs + 1} \times \R^{d_\ru}$ and the restriction of $D\zeta_x (x)$ to $E_\ru (x)$ is an isometry from the adapted metric defined in Section \ref{subsubsection adapted metric} to the Euclidean metric.
\item $D \zeta_x (x) E_\rs (x) = E_\rs (0_{\R^n}):= \{0\} \times \R^{d_\rs} \times \{0\}^{d_\ru}$ and the restriction of $D\zeta_x (x)$ to $E_\ru (x)$ is an isometry from the adapted metric to the euclidean metric.
\end{enumerate}
\end{definition}

For the sake of our analysis, we~consider such a family of adapted charts and we explain how to modify it to straighten the weak-stable manifold and the unstable manifold. Precisely, we~construct a family of charts satisfying the following assumptions.

\begin{definition}[Straightening charts]\label{def: straightening charts}
A family of $\Ccal^{1}$ diffeomorphisms
\[
\chi_x: U_x \to V_x, \quad x \in K,\quad x \in U_x \subset M, \quad 0 \in V_x \subset \R^n, \quad \chi_x (x) = 0,
\]
is called a family of straightening charts for $\varphi^t$ at $K$ if: for all $x \in K$, we~have
\begin{align*}
\chi_x (U_x \cap W_{\varepsilon_0}^{\rso} (x)) &\subset \R^{d_\rs + 1} \times \{0\}^{d_\ru} \\
\chi_x (U_x \cap W_{\varepsilon_0}^{\ru} (x)) &\subset \{0\}^{d_\rs} \times \R^{d_\ru},
\end{align*}
where $\varepsilon_0$ is given by the stable manifold theorem, where $W_{\varepsilon_0}^{\rso} (x)$ denotes the local weak-stable manifold at $x$ and where $W_{\varepsilon_0}^{\ru} (x)$ denotes the local unstable manifold at~$x$.
\end{definition}

First of all, we~begin with explaining how to deduce a family of straightening charts from a family of adapted charts. Arguing as in \cite[p.\,36]{dyatlov}, we~first consider a family of adapted charts $\left(\zeta_x: \Ucal_x \subset M \to \Vcal_x \subset \R^n \right)_{x \in K}$. Then, we~consider the rescaled chart
\[
\widetilde{\zeta}_x:= T \circ \zeta_x: \Ucal_x \to \widetilde{ \Vcal}_{x}:= T\left(\Vcal_{x} \right), \quad x \in K, \quad T = \delta_1^{-1} \Id_{\R_n},
\]
for some $\delta_1 > 0$ and we define the family of diffeomorphisms
\[
\widetilde{\psi}_x:= \widetilde{\zeta}_{\varphi^1 (x)} \circ \varphi^1 \circ \widetilde{\zeta}_{x}^{-1}: \widetilde{\Vcal}_x \subset \R^n \to \widetilde{\Vcal}_{\varphi^1 (x)} \subset \R^n, \qquad x \in K.
\]
As explained in \cite[\S 3.4]{dyatlov}, the stable manifold theorem applied to the family of maps $(\widetilde{\psi}_{\varphi^k (x)})_{k \in \Z}$ gives the existence of a stable manifold $W_{\loc,x}^\rs (0_{\R^n})$ and of an unstable manifold $W_{\loc,x}^\ru (0_{\R^n})$ which are invariant under the family of diffeomorphisms $(\widetilde{\psi}_{\varphi^k (x)})_{k \in \Z}$ as soon as $\delta_1 > 0$ is sufficiently small ($\delta_1$ only depends on $\varphi^t$, $K$). We also assume that $\delta_1$ is small enough so that $[-2,2]^n \subset \widetilde{\Vcal}_x$ for every $x \in K$. Recall that the local stable and unstable manifolds on $M$ are defined by pullback from these local stable and unstable manifolds on $\R^n$.
Moreover, the unstable and stable manifolds restricted to $[-2,2]^n$ can be seen as graph of functions $F_\ru: [-2,2]^{d_\ru} \to [-2,2]^{d_\rs + 1}$, $F_s: [-2,2]^{d_\rs} \to [-2,2]^{d_\ru + 1}$ which are of class $\Ccal^{1,1}$, defined as the set of $\Ccal^1$ functions whose differential is Lipschitz\footnote{The regularity is not optimal here and could be improved using adapted charts with higher regularity, as it was done in \cite{dyatlov}. But it will be enough for the proof of the $\lambda$-lemma.}, and which satisfy the bounds
\begin{equation}\label{eq: bound Fu and Fs}
\| F_\ru \|_{\Ccal^{1}} \leq C \delta_1, \quad \| F_s \|_{\Ccal^{1}} \leq C \delta_1
\end{equation}
for some constant $C > 0$ depending only on $\varphi^t$, $K$. Here, $\|.\|_{\Ccal^{1}}$ denotes the maximum of the $\Ccal^1$ norm and the Lipschitz norm of the differential. Following the construction of \cite{dyatlov}, the weak-unstable is then defined by
\[
W_{\loc,x}^{\ruo} (0):= \left\{ x + (s,0,\dots, 0), \text{ for } x \in W_{\loc,x}^\ru (0) \text{ and } -2 \leq s \leq 2 \right\} \cap [-2,2]^n
\]
and the weak-stable manifold $ W_{\loc,x}^{\rso} (0)$ has a very similar definition. Note that the weak-unstable manifold $W_{\loc,x}^{\rso} (0)$ is the graph of the $\Ccal^{1}$ map
\[
F_{\rso}: (x_1, x_\rs) \in [-2,2]_{x_1} \times [-2,2]_{x_\ru}^{d_\rs} \mto \pi_\ru (F_\rs (x_\rs)) \in [-2,2]^{d_\ru},
\]
where $\pi_\ru (x_1,x_\ru) = x_\ru$ for all $(x_1,x_\ru) \in [-2,2]_{x_1} \times [-2,2]_{x_\rs}^{d_\ru}$. Now, we~consider the straightening operator
\begin{align*}
& \qquad \qquad S: (x_\ro, x_\rs, x_\ru) \in [-2,2]^{n} \mto (\widetilde{x}_\ro, \widetilde{x}_\rs, \widetilde{x}_\ru), \\
&\text{ with } (\widetilde{x}_\ro, \widetilde{x}_\rs)= (x_\ro, x_\rs) - F_\ru (x_\ru), \quad \widetilde{x}_\ru= x_\ru - F_{\rso} (x_\ro,x_\rs),
\end{align*}
where $(x_\ro, x_\rs,x_\ru)$ denotes the decomposition $\R \times \R^{d_\rs} \times \R^{d_\ru} $ of $\R^n$. Moreover, the bounds (\ref{eq: bound Fu and Fs}) ensure that $S$ is a diffeomorphism on its image (which contains the ball $[-2 + C\delta_1,2 - C\delta_1]$ and is contained in $[-2 - C\delta_1,2 + C\delta_1]$). Furthermore, it is not hard to see that $S$ is indeed straightening the weak-stable and the unstable manifolds in the sense that $ S (W_{\loc,x}^{\rso} (0)) \subset \R^{d_\rs + 1} \times \{0\}^{d_\ru}$ and $ S (W_{\loc,x}^{\ru} (0)) \subset \{0\}^{d_\rs + 1} \times \R^{d_\ru}$. Finally, we~introduce the adapted chart
\[
\chi_x = S \circ \widetilde{\zeta}_x \:: \: \: U_x:= \Ucal_x \cap \widetilde{\zeta}_x^{-1} (]-2,2[^n) \to V_x:= \Ran (\chi_x)
\]
which satisfies all the claimed properties.

\subsubsection*{Another zoom for local estimates}
For $\delta > 0$, we~consider the family diffeomorphisms
\[
\psi_{x} = \widetilde{\chi}_{\varphi^1 (x)} \circ \varphi^1 \circ \widetilde{\chi}_x^{-1}: \widetilde{V}_x \to \widetilde{V}_{\varphi^1 (x)}, \quad \psi_{x} (0_{\R^n}) = 0_{\R^n}, \quad x \in K,
\]
where $\widetilde{V}_{x}= \delta^{-1} V_x$, $\widetilde{\chi}_x = \delta^{-1} \chi_x$.
denoted for all $y \in [-1,1]^n$ by
\[
D \psi_{x} (y) = \begin{pmatrix}
A_{\rso,\rso}^{x} (y) & A_{\ru,\rso}^{x} (y) \\
A_{\rso,\ru}^{x} (y) & A_{\ru,\ru}^{x} (y)
\end{pmatrix}
\]
has for coefficients the linear maps
\begin{align*}
A_{\rso,\rso}^{x} (y): \R^{d_\rs + 1} &\to \R^{d_\rs + 1},& A_{\rso,\ru}^{x} (y): \R^{d_\rs + 1} &\to \R^{d_\ru},\\
A_{\ru,\rso}^{x} (y): \R^{d_\ru} &\to \R^{d_\rs + 1},& A_{\ru,\ru}^{x} (y): \R^{d_\ru} &\to \R^{d_\ru},
\end{align*}
which satisfy for all $y \in [-1,1]^n$:
\begin{equation}\label{eq: appendix, estimates differential in adapted chart}
\begin{split}
& \| A_{\rso,\rso}^{x} (y) \|\leq 1 + C' \delta, \quad \| (A_{\ru,\ru}^{x} (y))^{-1} \| \leq e^{- \lambda/2} + C' \delta,\\
&\qquad \| A_{\rso,\ru}^{x} (y) \| \leq C' \delta |y_u|, \quad \| A_{\ru,\rso}^{x} (y) \| \leq C' \delta |y|,
\end{split}
\end{equation}
for some constant $C' > 0$. These estimates are a consequence of (\ref{hyperbolic estimates for the adapted metric}). Moreover, it is important to remark that $A_{\rso,\ru}^{(k)}$ vanishes on the stable manifold $W_{\loc,x}^{\rso} (0) \cap [-1,1]^n = [-1,1]^{d_\rs + 1} \times \{0\}^{d_\ru}$.

\begin{proof}[Proof of Proposition \ref{Prop: lambda lemma for non periodic orbit}]
We work in the family of charts $(\widetilde{\chi}_x)_{x \in K}$ defined just before the proof from a family of straightening charts $(\chi_x)_{x \in K}$ (see Definition \ref{def: straightening charts}) for some $\delta > 0$ sufficiently small so that
\begin{equation}\label{eq: estimees delta preuve lambda lemma}
\begin{split}
&(1 + 3 C' \delta) e^{- \lambda / 4} \leq e^{- \lambda /8}, \quad (1 + C' \delta) (e^{- \lambda/2 } + C' \delta) < e^{- \lambda / 4}, \\
&(e^{-\lambda/ 2} + C' \delta) \Bigl(1 - \frac{C' \delta}{2} (e^{- \lambda/ 2} + C' \delta) \Bigr)^{-1} < e^{- \lambda / 4}.
\end{split}
\end{equation}

Consider a point $z \in K$ whose orbit $(\varphi^k (z))_{k \in \Z}$ is not reduced to a fixed point and consider a disk $\D \subset [-1,1]^n$ which intersects the stable manifold $W_{\loc,z}^{\rso} (0) \cap [-1,1]^n =[-1,1]^{d_\rs +1} \times \{0\}^{d_\ru} $ at some point $q \in W_{\loc,z}^{s} (0)$ in the sense that
\[
\D \cap W_{\loc,z}^{\rso} (0) = \{q\} \quad \text{and} \quad \R^n = T_q \D \oplus T_q W_{\loc,z}^{\rso} (0).
\]
Since $q$ belongs to the stable manifold $ W_{\loc,z}^\rs (0)$, we~deduce from the stable manifold theorem that there exists $c_0 > 0$ which depends on $\varphi^t$, $K$ and $\delta$ such that
\[
\forall k \in \Z, \qquad | \widetilde{\chi}_{\varphi^k (z)} \circ \varphi^k \circ \widetilde{\chi}_{z}^{-1} (q)| \leq c_0 e^{- k\lambda /2} |q|,
\]
and therefore
\[
\left| \psi_z^k (q) \right| \leq c_0 e^{- k \lambda / 2} \quad \text{for } \psi_{z}^k:= \psi_{\varphi^{k-1}(z)} \circ \cdots \circ \psi_{z},
\]
since $\psi_x = \widetilde{\chi}_{\varphi^1 (x)} \circ \varphi^1 \circ \widetilde{\chi}_{x}^{-1} (q)$ ($x \in K$).
Contrary to the usual version of the $\lambda$-lemma, the diffeomorphism we apply to the disk $\D$ changes at each iteration. But since all the diffeomorphisms $\psi_{\varphi^k (z)}$ satisfy hyperbolic estimates which are uniform with respect to $k$, we~claim that the usual proof can be adapted in this setting. For the sake of completeness, we~briefly check that claim.

\Subsubsection*{Step 1: Lipschitz bound for the tangent space of the iterated disks $\varphi^k (\D)$ at $\psi_z^k (q)$}
First, we~can describe the tangent space $T_q \D$ at $q$ as the graph of a linear map $E_0: \R^{d_\ru} \to \R^{d_\rs + 1}$ whose norm is bounded by some constant $\kappa > 0$. Similarly, for every $k \in \N$, we~describe $T_{\varphi^k (q)} \varphi^k (\D)$ as the graph of a linear map $E_{k}: \R^{d_\ru} \to \R^{d_\rs +1}$. The usual strategy consists in using the recursive relation
\begin{equation}\label{eq: proof lamdbda lemma - link tangent spaces}
D \psi_{\varphi^k (z)} (\psi_{z}^k (q)) T_{\psi_{z}^k (q)} \psi_{z}^k (\D)
= T_{\psi_{z}^{k+1} (q)} \psi_{z}^{k+1}(\D), \quad \forall k \in \N,
\end{equation}
to deduce a recursive relation on the norm of $E_{k}$. To simplify the notations, let us denote by $q_k$ the point $\psi_{z}^k (q)$ and $z_k$ the point $\varphi^k (z)$. As a consequence of (\ref{eq: proof lamdbda lemma - link tangent spaces}) and
\begin{equation}\label{eq: proof lambda lemma - image du graphe par la differentielle}
D \psi_{z_k} (q_k) \begin{pmatrix}
E_k \\
I_{d_\ru} \\
\end{pmatrix}
\!\! \begin{pmatrix}
A_{\rso,\rso}^{z_k} (q_k) E_k+ A_{\ru,\rso}^{z_k} (q_k) \\
A_{\rso,\ru}^{z_k} (q_k) E_k + A_{\ru,\ru}^{z_k} (q_k)
\end{pmatrix}
\!=\! \begin{pmatrix}
A_{\rso,\rso}^{z_k} (q_k)E_k + A_{\ru,\rso}^{z_k} (q_k) \\
A_{\ru,\ru}^{z_k} (q_k)
\end{pmatrix}\!,
\end{equation}
since $ A_{\rso,\ru}^{z_k}$ vanishes on $[-1,1]^{d_\rs + 1}\times \{0\}^{d_\ru} $ by stability of the weak stable manifold. But since the terms in (\ref{eq: proof lambda lemma - image du graphe par la differentielle}) belong to the graph of $E_{k+1}$, we~must have
\begin{equation}\label{eq: proof lambda lemma - induction relation}
E_{k+1} = \left(A_{\rso,\rso}^{z_k} (q_k)E_k + A_{\ru,\rso}^{z_k} (q_k) \right) A_{\ru,\ru}^{z_k} (q_k)^{-1}.
\end{equation}
Now, it remains to see that the norms of $E_k$ satisfy an exponential bound. Thanks to~\eqref{eq: proof lambda lemma - induction relation} and thanks to (\ref{eq: appendix, estimates differential in adapted chart}) again, we~deduce that
\begin{align*}
\|E_{k+1}\| &\leq \left(\|A_{\rso,\rso}^{z_k} (q_k) \| \|E_k \| + \| A_{\ru,\rso}^{z_k} (q_k) \| \right) \times \| A_{\ru,\ru}^{z_k} (q_k)^{-1}\| \\
&\leq \big((1 + C' \delta) \|E_k \| + C' \delta |q_k| \big) (e^{-\lambda/ 2} + C' \delta) \\
&\leq e^{- \lambda /4} \|E_k \| + e^{- \lambda /4 } c_0 e^{-k \lambda/ 2},
\end{align*}
according to the estimates on $\delta$ given in \eqref{eq: estimees delta preuve lambda lemma}.
Now, if we consider the auxiliary sequence $u_k = \| E_k \| e^{k \lambda / 4}$, then it is not hard to see that $u_{k + 1} \leq u_k + c_0 e^{-k \lambda /4}$ and, thus, that
\[
\forall k \in \N, \quad \|E_k \| \leq e^{- k \lambda /4} \Bigl(\kappa + \frac{c_0}{1 - e^{-\lambda/ 4}} \Bigr).
\]
In particular, we~get that $\|E_k \|$ converges to $0$ as $k \to +\infty$ and there exists $k_{0} \in \N$, such that
\[
\forall k \geq k_0, \quad |q_k| \leq \frac{1}{2}, \qquad \|E_k \| \leq \frac{1}{4}.
\]
By continuity, there exists a disk $\widetilde{\D} \subset \D$ (of same dimension) which contains $q$ such that $\psi_z^{k_0} (\widetilde{\D})$ is contained in the adapted chart near $\varphi^{k_0} (z)$ and which is the graph of a $\Ccal^{1}$ map $F_{k_0}$. Precisely, there exists $F_{k_0}: [-\eta, \eta]^{d_\ru} \to [-1,1]^{d_\rs + 1}$ for some $\eta \in (0, 1)$ such that
\[
\Graph (F_{k_0}) = \psi_z^{k_0} (\widetilde{\D}), \quad \left(F_{k_0} (0_{\R^{d_\ru}}), 0_{\R^{d_\ru}}\right) = q_k, \quad \sup_{x_\ru \in [-\eta, \eta]^{d_\ru}}\|D F_{k_0} (x_\ru) \| \leq \frac{1}{2}.
\]

Now, we~consider a $\Ccal^1$ extension\footnote{Thanks to the expansivity of the unstable manifold, the $\Ccal^1$ convergence of the disk will not depend on the extension $\widetilde{F}_{k_0}$ of $F_{k_0}$. This fact will be detailed in step 3.} of $F_{k_0}$, that we denote by $\widetilde{F}_{k_0}$, such that
\[
\widetilde{F}_{k_0} \in \Ccal^1 \left([-1,1]^{d_\ru}; [-1,1]^{d_\rs + 1} \right), \quad \rest{\widetilde{F}_{k_0}}{[-\eta, \eta]^{d_\ru}}= F_{k_0}, \quad \sup_{x_\ru \in [-1, 1]^{d_\ru}}\|D \widetilde{F}_{k_0} (x_\ru) \| \leq \frac{1}{2}.
\]

\subsubsection*{Step 2: a graph transform argument for the map $\widetilde{F}_{k_0}$}
We first introduce the complete metric space
\begin{multline*}
\Ecal:= \Bigl\{ F \in \Ccal^{1} \bigl([-1,1]^{d_\ru}; [-1,1]^{d_\rs + 1} \bigr),\
\left(F(0_{\R^{d_\ru}}), 0_{\R^{d_\ru}} \right) \in W_{\loc,z}^\rs(0),\\
|F(0_{\R^{d_\ru}})| \leq \frac{1}{2}, \quad
\sup_{x \in [-1,1]^n} \|DF(x)\| \leq \frac{1}{2} \Bigr\}.
\end{multline*}
endowed with the distance $d_{\Ecal} (.,.)$ defined, for all $F, \widetilde{F} \in \Ecal$, by
\[
d_{\Ecal}(F, \widetilde{F}) = \max \Bigl(|F (0_{\R^{d_\ru}}) - \widetilde{F} (0_{\R^{d_\ru}})|, \sup_{x_\ru \in [-1,1]^n} \|DF (x_\ru) - D\widetilde{F} (x_\ru) \| \Bigr).
\]
Now, we~consider the family of graph transforms defined by
\begin{equation}\label{eq: definition graph transform}
\Graph \left(\Psi_k (F) \right) = \psi_{z_k} (\Graph F) \cap [-1, 1]^n,
\end{equation}
for all $ k \in \Z$ and for all $F \in \Ecal$. The fact that $[-1, 1]^{d_\ru} \subset \pi_\ru \left(\psi_{z_k} \left(\Graph F \right)\right)$ (which is implicitly used in \eqref{eq: definition graph transform}), where $\pi_\ru: [-1,1]^n \to [-1,1]^{d_\ru}$ denotes the orthogonal projection to $[-1,1]^{d_\ru}$, follows from the stronger fact: there exists $\Lambda > 1$ (which depends on $\varphi^t$, $K$, $\delta$) such that
\begin{equation}\label{eq: expansivity in the unstable direction}
\forall F \in \Ecal, \: \forall k \in \Z, \: \forall x_\ru \in [- 1, 1]^{d_\ru}, \quad \frac{d}{d x_\ru} \pi_\ru \left(\psi_{z_k} \left(F(x_\ru), x_\ru \right) \right) \geq \Lambda.
\end{equation}

This last inequality is a consequence of the expansivity of the unstable manifold. Now that the definition of the graph transforms given in \eqref{eq: definition graph transform} makes sense, our goal is to prove that for every $F \in \Ecal$, for all $k \in \Z$ and for all $\ell \in \N$,
\begin{equation}\label{eq: proof lambda lemma - contracting inequality}
d_{\Ecal} \left(\Psi_k^{\ell}(F), 0 \right) \leq \max (c_0, 1) e^{-\ell \lambda / 8} d_{\Ecal} \left(F, 0 \right),
\end{equation}
where $\Psi_k^{\ell}$ denotes the composed map $\Psi_{k + \ell} \circ \cdots \circ \Psi_{k}$.
We can remark that \eqref{eq: proof lambda lemma - contracting inequality} implies the $\Ccal^1$ convergence of $\Psi_k^{\ell}(F)$ to the null map, which is geometrically the unstable manifold $W_{\loc, z}^\ru (0)$.

In order to prove \eqref{eq: proof lambda lemma - contracting inequality}, let $k \in \Z$, $F \in \Ccal^{1} ([-1,1]^{d_\ru}, [-1,1]^{d_\rs + 1})$ and let $x = (F(x_\ru), x_\ru) \in [-1,1]^{n}$ for some $x_\ru \in [-1,1]^{d_\ru}$. By~definition of the graph transform given in \eqref{eq: definition graph transform}, we~have
\begin{equation}
D \psi_{z_k} (x) \begin{pmatrix}
DF (x_\ru) \\
I_{d_\ru} \\
\end{pmatrix}
= \begin{pmatrix}
A_{\rso,\rso}^{z_k} (x) DF (x_\ru) + A_{\ru,\rso}^{z_k} (x) \\
A_{\rso,\ru}^{z_k} (x) DF (x_\ru) + A_{\ru,\ru}^{z_k} (x)
\end{pmatrix}.
\end{equation}
From the definition of the graph transform $\Psi_k$, we~get that
\[
D \Psi_k (F) \left(\psi_{z_k} (x)_u \right) = \left(A_{\rso,\rso}^{z_k} (x) DF (x_\ru) + A_{\ru,\rso}^{z_k} (x) \right) \left(A_{\rso,\ru}^{z_k} (x) DF (x_\ru) + A_{\ru,\ru}^{z_k} (x) \right)^{-1},
\]
where $\psi_{z_k} (x) = (\psi_{z_k} (x)_{\rso}, \psi_{z_k} (x)_u)$. Taking the operator norm of $ D \Psi_k (F) \left(\psi_{z_k} (x)_u \right)$, we~get
\[
\left\| D \Psi_k (F) \left(\psi_{z_k} (x)_u \right) \right\| \leq \left\| A_{\rso,\rso}^{z_k} (x) DF (x_\ru) + A_{\ru,\rso}^{z_k} (x) \right\| \times I,
\]
with $I = \bigl\| \left(A_{\rso,\ru}^{z_k} (x) DF (x_\ru) + A_{\ru,\ru}^{z_k} (x) \right)^{-1} \bigr\|$. Moreover, we~have
\begin{equation}\label{eq: majoration inverse}
I \leq \bigl\| \left(A_{\ru,\ru}^{z_k} (x) \right)^{-1} \bigr\| \underbrace{\Bigl\| \left(\Id + A_{\rso,\ru}^{z_k} (x) DF (x_\ru) \left(A_{\ru,\ru}^{z_k} (x) \right)^{-1} \right)^{-1} \Bigr\|}_{J}.
\end{equation}
Since $F \in \Ecal$, we~have that
\begin{align*}
\bigl\| A_{\rso,\ru}^{z_k} (x) DF (x_\ru) \left(A_{\ru,\ru}^{z_k} (x) \right)^{-1} \bigr\| &\leq \left\| A_{\rso,\ru}^{z_k} (x)\right\| \left\| DF (x_\ru)\right\| \bigl\| \left(A_{\ru,\ru}^{z_k} (x) \right)^{-1} \bigr\| \\
&\leq \frac{C' \delta}{2} (e^{- \lambda/2 } + C' \delta) < 1.
\end{align*}
Therefore, we~deduce that
\begin{align*}
J &\leq \sum_{k = 0}^{+ \infty} \bigl\| A_{\rso,\ru}^{z_k} (x) DF (x_\ru) \left(A_{\ru,\ru}^{z_k} (x) \right)^{-1} \bigr\|^k \leq \sum_{k = 0}^{+ \infty} \Bigl(\frac{C' \delta}{2} (e^{- \lambda/2 } + C' \delta) \Bigr)^k \\
&\leq \Bigl(1 - \frac{C' \delta}{2} (e^{- \lambda/2 } + C' \delta) \Bigr)^{-1}.
\end{align*}
Putting all together the last inequalities with \eqref{eq: majoration inverse} and \eqref{eq: estimees delta preuve lambda lemma}, we~get that
\[
I \leq (e^{-\lambda/ 2} + C' \delta) \Bigl(1 - \frac{C' \delta}{2} (e^{- \lambda/ 2} + C' \delta) \Bigr)^{-1} < e^{- \lambda / 4}.
\]
Finally, if we denote by $\|F'\|_{\infty}$ the quantity $\sup_{y_u \in [-1,1]^n}\left\| DF (y_u) \right\|$ to simplify, then we get that
\begin{equation}\label{eq: proof of lambda lemma - upper bound graph transform derivative}
\begin{split}
\left\| D \Psi_k (F) \left(\psi_{z_k} (x)_u \right) \right\| &\leq \left(\left\| A_{\rso,\rso}^{z_k} (x)\right\| \left\| DF (x_\ru) \right\| + \left\| A_{\ru,\rso}^{z_k} (x) \right\| \right) e^{- \lambda/4}\\
&\leq \left((1 + C' \delta) \|F'\|_{\infty} + C' \delta |F(x_\ru)| \right)e^{- \lambda/4} \\
&\leq \left((1 + C' \delta) \|F'\|_{\infty} + C' \delta \left(|F(0_{\R^{d_\ru}})| + \|F'\|_{\infty} \right) \right) e^{- \lambda / 4}
\\
&\leq e^{- \lambda /8} d_{\Ecal}(F, 0),
\end{split}
\end{equation}
where we used the estimates on $\delta$ given in \eqref{eq: estimees delta preuve lambda lemma}.

Furthermore, since $(F(0_{\R^{d_\ru}}), 0_{\R^{d_\ru}}) \in W_{\loc,z}^\rs (0)$, we~get that
\[
\forall k \in \Z, \quad \psi_{z_k} (F(0_{\R^{d_\ru}}), 0_{\R^{d_\ru}}) = (\Psi_k F (0_{\R^{d_\ru}}), 0_{\R^{d_\ru}}) \in W_{\loc,z}^\rs (0).
\]
Since
\[
| \psi_{z_k}^{\ell} (F(0_{\R^{d_\ru}}), 0_{\R^{d_\ru}})| \leq c_0 e^{- \ell \lambda / 2} | (F(0_{\R^{d_\ru}}), 0_{\R^{d_\ru}})| = c_0 e^{- \ell \lambda / 2} |F(0_{\R^{d_\ru}})|,
\]
we~deduce that
\begin{equation}\label{eq: proof of lambda lemma - upper bound graph transform value at 0}
\forall k \in \Z, \quad \forall \ell \in \N \qquad |\Psi_k^{\ell} (F) (0_{\R^{d_\ru}})| \leq c_0 e^{- \ell\lambda / 2} |F(0_{\R^{d_\ru}})|.
\end{equation}
Fix $\ell_0 \in \N$ such that $c_0 e^{-\ell_0 \lambda /8} \leq 1$. A mean value inequality leads to the stability of the graph transforms:
\begin{equation}\label{eq: proof of lambda lemma - stability of the graph transforms}
\forall F \in \Ecal, \quad \forall k \in \Z, \quad \forall \ell \geq \ell_0 \qquad \Psi_k^{\ell} F \in \Ecal,
\end{equation}
and the inequalities \eqref{eq: proof of lambda lemma - upper bound graph transform derivative} and \eqref{eq: proof of lambda lemma - upper bound graph transform value at 0} leads to the claimed inequality \eqref{eq: proof lambda lemma - contracting inequality}.

\subsubsection*{Step 3: coming back to the disk}
Applying step 2 to the map $\widetilde{F}_{k_0}$ defined at the end of step 1, we~get in particular that
\[
d_{\Ecal} \bigl(\Psi_{k_0}^{\ell} (\widetilde{F}_{k_0}), 0 \bigr) \leq \max (c_0, 1) e^{- \ell \lambda/8} d_{\Ecal}(\widetilde{F}_{k_0}, 0) \leq \max (c_0, 1)\, \frac{e^{- \ell \lambda/8}}{2}.
\]
In particular, thanks to the mean value inequality, we~deduce that
\[
\bigl\| \Psi_{k_0}^{\ell} (\widetilde{F}_{k_0}) \bigr\|_{\Ccal^1} \underset{\ell \to +\infty}{\to} 0.
\]
Finally, thanks to the expansivity of the unstable manifold, to the stability given in~\eqref{eq: proof of lambda lemma - stability of the graph transforms} and to the above $\Ccal^1$-convergence, we~obtain
\[
\Graph \bigl(\Psi_{k_0}^{\ell} (\widetilde{F}_{k_0}) \bigr) \cap [-1,1]^n = \Graph \left(\Psi_{k_0}^{\ell} (F_{k_0}) \right) \cap [-1,1]^n = \psi_z^{k_0 + \ell}(\widetilde{\D}) \cap [-1,1]^n,
\]
for all $\ell \geq \ell_0$.

Now, if we fix $\varepsilon_1 > 0$ such that $\varepsilon_1 < \varepsilon_0$ and
\[
\forall y \in K, \qquad W_{\varepsilon_1}^\ru (y) \subset \widetilde{\chi}_{y} \bigl(\{0\}^{d_\rs + 1} \times [-1,1]^{d_\ru} \bigr),
\]
then we can deduce the following statement: for all $\varepsilon > 0$, there exists $k_1 \in \N$ such that for all $k \geq k_1$, there exists a disk $\widetilde{\D}_k \subset \widetilde{\D}$ (of same dimension) which contains $q$ such that $\varphi^k(\widetilde{\D}_k)$ is $\Ccal^{1}$ $\varepsilon$-close to $W_{\varepsilon_1}^\ru (\varphi^k (z))$.
\end{proof}

\end{appendix}
